% concatenated from Milnor_main.tex
\documentclass[11pt,reqno]{amsart}

\usepackage[T1]{fontenc}
\usepackage{lmodern}
\usepackage[margin=1.06in]{geometry}
\usepackage{microtype}
\usepackage{mathtools}
\usepackage{amssymb}
\usepackage{enumitem}
\usepackage{booktabs}
\usepackage{tabularx}
\usepackage{xcolor}
\usepackage[colorlinks=true,
 linkcolor=blue!55!black,citecolor=blue!55!black,
 urlcolor=blue!65!black]{hyperref}
\hypersetup{
 pdftitle={A unified Ricci-positive square-root cable in dimensions four and five},
 pdfauthor={Nan Wu},
 pdfsubject={Counterexamples to the Milnor conjecture in dimensions at least four},
 pdfkeywords={positive Ricci curvature, fundamental groups, Milnor conjecture,
 Pruefer group, complete noncompact manifolds}}

\newtheorem{theorem}{Theorem}[section]
\newtheorem{lemma}[theorem]{Lemma}
\newtheorem{proposition}[theorem]{Proposition}
\newtheorem{corollary}[theorem]{Corollary}
\theoremstyle{definition}
\newtheorem{definition}[theorem]{Definition}
\newtheorem{remark}[theorem]{Remark}
\numberwithin{equation}{section}

\DeclareMathOperator{\Ric}{Ric}
\DeclareMathOperator{\Rm}{Rm}
\DeclareMathOperator{\II}{II}
\DeclareMathOperator{\Hess}{Hess}
\DeclareMathOperator{\tr}{tr}
\DeclareMathOperator{\divsymb}{div}
\DeclareMathOperator{\sgn}{sgn}
\DeclareMathOperator{\Id}{Id}
\DeclareMathOperator{\Sym}{Sym}
\DeclareMathOperator{\vspan}{span}
\newcommand{\dd}{\mathop{}\!\mathrm d}
\newcommand{\rhoCab}{{\rho_{\mathrm{cab}}}}
\newcommand{\CPrufer}{C_{2^\infty}}

\title[Counterexamples to Milnor's Conjecture]
{Pr\"ufer $2$-group and Milnor's Conjecture on Fundamental Groups}

\author[N. Wu]{Nan Wu}
\address[N. Wu]{Department of Mathematics \& IMS \\ Nanjing University \\ Nanjing \\ 210093 \\ China}
\email{nwu2@nju.edu.cn}
\author[Z. Yan]{Zetian Yan}
\address[Z. Yan]{Department of Mathematics \\ The Chinese University of Hong Kong \\ Shatin, N.T. \\ Hong Kong}
\email{zetianyan@cuhk.edu.hk}

%\author[X. Zhu]{Xingyu Zhu}
%\address[X. Zhu]{Department of Mathematics \\ Michigan State University \\ East Lansing \\ MI 48824 \\ USA}
%\email{zhuxing3@msu.edu}

\subjclass[2020]{Primary 53C20; Secondary 53C21, 57R19}
\keywords{positive Ricci curvature, fundamental groups, Milnor conjecture,
Pr\"ufer group, complete noncompact manifolds}
\date{August 15, 2026}

\begin{document}
\raggedbottom

\begin{abstract}
We construct complete, one-ended Riemannian manifolds in dimensions four
and five having strictly positive Ricci curvature, with fundamental group
isomorphic to the Pr\"ufer $2$-group $C_{2^\infty}$. This gives counterexamples
to Milnor's conjecture in the two remaining dimensions. 
\end{abstract}

\maketitle
\tableofcontents
\enlargethispage{4pt}



\section{Introduction}
\label{sec:statement}
Milnor conjectured that the fundamental group of every complete
Riemannian manifold with nonnegative Ricci curvature is finitely
generated \cite{Milnor1968}.  The conjecture is known to hold in
dimensions at most three
\cite{CohnVossen1935,Liu2013,Pan2020}.  More recently, Bru\`e, Naber,
and Semola constructed counterexamples in dimensions great than six 
\cite{BNS7,BNS6}.  Their construction proceeds on the universal cover
through an equivariant smooth snowflake, leaving dimensions four and
five open.

The aim of this paper is to handle the remaining case by constructing counterexamples
in dimension four and five via the direct quotient.  Moreover, the manifold we constructed actually have strictly positive Ricci
curvature. The fundamental group shown in our construction is the following Pr\"ufer $2$-group.
\begin{definition}
We define the Pr\"ufer $2$-group as
\[
 \CPrufer
 :=\varinjlim_j
 \left(
 \mathbb Z/2^{j+1}\mathbb Z
 \longrightarrow
 \mathbb Z/2^{j+2}\mathbb Z
 \right),
 \qquad
 [k]\longmapsto[2k].
\]
Equivalently,
\[
 \CPrufer
 =
 \bigcup_{j\ge0}
 \frac{1}{2^{j+1}}\mathbb Z\big/\mathbb Z
 \subset \mathbb Q/\mathbb Z.
\]
\end{definition}
Every finitely generated subgroup of $\CPrufer$ is contained in a finite
cyclic subgroup of order $2^N$ for some $N$, whereas $\CPrufer$ contains
elements of arbitrarily large order.  Consequently, $\CPrufer$ is not
finitely generated.

Our main Theorem can be stated as follows.
\begin{theorem}\label{thm:main}
For each $m\in\{2,3\}$, there exists a connected, orientable, one-ended
smooth manifold $M^{m+2}$ admitting a complete Riemannian metric $g$, such
that
\[
 \Ric_g>0,
 \qquad
 \pi_1(M^{m+2})\cong\CPrufer.
\]
In particular, Milnor's conjecture fails in dimensions four and five, respectively.
\end{theorem}

\begin{corollary}\label{cor:all-dimensions}
For every integer $n\ge4$, there exists a connected, orientable, one-ended
smooth $n$-manifold admitting a complete Riemannian metric of strictly
positive Ricci curvature and having fundamental group isomorphic to
$\CPrufer$.
\end{corollary}

\begin{proof}
For $n=4,5$, this follows directly from
Theorem~\ref{thm:main}.  Let $n\ge6$, and let $(M^4,g_4)$ be the
four-dimensional example furnished by Theorem~\ref{thm:main}.  Equip
\[
 N^n:=M^4\times\mathbb S^{n-4}
\]
with the standard product metric $g_4\times g_{\mathbb S^{n-4}}$. The resulting metric is complete, and $N^n$
remains connected, orientable, and one-ended.  Since $n-4\ge2$, the
sphere is simply connected, and hence
\[
 \pi_1(N^n)
 \cong\pi_1(M^4)
 \cong\CPrufer.
\]
Moreover, for every $(v,w)\in T_xM^4\oplus T_y\mathbb S^{n-4}$,
\[
 \Ric_{N^n}\bigl((v,w),(v,w)\bigr)
 =
 \Ric_{g_4}(v,v)
 +(n-5)\lvert w\rvert^2.
\]
Because $\Ric_{g_4}>0$ and $n-5>0$, this quantity is positive whenever
$(v,w)\ne0$.
\end{proof}

Before sketching the construction in Theorem~\ref{thm:main}, we first
briefly recall several major developments surrounding Milnor's
conjecture.  Milnor's original argument showed that every finitely
generated subgroup of the fundamental group has polynomial growth
\cite{Milnor1968}; Gromov's theorem then implies that every such
subgroup is virtually nilpotent \cite{Gromov1981}.  Wilking subsequently
reduced the search for counterexamples to the case of abelian
fundamental groups, while the generalized Margulis lemma of
Kapovitch--Wilking provided dimensionally uniform control of the
underlying nilpotent structure
\cite{Wilking2000,KapovitchWilking2011}.  On the affirmative side, the
conjecture holds in dimension two by the work of Cohn--Vossen and in
dimension three by the work of Schoen--Yau, Liu, and Pan
\cite{CohnVossen1935,SchoenYau1982,Liu2013,Pan2020}.  A parallel line
of results established finite generation under additional geometric
hypotheses: Li and Anderson proved finiteness under Euclidean volume
growth, Sormani treated small linear diameter growth, and subsequent
work of Pan and Huang imposed various stability, conicality, and
polarity assumptions on the universal cover
\cite{Li1986,AndersonTopology1990,SormaniDiameter,PanStability,
PanAlmostStability,HuangPolar}.  More recently, Huang--Huang proved
finite generation in dimension four when the universal cover has
Euclidean volume growth, while Navarro--Pan--Zhu obtained a stronger
structural description in the linear-volume-growth setting
\cite{HuangHuang2025,NavarroPanZhu2024}.  The picture changed
decisively with the constructions of Bru\`e--Naber--Semola
\cite{BNS7,BNS6}: their
seven- and six-dimensional examples realize arbitrary subgroups of
\(\mathbb Q/\mathbb Z\) and \(\mathbb Q/\mathbb Z\), respectively,
leaving only dimensions four and five open.
Theorem~\ref{thm:main}
closes this dimensional gap and shows, moreover, that the conjecture
already fails under strictly positive Ricci curvature.

We now give a sketch of the construction in Theorem~\ref{thm:main}. The starting point of the construction is a simple algebraic feature of
the Pr\"ufer $2$-group.  It admits the presentation
\[
 \CPrufer
 =
 \left\langle
 a_0,a_1,a_2,\ldots
 \ \middle|\
 a_0^2=1,\quad a_j=a_{j+1}^2\ \text{for all }j\ge0
 \right\rangle.
\]
Thus each $a_{j+1}$ can be viewed as a square root of $a_j$, while the
relation $a_0^2=1$ supplies the initial torsion.  Such presentation
suggests the following geometric construction.  A \emph{product parametrization of the boundary} of a compact cobordism $B$ with boundary components
$\partial_iB$ and $\partial_oB$, each diffeomorphic to
$\mathbb S^1\times\mathbb S^m$, is a specified pair of diffeomorphisms
\[
 \Phi_I:\mathbb S^1\times\mathbb S^m\longrightarrow\partial_iB,
 \qquad
 \Phi_O:\mathbb S^1\times\mathbb S^m\longrightarrow\partial_oB.
\]
For a fixed $y\in\mathbb S^m$, these parametrizations  determine the
incoming and outgoing product circles
\[
 \lambda_{\mathrm{in}}(\theta)=\Phi_I(\theta,y),
 \qquad
 \lambda_{\mathrm{out}}(\theta)=\Phi_O(\theta,y).
\]
We call $(B;\Phi_I,\Phi_O)$ a \emph{square-root block} if
\begin{equation}\label{def:square-root-block}
 \pi_1(B)\cong\mathbb Z=\langle a\rangle,
 \qquad
 [\lambda_{\mathrm{out}}]=a,
 \qquad
 [\lambda_{\mathrm{in}}]=a^2.
\end{equation}
Thus attaching the incoming boundary of one such block to the outgoing
boundary of the preceding block introduces the relation
$a_j=a_{j+1}^2$.  The construction below supplies the boundary metrics
and second fundamental forms required to repeat this attachment.

The incoming boundary of the first block is filled by a compact
manifold diffeomorphic to \(D^2\times\mathbb S^m\).  Under the natural
identification
\[
 \partial(D^2\times\mathbb S^m)
 =\mathbb S^1\times\mathbb S^m,
\]
each circle \(\mathbb S^1\times\{y\}\) bounds the embedded disk
\(D^2\times\{y\}\).  We call this filling the \emph{initial cap}. Its
attachment imposes $a_0^2=1$.  Repeating the block construction then
generates the entire group.

If the relation $a_0^2=1$ is omitted, the remaining presentation is
\[
 \left\langle a_0,a_1,\ldots
 \ \middle|\
 a_j=a_{j+1}^2\ \text{for all }j\ge0
 \right\rangle
 \cong\mathbb Z[1/2].
\]
This presentation has a classical three-dimensional realization.
In the dyadic solenoid
construction, one considers a nested sequence of unknotted solid tori,
each winding twice inside its predecessor.  The fundamental group of
each successive complement is infinite cyclic, and its distinguished
generators satisfy
\[
 [\gamma_j]=2[\gamma_{j+1}].
\]
The resulting direct limit is $\mathbb Z[1/2]$.  This model was discussed
by Shen and Sormani in connection with the Milnor conjecture
\cite{ShenSormani2006}, appears in Semola's lecture notes as a
topological precursor to the construction of Bru\`e, Naber, and Semola
\cite{SemolaNotes2025}, and belongs to the broader theory of solenoid
complements studied in \cite{ConnerMeilstrupRepovs2015}.

In the notation used below, this three-dimensional model arises formally
by setting $m=1$.  In that case,
\[
 P_1
 =
 D^2\setminus
 \left(
 \operatorname{int}D^2_+
 \sqcup
 \operatorname{int}D^2_-
 \right)
\]
is an ordinary pair of pants, and the exterior of the winding-two circle
is its mapping torus under an involution exchanging the two inner
boundary circles.  This is precisely the elementary block underlying
the dyadic solenoid complement.  We point out that this
model is different with the Whitehead manifold.  Such dyadic pattern has algebraic
winding number two and gives rise to the direct-limit group
$\mathbb Z[1/2]$, whereas the Whitehead construction has algebraic winding
number zero and produces a contractible open manifold.

The formal case $m=1$ also identifies the key modification required for
our construction.  Let
\[
 c_m\subset \mathbb S^1\times D^{m+1}
\]
denote the winding-two circle, and let $\mathcal N(c_m)$ be a closed
tubular neighborhood of $c_m$.  When $m=1$, the circle $c_1$ has
codimension two, and removing the interior of $\mathcal N(c_1)$ introduces
an additional meridian.  For $m\ge2$, by contrast, $c_m$ has
codimension $m+1\ge3$, and general position shows that the inclusion
induces an isomorphism
\[
 \pi_1\left(
   (\mathbb S^1\times D^{m+1})
   \setminus\operatorname{int}\mathcal N(c_m)
 \right)
 \xrightarrow{\;\cong\;}
 \pi_1(\mathbb S^1\times D^{m+1})
 \cong\mathbb Z.
\]
Let $\lambda_{\mathrm{out}}$ and $\lambda_{\mathrm{in}}$ denote the
distinguished outer and inner longitudes, respectively.  If
\[
 \gamma:=[\lambda_{\mathrm{out}}],
\]
then $\gamma$ generates the fundamental group of the exterior, whereas
\[
 [\lambda_{\mathrm{in}}]=\gamma^2.
\]
Thus, viewed as a cobordism from its inner boundary to its outer
boundary, this exterior carries the incoming longitude class
\(\gamma^2\) to the outgoing longitude class \(\gamma\).  In other
words, it realizes the square-root relation without introducing an
additional meridian; see Proposition~\ref{prop:double-cable}.


Under the preceding identification with $\mathbb Z[1/2]$, the relation
$a_0^2=1$ imposed by the initial cap kills the integral subgroup
$\mathbb Z$, leaving
\[
 \mathbb Z[1/2]\big/\mathbb Z
 \cong
 \CPrufer.
\]
This quotient is the algebraic cornerstone for the construction; the
resulting manifold is not obtained by taking a literal quotient of the
three-dimensional solenoid complement.

It remains to realize such a building block in dimension $m+2$ with
positive Ricci curvature and with boundary metrics and second fundamental
forms that permit successive attachments. We carry out this realization in three steps: we construct an equivariant
Ricci-positive metric on the directed base \(P_m\) (see
(\ref{def:directed-base})); form the corresponding screw
quotient, which realizes the square-root relation (see
(\ref{def:screw-quotient})); and attach a Ricci-positive
holonomy connector that removes the resulting boundary holonomy (see (\ref{def:holonomy-connector})).


For each $m\in\{2,3\}$, we first construct
\[
 P_m\cong
 D^{m+1}\setminus
 \left(
 \operatorname{int}D^{m+1}_+
 \sqcup
 \operatorname{int}D^{m+1}_-
 \right),
\]
with boundary decomposition
\[
 \partial P_m=\partial_oP_m\sqcup I_+\sqcup I_-.
\]
We designate $\partial_oP_m$ as the outgoing boundary and $I_+$ and
$I_-$ as the incoming boundaries.  The tuple
\begin{equation}\label{def:directed-base}
 (P_m,\sigma_m;\partial_oP_m,I_+,I_-),
\end{equation}
where $\sigma_m$ is an involution preserving $\partial_oP_m$ and
exchanging $I_+$ and $I_-$, is called a \emph{directed base}; the word
``directed'' refers precisely to this assignment of incoming and outgoing
boundary components.

Given any $R_*>0$ and $\varepsilon_*>0$, we will show that the manifold $P_m$ carries a
metric $h_m$ with $\Ric_{h_m}>0$ and an isometric involution $\sigma_m$
with these properties.  The incoming boundary components
are round $m$-spheres of radius $R_*$ and satisfy
\[
 \mathrm{II}_{I_\pm}
 >
 -\varepsilon_*\,h_m|_{I_\pm},
\]
whereas the outer boundary is round and strictly convex. In this paper, we fix the convention
\[
 \II(X,Y)=g(\nabla_X\nu,Y),
\]
where \(\nu\) is the outward unit normal; thus \(\II>0\) means strict convexity.  
On the outer
boundary, the involution restricts to an axial half-turn when $m=2$ and
to the antipodal map when $m=3$; see
Theorem~\ref{thm:unified-directed-base}.

Although the underlying topological model also makes sense for $m=1$,
the boundary geometry required above does not.  Indeed, $P_1$ is a pair
of pants, and positive Ricci curvature in dimension two is equivalent
to positive Gaussian curvature.  Suppose that an analogous metric
existed on $P_1$, with two inner boundary circles of radius $R_*$,
strictly convex outer boundary, and
\[
 \mathrm{II}_{I_\pm}
 >
 -\varepsilon_*\,h_1|_{I_\pm}.
\]
With our sign convention, the geodesic curvature of each boundary
component is $\mathrm{II}(T,T)$ for a unit tangent vector $T$.
Since $\chi(P_1)=-1$ and
$\operatorname{Length}(I_\pm)=2\pi R_*$, the Gauss--Bonnet theorem would
give
\[
 -2\pi
 =
 2\pi\chi(P_1)
 =
 \int_{P_1}K_{h_1}\,\dd A
 +
 \int_{\partial P_1}k_g\,\dd s
 >
 -4\pi R_*\varepsilon_*.
\]
This is impossible whenever
$\varepsilon_*\le (2R_*)^{-1}$.  Thus the restriction $m\ge2$ is
essential in the construction of the directed base: the failure of
such estimates when $m=1$ reflects a genuine
two-dimensional obstruction.

For $m=2,3$, the metric on the directed base $P_m$ is constructed using Perelman's
angular-neck construction \cite{Perelman1997}: a sectionally positively
curved angular ball contains two smaller angular balls exchanged by an
involution.  After removing the smaller balls, we attach identical
Ricci-positive necks along the resulting inner hypersurfaces, thereby
normalizing the induced metrics and second fundamental forms on
$I_+$ and $I_-$.  The metric induced on $\partial_oP_m$ is
then joined to a round metric by the equivariant Ricci-positive cylinder
construction of Theorem~\ref{thm:path-cylinder-unified}.  The required
invariant paths of boundary metrics are
provided by Ricci flow: Hamilton's surface theorem is used when $m=2$,
and his three-dimensional convergence theorem when $m=3$
\cite{Hamilton1982, Hamilton1988Surfaces}.

Starting from $(P_m,h_m,\sigma_m)$, we introduce a circle direction and
choose a positive $\sigma_m$-invariant function $f$, such that the warped
metric
\[
 G=h_m+f^2\dd t^2
\]
on $\mathbb R\times P_m$ has positive Ricci curvature and the required
boundary second fundamental forms, by Lemma~\ref{lem:un-warp}.  For
$T>0$, the map
\[
 \mathcal T(t,x):=(t+T,\sigma_m x)
\]
is an isometry of $(\mathbb R\times P_m,G)$.  We call
\begin{equation}\label{def:screw-quotient}
 W_m:=(\mathbb R\times P_m)/\langle\mathcal T\rangle
\end{equation}
the \emph{screw quotient}.  On its outer boundary, the distinguished
mapping-torus loop closes after one application of $\mathcal T$.  On the
inner boundary, the first application of $\mathcal T$ exchanges the two
components, so the first return to either component is given by
$\mathcal T^2$.  This difference in
return times is the geometric origin of the relation
\[
 [\lambda_{\mathrm{in}}]
 =
 [\lambda_{\mathrm{out}}]^2.
\]
The product parametrization of the boundary and this square relation are recorded
in Lemma~\ref{lem:un-boundary-components}.

The two quotient boundaries have the desired relative periods but
with different holonomy: the outer boundary carries a half-turn, whereas
the inner boundary is an untwisted product.  To interpolate between
them, we introduce the \emph{holonomy connector}
\begin{equation}\label{def:holonomy-connector}
 \mathcal C=[0,s_*]\times\mathbb S^1\times\mathbb S^m,
 \qquad
 g_{\mathcal C}=\dd s^2+g_{\ell(s),\rho(s),\omega(s)},
\end{equation}
where $\omega(0)=1/2$ and $\omega(s_*)=0$.  Thus the twist parameter $\omega$ changes
the half-turn holonomy on the outer boundary of $W_m$ into the untwisted
product structure required at the outgoing boundary of the square-root
block.  Once the circle and sphere radius parameters $\ell$ and $\rho$
have been fixed, the twist parameter $\omega$ is chosen to minimize the
corresponding weighted Dirichlet energy \eqref{eq:weighted Dirichlet energy}.  Its Euler--Lagrange equation
is
\[
 \left(
   \frac{\rho^{m+2}\omega'}{\ell}
 \right)'=0,
\]
which makes the mixed Ricci component vanish identically. The variational characterization is given in Lemma~\ref{lem:least-energy-twist}, while the corresponding
mixed-Ricci cancellation follows from Lemma~\ref{lem:un-ricci}. The resulting Ricci-positive holonomy connector is constructed in
Proposition~\ref{prop:un-connector}.  Its compatibility with the screw
quotient and the subsequent smoothing follow from
Lemmas~\ref{lem:un-left-boundary-gluing}
and~\ref{lem:relative-gluing-unified}.

The screw quotient and the holonomy connector together form a compact building block whose untwisted product parametrizations of the boundary permit the
construction to be iterated, as established in
Proposition~\ref{prop:square-root-block}.  The required initial cap is
supplied by Lemma~\ref{lem:un-cap}.  Attaching this cap and then an
infinite chain of the blocks from
Proposition~\ref{prop:square-root-block} produces a locally finite,
one-ended smooth manifold with positive Ricci curvature. The
fundamental group of its $N$-th compact stage is
\[
 \mathbb Z/2^{N+1}\mathbb Z,
\]
and the inclusion into the next stage is multiplication by two.
Consequently,
\[
 \pi_1(M)
 \cong
 \varinjlim_N
 \left(
   \mathbb Z/2^{N+1}\mathbb Z,
   [k]\longmapsto[2k]
 \right)
 \cong
 \CPrufer.
\]
The finite-stage computation and the direct limit are
carried out in
\eqref{eq:un-finite-pi1}--\eqref{eq:un-prufer}.
Finally, each holonomy connector contains a central subcylinder disjoint
from all subsequent smoothing collars.  Choosing these subcylinders so
that their accumulative
crossing distance diverges forces every curve passing through
infinitely many stages to have infinite length. The resulting metric is therefore complete; see Section~\ref{sec:infinite-telescope} for detailed argument.

The present quotient construction is related to the smooth
snowflake construction of Bru\`e, Naber, and Semola
\cite{BNS7,BNS6}, but the two
constructions organize the branching in fundamentally different ways.
In their seven-dimensional construction \cite{BNS7}, one begins with an arbitrary
subgroup $\Gamma\le\mathbb Q/\mathbb Z$ and an exhaustion by finite
cyclic subgroups whose generators satisfy
\[
 \gamma_j^{k_j}=\gamma_{j-1}
\]
for a sequence of integers $k_j\ge2$.  The universal cover is then
constructed together with a compatible free action of $\Gamma$, and
the variable branching orders $k_j$ are encoded by an equivariant smooth
snowflake and a highly twisted gluing procedure.  Their six-dimensional
construction carries out a similar argument for
$\Gamma=\mathbb Q/\mathbb Z$ \cite{BNS6}.

Our construction instead works directly on the quotient and specializes
to the choices $k_j\equiv2$.  The involution $\sigma_m$ exchanging
the two incoming boundary components realizes this order-two interchange
directly on the quotient, and the same construction is
repeated at
every stage.  Consequently, we do not need to construct compatible
actions of increasingly large cyclic groups on the universal cover or
to untwist general elements of a higher-dimensional mapping class group.
This reduction to a single involution and a single block construction is
what makes the direct four- and five-dimensional construction possible.

\begin{remark}\label{rem:beyond-binary}
It is natural to ask whether such quotient construction can realize the
full range of groups obtained by Bru\`e, Naber, and Semola
\cite{BNS7,BNS6}.  Given a
relation
\[
 \gamma_j^{k_j}=\gamma_{j-1},
\]
the topological analogue of our building block would be a directed base in the form of
\[
 P_{m,k_j}
 \cong
 D^{m+1}\setminus
 \biggl(
   \bigsqcup_{\alpha=1}^{k_j}
   \operatorname{int}D^{m+1}_\alpha
 \biggr),
\]
together with an order-$k_j$ symmetry cyclically permuting its inner
boundary components.  The associated screw quotient would have first
return time $k_jT$ on the inner boundary and would therefore encode the
required $k_j$-th-root relation.

The main difficulty is geometric.  One would need Ricci-positive
metrics on $P_{m,k}$, for arbitrary $k$, with cyclic symmetry and with
the same roundness and second fundamental form control required for
iteration.  One would also need holonomy connectors that correct the
resulting order-$k$ outer holonomy while preserving positive Ricci
curvature.  The angular-neck and holonomy-connector constructions developed
here are tailored to the two-fold symmetry of a half-turn or antipodal
map and do not presently provide the corresponding quantitative control
as $k$ varies.  Even a fixed prime $p>2$ requires a new $p$-fold directed
base, while realizing $\mathbb Q/\mathbb Z$ requires a sequence of
branching orders that may be unbounded.

We nevertheless expect that for every $m\in\{2,3\}$ and every subgroup
$\Gamma\le\mathbb Q/\mathbb Z$, there exists a complete, one-ended
Riemannian manifold $M^{m+2}_\Gamma$ satisfying
\[
 \Ric_{M^{m+2}_\Gamma}>0,
 \qquad
 \pi_1(M^{m+2}_\Gamma)\cong\Gamma.
\]
\end{remark}

The final section records several further features of the constructed
manifolds. First, their universal covers $\widetilde M$ are diffeomorphic to
\[
 \mathbb R^{m+2}\#_{\infty}
 (\mathbb S^2\times\mathbb S^m),
\]
and are therefore one-ended and simply connected at infinity.  Next, the
free parameters in the holonomy connectors also allow any prescribed
sequence in \([1,3/2]\) to be realized as logarithmic volume-growth
exponents along suitable subsequences of radii, simultaneously on the
manifold $M$ and its universal cover $\widetilde M$. Furthermore, for the asymptotic volume ratio defined by
\begin{equation*}
    \operatorname{AVR}(M):=\lim_{r\to\infty}
 \frac{\operatorname{Vol_M}(B_r(x))}
 {\operatorname{Vol}(B_r^{\mathbb R^n})},
\end{equation*}
the manifold $M$ has zero $\operatorname{AVR}(M)$ for every admissible choice of
parameters and has super linear growth:
\begin{equation*}
    \limsup_{r\to \infty} \frac{\operatorname{Vol_M}(B_r(x))}{r}=\infty.
\end{equation*}
When \(m=2\), the universal cover $\widetilde M$ has zero $\operatorname{AVR}(\widetilde M)$
for every admissible choice of parameters as well. However, for arbitrary parameter choices when \(m=3\), the present argument does not determine the asymptotic volume ratio of \(\widetilde M\). Nevertheless, the alternating choice \[ e_{2j}=1, \qquad e_{2j-1}=\frac32 \qquad (j\ge1) \] in Proposition~\ref{prop:volume-growth-flexibility} produces examples with \(m=3\) for which \[ \operatorname{AVR}(\widetilde M)=0. \]


These results give a concrete
description of the topology and large-scale geometry of the examples
and clarify how they lie outside several known finite-generation
regimes.

\subsection*{Organization of the paper}

Section~2 isolates the topological square-root relation.
Section~3 establishes the relative gluing and path-cylinder tools.
Section~4 constructs the Ricci-positive directed base, with the
angular-neck calculation deferred to Appendix~A.
Sections~5 and~6 construct the screw quotient and the holonomy
connector, respectively.
Section~7 assembles the initial cap and the square-root blocks into a
complete infinite telescope.
Finally, Section~8 studies the topology of the universal cover, the
flexibility of the volume growth, and the asymptotic volume ratio.

\subsection*{Acknowledgments}
The authors would like to thank Xingyu Zhu for many helpful discussions, Xianzhe Dai and Guofang Wei for bringing this problem to their attention, and Xuezhang Chen, Qing Han, and Jiayin Pan for valuable comments and suggestions.



\subsection*{Disclosure on AI assistance}

The authors used AI-assisted tools, principally ChatGPT/Codex, for
brainstorming, literature searches and source triage, objection
generation, preliminary checks of calculations and parameter
estimates, notation and exposition review, and manuscript editing.
AI output was not treated as mathematical authority.  The authors
independently worked and verified all theorem statements, proofs,
calculations, constants, citations, and final prose, and take full
responsibility for the contents of the paper.


\section{The topological square-root relation}
\label{sec:double-cable}

Before introducing the metric, we isolate the topological origin of the
square-root relation. Consider an embedded circle
\[
c\subset\mathbb S^{1}\times D^{m+1}
\]
such that the restriction to \(c\) of the projection
\(\mathbb S^{1}\times D^{m+1}\to\mathbb S^{1}\) has degree two, whereas its projection
to \(D^{m+1}\) parametrizes an embedded circle once. Over each point of
\(S^{1}\), the curve \(c\) meets the corresponding fiber in two points.
Consequently, the complement of a sufficiently small tubular
neighborhood of \(c\) fibers over \(\mathbb S^{1}\) with fiber
\[
D^{m+1}\setminus
\left(
\operatorname{int}D^{m+1}_{+}
\sqcup
\operatorname{int}D^{m+1}_{-}
\right).
\]
Its monodromy interchanges the two inner boundary components. We now
describe this fibration explicitly; the resulting fiber, together with
its order-two monodromy, will serve as the topological model for the
directed base constructed later.

Let $m\ge2$.  For every $\rho>0$, denote by
\[
 D^{m+1}_\rho:=\{x\in\mathbb R^{m+1}:|x|\le\rho\}
\]
the closed Euclidean $(m+1)$-ball of radius $\rho$.  Abbreviate
$D^{m+1}:=D^{m+1}_1$ and put
\[
 V:=\mathbb S^1\times D^{m+1}.
\]
Let $e_1,e_{m+1}$ be the first and last standard basis vectors of
$\mathbb R^{m+1}$, write $\mathrm i=\sqrt{-1}$, and define the
$\mathbb R^{m+1}$-valued function
\[
 u:\mathbb R/2\pi\mathbb Z\longrightarrow\mathbb R^{m+1},
 \qquad
 u(\theta):=(\cos\theta,\sin\theta,0,\ldots,0).
\]
Further, we fix
\[
 0<\varepsilon<1,
 \qquad
 0<\rhoCab<\min\{\varepsilon,1-\varepsilon\},
\]
and define
\begin{align}
 c_2(\theta)
   &:=\bigl(e^{2\mathrm i\theta},\varepsilon u(\theta)\bigr),\label{eq:double-cable}\\
 \mathcal F_{\rhoCab}(\theta,v)
   &:=\bigl(e^{2\mathrm i\theta},\varepsilon u(\theta)+v\bigr),\label{eq:explicit-tube}\\
 \nu_{\rhoCab}(c_2)
   &:=\mathcal F_{\rhoCab}
     \bigl((\mathbb R/2\pi\mathbb Z)\times
                     D^{m+1}_{\rhoCab}\bigr),\notag\\
 C_m
   &:=V\setminus\mathcal F_{\rhoCab}
     \bigl((\mathbb R/2\pi\mathbb Z)\times
              \operatorname{int}D^{m+1}_{\rhoCab}\bigr).\label{eq:explicit-exterior}
\end{align}
The first assertion of Proposition~\ref{prop:double-cable} below shows
that $\mathcal F_{\rhoCab}$ is a tubular embedding.  We introduce the two boundary components of $C_m$:
\[
 \partial_oC_m:=\mathbb S^1\times\partial D^{m+1},
 \qquad
 \partial_iC_m:=\mathcal F_{\rhoCab}
 \bigl((\mathbb R/2\pi\mathbb Z)\times
                         \partial D^{m+1}_{\rhoCab}\bigr).
\]

We use the base points, longitudes, and connecting path
\begin{align*}
 x_O&:=(1,e_{m+1}),
 &\lambda_O(\varphi)&:=\bigl(e^{\mathrm i\varphi},e_{m+1}\bigr),\\
 x_I&:=(1,\varepsilon e_1+\rhoCab e_{m+1}),
 &\lambda_I(\theta)&:=
   \mathcal F_{\rhoCab}(\theta,\rhoCab e_{m+1}),\\
 \beta(a)&:=\left(1,
 a\varepsilon e_1+(1-a+a\rhoCab)e_{m+1}\right),
 &\theta,\varphi&\in\mathbb R/2\pi\mathbb Z,
 &0&\le a\le1.
\end{align*}
Here $\lambda_O$ is the positively oriented outer longitude based at
$x_O$, while $\lambda_I$ is the positively oriented longitude of the
drilled tube based at $x_I$ and determined by the constant normal vector
$\rhoCab e_{m+1}$.

For $\alpha\in\mathbb R$, let $R_\alpha\in SO(m+1)$ be rotation through
angle $\alpha$ in the $(e_1,e_2)$-plane and the identity on its
orthogonal complement.  Explicitly,
\[
 R_\alpha(x_1,x_2,x_3,\ldots,x_{m+1})
 =\bigl(x_1\cos\alpha-x_2\sin\alpha,
        x_1\sin\alpha+x_2\cos\alpha,
        x_3,\ldots,x_{m+1}\bigr).
\]
Define
\begin{align*}
 D^{m+1}_+&:=\{x\in D^{m+1}:|x-\varepsilon e_1|\le\rhoCab\},
 &D^{m+1}_-&:=\{x\in D^{m+1}:|x+\varepsilon e_1|\le\rhoCab\},\\
 P_m&:=D^{m+1}\setminus
 \bigl(\operatorname{int}D^{m+1}_+
       \sqcup\operatorname{int}D^{m+1}_-\bigr),
 &I_\pm&:=\partial D^{m+1}_\pm.
\end{align*}
The inequalities imposed on $\rhoCab$ ensure that $D^{m+1}_+$ and
$D^{m+1}_-$ are disjoint and contained in the interior of $D^{m+1}$.
Thus $P_m$ is an $(m+1)$-disk with the interiors of two disjoint
$(m+1)$-disks removed.  The half-turn $R_\pi$ preserves $P_m$ and
exchanges $I_+$ and $I_-$.  Put
\[
 \phi:=R_\pi|_{P_m},
 \qquad
 T_\phi:=\bigl(P_m\times[0,2\pi]\bigr)/
       \bigl((x,2\pi)\sim(\phi(x),0)\bigr).
\]
This Euclidean $P_m$ is the fixed topological model used throughout the
paper.  In Section~\ref{sec:unified-directed-base}, the same symbol
denotes a manifold equipped with a Ricci-positive metric and a specified
diffeomorphism to this model.

\begin{proposition}
\label{prop:double-cable}
The exterior $C_m$ defined above has the following properties.
\begin{enumerate}[label=\textup{(\roman*)}]
\item $\mathcal F_{\rhoCab}$ is a smooth embedding,
\[
 \partial C_m=\partial_oC_m\sqcup\partial_iC_m,
 \qquad
 \partial_oC_m\cong\mathbb S^1\times\mathbb S^m,
 \qquad
 \partial_iC_m\cong\mathbb S^1\times\mathbb S^m.
\]
\item If $j:C_m\hookrightarrow V$ is the inclusion, then
\[
 j_*:\pi_1(C_m,x_O)\xrightarrow{\ \cong\ }\pi_1(V,x_O),
 \qquad
 \pi_1(C_m,x_O)=\langle g\rangle\cong\mathbb Z,
 \qquad
 g=[\lambda_O].
\]
\item The path $\beta$ is contained in $C_m$, satisfies
\[
 \beta(0)=x_O,
 \qquad
 \beta(1)=x_I,
\]
and, with $\beta^{-1}(a)=\beta(1-a)$,
\[
 \bigl[\beta*\lambda_I*\beta^{-1}\bigr]=g^2
 \quad\text{in }\pi_1(C_m,x_O).
\]
For every path $\widehat\beta:[0,1]\to C_m$ from $x_O$ to $x_I$,
\[
 \bigl[\widehat\beta*\lambda_I*\widehat\beta^{-1}\bigr]=g^2.
\]
\item The map
\[
 \Phi:T_\phi\longrightarrow C_m,
 \qquad
 \Phi([x,\varphi])=(e^{\mathrm i\varphi},R_{\varphi/2}x),
\]
is a diffeomorphism, and
\[
 \phi(I_+)=I_-,
 \qquad
 \phi(I_-)=I_+.
\]
\end{enumerate}
\end{proposition}

\begin{proof}
\emph{Part \textup{(i)}.}
For $(\theta,v),(\theta',v')\in
(\mathbb R/2\pi\mathbb Z)\times D^{m+1}_{\rhoCab}$, the equality
$\mathcal F_{\rhoCab}(\theta,v)
 =\mathcal F_{\rhoCab}(\theta',v')$ gives
\[
 e^{2\mathrm i\theta}=e^{2\mathrm i\theta'},
 \qquad
 \theta'-\theta\in\{0,\pi\}\pmod{2\pi}.
\]
If $\theta'=\theta$, then $v'=v$.  If
$\theta'=\theta+\pi$, then
\[
 u(\theta')=-u(\theta),
 \qquad
 v'-v=2\varepsilon u(\theta),
\]
and hence
\[
 2\varepsilon=|v'-v|\le |v'|+|v|\le2\rhoCab<2\varepsilon,
\]
a contradiction.  Thus $\mathcal F_{\rhoCab}$ is injective.  For
$a\in\mathbb R$ and $w\in\mathbb R^{m+1}$, its differential is
\[
 d\mathcal F_{\rhoCab,(\theta,v)}(a,w)
 =\bigl(2a\mathrm i e^{2\mathrm i\theta},
         \varepsilon a u'(\theta)+w\bigr),
\]
and
\[
 d\mathcal F_{\rhoCab,(\theta,v)}(a,w)=0
 \quad\Longrightarrow\quad a=0\quad\Longrightarrow\quad w=0.
\]
Consequently
\[
 \operatorname{rank}d\mathcal F_{\rhoCab,(\theta,v)}=m+2.
\]
Since the domain of $\mathcal F_{\rhoCab}$ is compact and $V$ is Hausdorff,
$\mathcal F_{\rhoCab}$ is a smooth embedding.  The estimate
\[
 |\varepsilon u(\theta)+v|\le\varepsilon+\rhoCab<1
\]
gives
\[
 \nu_{\rhoCab}(c_2)\subset\mathbb S^1\times\operatorname{int}D^{m+1}.
\]
It follows from \eqref{eq:explicit-exterior} that
\[
 \partial C_m
 =\bigl(\mathbb S^1\times\partial D^{m+1}\bigr)
  \sqcup
  \mathcal F_{\rhoCab}
  \bigl((\mathbb R/2\pi\mathbb Z)\times
                         \partial D^{m+1}_{\rhoCab}\bigr)
 =\partial_oC_m\sqcup\partial_iC_m,
\]
and
\[
 \partial_oC_m\cong\mathbb S^1\times\mathbb S^m,
 \qquad
 \partial_iC_m\cong\mathbb S^1\times\mathbb S^m.
\]

\emph{Part \textup{(ii)}.}
Define
\[
\mathcal D_\tau:V\setminus\operatorname{im}c_2\longrightarrow
                  V\setminus\operatorname{im}c_2,
 \qquad 0\le\tau\le1,
\]
by
\[
\mathcal D_\tau(y)=
 \begin{cases}
 \displaystyle
 \mathcal F_{\rhoCab}\left(\theta,
  \left((1-\tau)|v|+\tau\rhoCab\right)\frac{v}{|v|}\right),
 &y=\mathcal F_{\rhoCab}(\theta,v),\quad0<|v|\le\rhoCab,\\[3mm]
 y,&y\in C_m.
 \end{cases}
\]
On $|v|=\rhoCab$ the two formulas agree, so $\mathcal D_\tau$ is
continuous.  Hence $\mathcal D_\tau$ is a strong deformation retraction
of $V\setminus\operatorname{im}c_2$ onto $C_m$; explicitly,
\begin{align*}
 \mathcal D_0&=\operatorname{id}_{V\setminus\operatorname{im}c_2},\\
 \mathcal D_\tau|_{C_m}&=\operatorname{id}_{C_m},\\
 \mathcal D_1\bigl(V\setminus\operatorname{im}c_2\bigr)&=C_m.
\end{align*}

Let $[\gamma]\in\pi_1(V,x_O)$.  Since $x_O\notin\operatorname{im}c_2$,
relative transversality gives
$\gamma'\simeq\gamma$ relative to $x_O$ with
$\gamma'\pitchfork\operatorname{im}c_2$.  Therefore
\[
 \dim (\gamma')^{-1}\bigl(\operatorname{im}c_2\bigr)
 =1+1-(m+2)=-m<0,
\]
so
\[
 (\gamma')^{-1}\bigl(\operatorname{im}c_2\bigr)=\varnothing,
 \qquad
 \mathcal D_1\circ\gamma':\mathbb S^1\longrightarrow C_m.
\]
Thus $j_*$ is surjective.

If $[\gamma]\in\ker j_*$, there is a map
\[
 H:(D^2,\partial D^2)\longrightarrow(V,C_m),
 \qquad H|_{\partial D^2}=\gamma.
\]
The boundary of $H$ is contained in $C_m$ and is therefore disjoint from
$\operatorname{im}c_2$.  Relative transversality gives
$H'\simeq H$ relative to $\partial D^2$ with
$H'\pitchfork\operatorname{im}c_2$.  Hence
\[
 \dim (H')^{-1}\bigl(\operatorname{im}c_2\bigr)
 =2+1-(m+2)=1-m<0,
\]
and therefore
\[
 (H')^{-1}\bigl(\operatorname{im}c_2\bigr)=\varnothing.
\]
The map $\mathcal D_1\circ H':D^2\to C_m$ satisfies
\[
 (\mathcal D_1\circ H')|_{\partial D^2}=\gamma,
\]
so $[\gamma]=1$ in $\pi_1(C_m,x_O)$.  Thus $j_*$ is injective.
Since
\[
 V\simeq\mathbb S^1,
 \qquad
 \pi_1(V,x_O)\cong\mathbb Z,
\]
we obtain
\[
 j_*:\pi_1(C_m,x_O)\xrightarrow{\ \cong\ }\pi_1(V,x_O),
 \qquad
 \pi_1(C_m,x_O)\cong\mathbb Z.
\]

\emph{Part \textup{(iii)}.}
The equation $e^{2\mathrm i\theta}=1$ has the two solutions
$\theta=0,\pi$ modulo $2\pi$.  Consequently,
\[
 \nu_{\rhoCab}(c_2)\cap(\{1\}\times D^{m+1})
 =\{1\}\times(D^{m+1}_+\sqcup D^{m+1}_-).
\]
Write
\[
 b(a)=a\varepsilon e_1+(1-a+a\rhoCab)e_{m+1},
 \qquad \beta(a)=(1,b(a)).
\]
Then
\[
 |b(a)|\le(1-a)+a(\varepsilon+\rhoCab)\le1,
\]
and
\begin{align*}
 |b(a)-\varepsilon e_1|^2
 &=\varepsilon^2(1-a)^2+(1-a+a\rhoCab)^2
 \ge\rhoCab^2,\\
 |b(a)+\varepsilon e_1|
 &\ge(1+a)\varepsilon>\rhoCab.
\end{align*}
Therefore
\[
 \beta([0,1])\subset C_m,
 \qquad
 \beta(0)=x_O,
 \qquad
 \beta(1)=x_I.
\]

Let
\[
 p_V:V\longrightarrow\mathbb S^1,
 \qquad p_V(z,x)=z,
 \qquad p=p_V|_{C_m}.
\]
Since $p_*=(p_V)_*\circ j_*$ and
$(p_V)_*:\pi_1(V,x_O)\to\pi_1(\mathbb S^1,1)$ is an isomorphism,
\[
 p_*:\pi_1(C_m,x_O)\xrightarrow{\ \cong\ }\mathbb Z.
\]
The identities
\begin{align*}
 (p\circ\lambda_O)(\varphi)&=e^{\mathrm i\varphi},\\
 p\circ\beta&\equiv1,\\
 (p\circ\lambda_I)(\theta)&=e^{2\mathrm i\theta}
\end{align*}
give
\begin{align*}
 p_*[\lambda_O]&=1,\\
 p_*\bigl[\beta*\lambda_I*\beta^{-1}\bigr]&=2.
\end{align*}
Consequently, for $g=[\lambda_O]$,
\[
 \bigl[\beta*\lambda_I*\beta^{-1}\bigr]=g^2.
\]
If $\widehat\beta$ is another path from $x_O$ to $x_I$, then
\[
 \bigl[\widehat\beta*\lambda_I*\widehat\beta^{-1}\bigr]
 =
 [\widehat\beta*\beta^{-1}]\,
 \bigl[\beta*\lambda_I*\beta^{-1}\bigr]\,
 [\widehat\beta*\beta^{-1}]^{-1}.
\]
The group $\pi_1(C_m,x_O)\cong\mathbb Z$ is abelian, so the right-hand
side equals $g^2$.

\emph{Part \textup{(iv)}.}
For $0\le\varphi\le2\pi$, the two solutions of
$e^{2\mathrm i\theta}=e^{\mathrm i\varphi}$ are
\[
 \theta=\frac\varphi2,
 \qquad
 \theta=\frac\varphi2+\pi
 \pmod{2\pi}.
\]
Since
\[
 \varepsilon u(\varphi/2)=R_{\varphi/2}(\varepsilon e_1),
 \qquad
 \varepsilon u(\varphi/2+\pi)=R_{\varphi/2}(-\varepsilon e_1),
\]
we have
\begin{align}
 \nu_{\rhoCab}(c_2)\cap
 \bigl(\{e^{\mathrm i\varphi}\}\times D^{m+1}\bigr)
 &=\{e^{\mathrm i\varphi}\}\times
 R_{\varphi/2}\bigl(D^{m+1}_+\sqcup D^{m+1}_-\bigr),
                                                   \label{eq:double-cable-tube-fiber}\\
 C_m\cap\bigl(\{e^{\mathrm i\varphi}\}\times D^{m+1}\bigr)
 &=\{e^{\mathrm i\varphi}\}\times R_{\varphi/2}P_m.
                                                   \label{eq:double-cable-exterior-fiber}
\end{align}
Furthermore,
\begin{align*}
 R_\pi(D^{m+1}_+)&=D^{m+1}_-,
 &R_\pi(D^{m+1}_-)&=D^{m+1}_+,\\
 \phi(I_+)&=I_-,
 &\phi(I_-)&=I_+.
\end{align*}
The map
\[
 \widetilde\Phi:P_m\times\mathbb R\longrightarrow C_m,
 \qquad
 \widetilde\Phi(x,\varphi)
   =(e^{\mathrm i\varphi},R_{\varphi/2}x),
\]
satisfies
\[
 \widetilde\Phi(x,\varphi+2\pi)
 =(e^{\mathrm i\varphi},R_{\varphi/2}R_\pi x)
 =\widetilde\Phi(\phi(x),\varphi).
\]
It therefore induces
\[
 \Phi:T_\phi\longrightarrow C_m,
 \qquad
 \Phi([x,\varphi])
   =(e^{\mathrm i\varphi},R_{\varphi/2}x).
\]
For fixed $\varphi$, the differential in the $P_m$-directions is the
isomorphism induced by $R_{\varphi/2}$, whereas the differential in the
$\varphi$-direction has nonzero component tangent to the first circle
factor.  Thus $\widetilde\Phi$, and hence $\Phi$, is a local
diffeomorphism.  Equation~\eqref{eq:double-cable-exterior-fiber} shows
that for $0\le\varphi<2\pi$ every point
$(e^{\mathrm i\varphi},y)\in C_m$ has the unique preimage
\[
 [R_{-\varphi/2}y,\varphi]\in T_\phi.
\]
Therefore $\Phi$ is bijective.  A bijective local diffeomorphism is a
diffeomorphism.
\end{proof}

The preceding proposition isolates the square-root relation in a form
adapted to the subsequent geometric construction.  Indeed, the
diffeomorphism in part~(iv) identifies \(C_m\) with the mapping torus of
\(\phi\colon P_m\to P_m\), where \(\phi\) interchanges the two inner
boundary components \(I_+\) and \(I_-\).  Under this identification, a
loop traversing the mapping-torus direction once represents the
generator \(\gamma\), whereas the distinguished longitude of the resulting
inner boundary represents \(\gamma^2\).  Thus, when regarded as a cobordism
from its inner boundary to its outer boundary, this topological model
replaces the distinguished inner loop representing \(\gamma^2\) by the
distinguished outer loop representing \(\gamma\).  It remains to realize this
model by a Ricci-positive metric whose induced boundary metrics and
second fundamental forms satisfy the hypotheses for successive
Ricci-positive gluings.



\section{Relative boundary tools}
\label{sec:relative-tools}
The construction of Section~\ref{sec:double-cable} determines the topology of
the basic block and the relation between its distinguished boundary
loops, but
it does not yet provide the boundary geometry needed for Ricci-positive
gluing.  To attach the basic blocks successively, we must be able to
modify a metric inside prescribed collar neighborhoods while controlling
both the induced boundary metric and the corresponding second
fundamental form.  Two such operations will be used repeatedly below:
gluing Ricci-positive manifolds along isometric boundary components, and
transporting a boundary metric along a path of Ricci-positive metrics
with controlled convexity at the two ends.  Both constructions must also
respect the relevant symmetries.

We develop these tools in this section.  Their proofs are based on the
Ricci-curvature formulas for a one-parameter family of metrics on the
slices of a collar.  Thus, consider a metric of the form
\[
g=ds^2+h(s),
\]
and let \(\nu=\partial_s\).  With the convention fixed above, the second
fundamental form of an \(s\)-slice with respect to \(\nu\) is
\[
L:=\II_{\nu}=\frac12\dot h.
\]
We begin by recording the corresponding Ricci-curvature formulas.


\begin{lemma}\label{lem:gaussian-ricci-unified}
For vectors $X,Y$ tangent to an $s$-slice,
\begin{align}
 \Ric_g(\partial_s,\partial_s)
 &=-\frac12\tr_h\ddot h
   +\frac14\tr(h^{-1}\dot h\,h^{-1}\dot h),\label{eq:gaussian-normal}\\
 \Ric_g(\partial_s,X)
 &=(\divsymb_hL)(X)-\dd(\tr_hL)(X),\label{eq:gaussian-mixed}\\
 \Ric_g(X,Y)
 &=\Ric_h(X,Y)-\frac12\ddot h(X,Y)
   +\frac12(\dot h\,h^{-1}\dot h)(X,Y)\notag\\
 &\hspace{31mm}
   -\frac14\tr_h(\dot h)\dot h(X,Y).\label{eq:gaussian-tangential}
\end{align}
\end{lemma}

\begin{proof}
Extend the tangent vectors to be independent of $s$ in product
coordinates.  Then the Christoffel symbols of $g$ satisfy
\[
 \Gamma^s_{ij}=-L_{ij},\qquad
 \Gamma^k_{si}=\Gamma^k_{is}=(h^{-1}L)^k{}_i,
 \qquad \Gamma^k_{ij}=\Gamma^k_{ij}(h).
\]
The Weingarten equation and the Gauss--Codazzi equation yield
\begin{align*}
 \Rm_g(\partial_s,X,Y,\partial_s)
   &=-\dot L(X,Y)+(Lh^{-1}L)(X,Y),\\
 \Ric_g(\partial_s,X)&=(\divsymb_hL)(X)-\dd(\tr_hL)(X),\\
 \Rm_g(X,Y,Z,W)&=\Rm_h(X,Y,Z,W)
 -L(X,W)L(Y,Z)+L(X,Z)L(Y,W).
\end{align*}
Taking the $h$-trace of the first identity gives
\eqref{eq:gaussian-normal}.  The second identity is
\eqref{eq:gaussian-mixed}; combining the first identity with the
$h$-trace of the Gauss equation gives
\[
 \Ric_g(X,Y)=\Ric_h(X,Y)-\dot L(X,Y)
 +2(Lh^{-1}L)(X,Y)-(\tr_hL)L(X,Y).
\]
Substituting $L=\tfrac12\dot h$ yields
\eqref{eq:gaussian-tangential}.
\end{proof}
We next reprove Perelman's Ricci-positive gluing theorem
\cite[Section~4]{Perelman1997} in a form adapted to our later
applications, keeping track of the support of the metric modification
and the equivariance of the construction.  A more general form of this
gluing theorem is proved in \cite{ReiserWraithGluing}.
\begin{lemma}
\label{lem:relative-gluing-unified}
Let $(M_i,g_i)$, $i=1,2$, be compact Riemannian manifolds with boundary
and $\Ric_{g_i}>0$.  For each $i$, let $N_i$ be a union of connected
components of $\partial M_i$, and let
\[
 \Psi:(N_1,g_1|_{TN_1})\longrightarrow(N_2,g_2|_{TN_2})
\]
be an isometry.  Let $\II_i$ denote the second fundamental form of $N_i$
with respect to the outward unit normal of $M_i$.  Suppose that
$\II_1+\Psi^*\II_2$ is positive definite on $TN_1$.

Then for every choice of collar neighborhoods $U_i\subset M_i$ of $N_i$,
the smooth gluing manifold $M_1\cup_\Psi M_2$ admits a Riemannian metric
$\widetilde g$ such that
\[
 \Ric_{\widetilde g}>0,
 \qquad \widetilde g|_{M_i\setminus U_i}=g_i|_{M_i\setminus U_i}
 \quad(i=1,2).
\]
Suppose, in addition, that a finite group $\Gamma$ acts isometrically on
each $(M_i,g_i)$, preserves $N_i$, and satisfies
$\Psi(\gamma x)=\gamma\Psi(x)$ for all $\gamma\in\Gamma$ and
$x\in N_1$.  Then $\widetilde g$ can be chosen to be $\Gamma$-invariant.  The
same conclusions hold for finitely many pairwise disjoint boundary
identifications, with the modifications contained in prescribed
pairwise disjoint collar neighborhoods.
Moreover, for every \(\varepsilon_{\mathrm{vol}}>0\), the collar
neighborhoods and the smoothing parameters can be chosen so that
\begin{equation}\label{eq:volume pinching}
 \left|
 \operatorname{Vol}_{\widetilde g}(M_1\cup_\Psi M_2)
 -\operatorname{Vol}_{g_1}(M_1)
 -\operatorname{Vol}_{g_2}(M_2)
 \right|<\varepsilon_{\mathrm{vol}}.
\end{equation}
For finitely many pairwise disjoint identifications, the sum of the
absolute volume changes can be required to be less than
\(\varepsilon_{\mathrm{vol}}\).
\end{lemma}

\begin{proof}
We first introduce one identification to simplify the argument.  Use $\Psi$ to identify
$N_1$ and $N_2$ as one closed manifold $N$, and choose a signed Fermi
coordinate $s$ that increases from $M_1$ to $M_2$.  After pulling back
the tensors on $N_2$ by $\Psi$, the metrics have the form
\[
 g=\dd s^2+h_-(s)\quad(s\le0),\qquad
 g=\dd s^2+h_+(s)\quad(s\ge0),\qquad h_-(0)=h_+(0).
\]
Because the outward normals are $+\partial_s$ and $-\partial_s$,
\begin{equation}
 \dot h_-(0)-\dot h_+(0)
   =2(\II_1+\Psi^*\II_2)>0.                    \label{eq:normal-derivative-gap}
\end{equation}

For sufficiently small $\varepsilon>0$, let $h_\varepsilon(s)$ be the
unique cubic polynomial in $s$, with coefficients in
$\Gamma(\Sym^2T^*N)$, satisfying
\[
 \begin{aligned}
 h_\varepsilon(-\varepsilon)&=h_-(-\varepsilon),&
 \dot h_\varepsilon(-\varepsilon)&=\dot h_-(-\varepsilon),\\
 h_\varepsilon(\varepsilon)&=h_+(\varepsilon),&
 \dot h_\varepsilon(\varepsilon)&=\dot h_+(\varepsilon).
 \end{aligned}
\]
Taylor expansion at $s=0$ gives, uniformly for
$|s|\le\varepsilon$,
\begin{equation}
 \begin{aligned}
 h_\varepsilon(s)&=h_-(0)+O(\varepsilon),\\
 \dot h_\varepsilon(s)&=O(1),\\
 \ddot h_\varepsilon(s)
   &=-\frac{\dot h_-(0)-\dot h_+(0)}{2\varepsilon}+O(1).
 \end{aligned}                                      \label{eq:hermite-estimates}
\end{equation}
 Here the three estimates are taken in the $C^2(N)$, $C^1(N)$, and $C^0(N)$ senses,
respectively.  In particular, $h_\varepsilon(s)$ is positive definite
when $\varepsilon$ is small.

Replace the two original metrics on $|s|\le\varepsilon$ by
$\dd s^2+h_\varepsilon(s)$, and retain the original metrics for
$|s|\ge\varepsilon$.  The resulting metric $g_\varepsilon$ is $C^1$
and is smooth away from $s=\pm\varepsilon$.  Equations
\eqref{eq:gaussian-normal}--\eqref{eq:gaussian-tangential} and
\eqref{eq:hermite-estimates} give, on $|s|<\varepsilon$,
\begin{align}
 \Ric_{g_\varepsilon}(\partial_s,\partial_s)
  &=\frac{1}{4\varepsilon}
    \tr_{h_\varepsilon}\!\bigl(\dot h_-(0)-\dot h_+(0)\bigr)+O(1),
    \notag\\
 \left.\Ric_{g_\varepsilon}\right|_{TN\times TN}
  &=\frac{\dot h_-(0)-\dot h_+(0)}{4\varepsilon}+O(1),
    \label{eq:interpolated-ricci}\\
 \Ric_{g_\varepsilon}(\partial_s,\cdot)&=O(1).\notag
\end{align}
All three estimates in \eqref{eq:interpolated-ricci} are uniform on $[-\varepsilon,\varepsilon]\times N$.
Set $Q_\varepsilon:=\left.\Ric_{g_\varepsilon}\right|_{TN\times TN}$ and
$\beta_\varepsilon:=\Ric_{g_\varepsilon}(\partial_s,\mathord\cdot)|_{TN}$.
By \eqref{eq:normal-derivative-gap}, compactness of $N$, and the uniform
equivalence of $h_\varepsilon$ and $h_-(0)$, there are constants $c,C>0$,
independent of sufficiently small $\varepsilon$, such that
\[
 \Ric_{g_\varepsilon}(\partial_s,\partial_s)
    \ge \frac{c}{\varepsilon}-C,
 \qquad
 Q_\varepsilon\ge\left(\frac{c}{\varepsilon}-C\right)h_\varepsilon,
 \qquad
 |\beta_\varepsilon|_{h_\varepsilon}\le C.
\]
After decreasing $\varepsilon$, the first two lower bounds are at least
$c/(2\varepsilon)$.  In particular, $Q_\varepsilon$ is positive definite
and, as a quadratic form on $T^*N$,
\[
 Q_\varepsilon^{-1}\le \frac{2\varepsilon}{c}\,h_\varepsilon^{-1}.
\]
Consequently, the Schur complement of $Q_\varepsilon$ in the decomposition
$\mathbb R\partial_s\oplus TN$ satisfies
\[
 \Ric_{g_\varepsilon}(\partial_s,\partial_s)
 -Q_\varepsilon^{-1}(\beta_\varepsilon,\beta_\varepsilon)
 \ge \frac{c}{2\varepsilon}-\frac{2C^2\varepsilon}{c}>0
\]
for all sufficiently small $\varepsilon$.  Since $Q_\varepsilon>0$, the
block-matrix criterion for positive definiteness now shows that
$\Ric_{g_\varepsilon}>0$ wherever $g_\varepsilon$ is smooth.

It remains to smooth $g_\varepsilon$ at the two hypersurfaces
$s=\pm\varepsilon$.  After translating either one to $t=0$, we write
$g_\varepsilon=\dd t^2+k(t)$, where $k$ is $C^1$ and is smooth on each
of the half-intervals $t\le0$ and $t\ge0$.  The two one-sided Ricci
tensors extend continuously to positive-definite tensors at $t=0$;
after restricting the interval, they therefore have a common positive
lower bound.

Choose a nonnegative even function
$\varphi\in C_c^\infty((-1,1))$ with $\int\varphi\dd t=1$, set
$\varphi_\eta(t)=\eta^{-1}\varphi(t/\eta)$, and define
\[
 k_\eta(t)=\int_{\mathbb R}\varphi_\eta(r)k(t-r)\dd r.
\]
Since $k$ is $C^1$, its distributional second derivative $D^2 k$ has the expression
\[
D^2 k=\ddot{k}_{\mathrm{pw}},
\]
where
\[
\ddot{k}_{\mathrm{pw}}(t)= \begin{cases}\ddot{k}_{-}(t), & t<0, \\ \ddot{k}_{+}(t), & t>0 .\end{cases}
\]
Also,
\begin{equation}
 \ddot k_\eta(t)
   =\int_{\mathbb R}\varphi_\eta(r)\ddot{k}_{\mathrm{pw}}(t-r)\dd r.
\label{eq:mollified-second-derivative}
\end{equation}
The formulas in Lemma~\ref{lem:gaussian-ricci-unified} are affine in
$\ddot k$, $k_\eta$ and $\dot k_\eta$ converge
uniformly to $k$ and $\dot k$ as $\eta\to0$.  Consequently,
\begin{equation}
 \Ric_{\dd t^2+k_\eta(t)}
 =\int_{\mathbb R}\varphi_\eta(r)
      \Ric_{\dd t^2+k(t-r)}\dd r+o_\eta(1),
                                                        \label{eq:mollified-ricci}
\end{equation}
uniformly near $t=0$.  The integral on the right is positive definite
with a uniform lower bound, so $\dd t^2+k_\eta(t)$ has positive Ricci
curvature near $t=0$ when $\eta$ is sufficiently small.

To make the metric unchanged outside a smaller collar, choose
$\chi_\eta\in C_c^\infty((-\sqrt\eta,\sqrt\eta))$ with
$0\le\chi_\eta\le1$, $\chi_\eta=1$ on $[-2\eta,2\eta]$, and
\[
 |\dot\chi_\eta|\le C\eta^{-1/2},
 \qquad |\ddot\chi_\eta|\le C\eta^{-1}.
\]
Set
\[
 \widetilde k_\eta
   =k+\chi_\eta(k_\eta-k).
\]
On $2\eta\le|t|\le\sqrt\eta$, the convolution defining $k_\eta$ uses
only one smooth half-interval.  The evenness of $\varphi$ and Taylor's
expansion imply
\[
 \|\partial_t^j(k_\eta-k)\|_{C^{2-j}(N)}=O(\eta^2),
 \qquad j=0,1,2.
\]
It follows that $\widetilde k_\eta-k=O(\eta)$ in $C^2$ on this region.
Since positive Ricci curvature is open in the $C^2$ topology,
$\dd t^2+\widetilde k_\eta(t)$ has positive Ricci curvature there for
small $\eta$.  It agrees with $\dd t^2+k_\eta(t)$ for $|t|\le2\eta$
and with $\dd t^2+k(t)$ for $|t|\ge\sqrt\eta$, so it gives the required
smooth local replacement.

Choose the initial Fermi collars inside $U_1$ and $U_2$, and perform the
preceding smoothing argument at both hypersurfaces $s=\pm\varepsilon$.  This
produces a smooth metric $\widetilde g$ satisfying
$\widetilde g=g_i$ on $M_i\setminus U_i$.

Under the hypotheses involving $\Gamma$, the normal exponential maps
give $\Gamma$-equivariant Fermi collars.  Hermite interpolation,
convolution in the normal coordinate, and multiplication by the scalar
function $\chi_\eta$ commute with the $\Gamma$-action.  Thus
$\widetilde g$ is $\Gamma$-invariant.  For finitely many pairwise
disjoint identifications, carry out the construction independently in
the prescribed pairwise disjoint collars.

It remains to verify \eqref{eq:volume pinching}. The first estimate in
\eqref{eq:hermite-estimates} shows that the interpolated tangential
metrics converge uniformly to the common boundary metric while their
supports shrink to \(N\).  At the two auxiliary hypersurfaces, the
mollified tensors and the final cutoff tensors also converge uniformly
to the piecewise smooth tensor.  Hence the associated volume densities
are uniformly bounded and converge uniformly on their supports, while
the widths of all modified collars tend to zero.  The absolute change
of volume therefore tends to zero.  Choose the collar widths and the
smoothing parameters so that this change is less than
\(\varepsilon_{\mathrm{vol}}\).  In the case of finitely many disjoint
identifications, assign a positive volume bound to each identification
whose sum is less than \(\varepsilon_{\mathrm{vol}}\).
\end{proof}

We next prove an ``isotopy implies concordance'' theorem for metrics of
positive Ricci curvature.  It realizes a path of Ricci-positive metrics
as a Ricci-positive cylinder whose induced boundary metrics agree with
the endpoints up to scaling and whose outward second fundamental forms
satisfy the required one-sided bounds.  The argument follows
\cite[Theorem~A.1]{Burdick} and \cite[Assertion]{Perelman1997}.

\begin{theorem}
\label{thm:path-cylinder-unified}
Let $N^n$ be a closed manifold, and let $(g_s)_{s\in[0,1]}$ be a smooth
path of Riemannian metrics satisfying $\Ric_{g_s}>0$.  Suppose that the
path is constant on a neighborhood of each endpoint.  For every
$\delta_{\mathrm{cyl}}>0$ there are $\Lambda_{\mathrm{cyl}}>0$ and a
smooth Riemannian metric $\widehat g$ on
$[0,1]\times N$ with $\Ric_{\widehat g}>0$ such that:
\begin{enumerate}[label=\textup{(\roman*)}]
\item the metric induced by $\widehat g$ on $\{0\}\times N$ is $g_0$,
and the second fundamental form with respect to the outward unit normal
is greater than $-\delta_{\mathrm{cyl}}g_0$;
\item the metric induced by $\widehat g$ on $\{1\}\times N$ is
$\Lambda_{\mathrm{cyl}}^2g_1$, and the second fundamental form with respect to the outward
unit normal is positive definite.
\end{enumerate}
Moreover, the construction is equivariant under every compact group of
diffeomorphisms of $N$ that preserves each metric $g_s$.
\end{theorem}

\begin{proof}
Compactness of $[0,1]\times N$ gives constants $\mu,C>0$ such that,
after increasing $C$ when necessary,
\begin{equation}
 \Ric_{g_s}\ge\mu g_s,
 \qquad
 |\partial_sg_s|_{g_s}
 +|\nabla^{g_s}(\partial_sg_s)|_{g_s}
 +|\partial_s^2g_s|_{g_s}\le C
 \quad (0\le s\le1).
 \label{eq:path-uniform-bounds}
\end{equation}
Choose $a>0$ and $\sigma>0$ so that
\begin{equation}
 a<\delta_{\mathrm{cyl}},
 \qquad \frac{\mu}{a^2}>2n,
 \qquad
 n\sigma-\frac12\sup_{[0,1]\times N}
   |\tr_{g_s}(\partial_sg_s)|\ge1.
 \label{eq:path-parameter-choice}
\end{equation}

For a sufficiently large number $\ell$, consider
$[e^\ell,e^{2\ell}]\times N$ and define
\begin{equation}
 \lambda(t)=2\left(1-\frac{\ell}{\log t}\right),
 \qquad
 f(t)=at\,e^{-\sigma\lambda(t)},
 \qquad
 \overline g=\dd t^2+f(t)^2g_{\lambda(t)}.
 \label{eq:path-cylinder-metric}
\end{equation}
Then $\lambda(e^\ell)=0$, $\lambda(e^{2\ell})=1$, and
\begin{equation}
 \lambda'(t)=\frac{2\ell}{t(\log t)^2},
 \qquad
 \lambda''(t)=-\frac{\lambda'(t)}t
                   \left(1+\frac2{\log t}\right),
 \qquad
 \frac{f'(t)}{f(t)}=\frac1t-\sigma\lambda'(t).
 \label{eq:path-cylinder-derivatives}
\end{equation}
The shape endomorphism of a $t$-slice with respect to $\partial_t$ is
therefore
\begin{equation}
 \frac12\bigl(f^2g_{\lambda(t)}\bigr)^{-1}
       \partial_t\bigl(f^2g_{\lambda(t)}\bigr)
 =\left(\frac1t-\sigma\lambda'(t)\right)\Id
  +\frac{\lambda'(t)}2g_{\lambda(t)}^{-1}
       \left.\partial_sg_s\right|_{s=\lambda(t)}.
 \label{eq:path-slice-shape}
\end{equation}

Substitute \eqref{eq:path-cylinder-derivatives} and
\eqref{eq:path-slice-shape} directly into the three formulae of
Lemma~\ref{lem:gaussian-ricci-unified}.  The normal component is
\begin{align}
 \Ric_{\overline g}(\partial_t,\partial_t)
 ={}&\left(n\sigma-\frac12\tr_{g_{\lambda(t)}}
       \left(\left.\partial_sg_s\right|_{s=\lambda(t)}\right)\right)
       \frac{\lambda'(t)}t\left(1-\frac2{\log t}\right)
       +O\bigl(\lambda'(t)^2\bigr).
 \label{eq:path-normal-ricci}
\end{align}
The constant in the error term is independent of $\ell$.  Since
$\log t\in[\ell,2\ell]$, 
\eqref{eq:path-parameter-choice} and \eqref{eq:path-normal-ricci} give
\begin{equation}
 \Ric_{\overline g}(\partial_t,\partial_t)
 \ge\frac1{8t^2\ell}
 \label{eq:path-normal-lower}
\end{equation}
when $\ell$ is sufficiently large.

For the tangential component, use $f(t)^2g_{\lambda(t)}$ to raise one
index in \eqref{eq:gaussian-tangential}.  Since
$\lambda'(t)=O((t\ell)^{-1})$ and
$\lambda''(t)=O((t^2\ell)^{-1})$, 
\eqref{eq:path-uniform-bounds} and
\eqref{eq:path-cylinder-derivatives} give
\begin{align*}
 &\bigl(f(t)^2g_{\lambda(t)}\bigr)^{-1}
   \partial_t\bigl(f(t)^2g_{\lambda(t)}\bigr)
   =\frac2t\Id+O\left(\frac1{t\ell}\right),\\
 &\bigl(f(t)^2g_{\lambda(t)}\bigr)^{-1}
   \partial_t^2\bigl(f(t)^2g_{\lambda(t)}\bigr)
   =\frac2{t^2}\Id+O\left(\frac1{t^2\ell}\right),\\
 &\tr_{f(t)^2g_{\lambda(t)}}
   \partial_t\bigl(f(t)^2g_{\lambda(t)}\bigr)
   =\frac{2n}{t}+O\left(\frac1{t\ell}\right).
\end{align*}
The endomorphism-valued error terms are measured with respect to
$g_{\lambda(t)}$, and all constants are independent of $t$ and $\ell$.
Thus the last three terms of \eqref{eq:gaussian-tangential}, after one
index is raised, equal
\[
 -\frac1{t^2}\Id+\frac2{t^2}\Id-\frac n{t^2}\Id
 +O\left(\frac1{t^2\ell}\right)
 =-\frac{n-1}{t^2}\Id+O\left(\frac1{t^2\ell}\right).
\]
Moreover, \eqref{eq:path-uniform-bounds} implies
\[
 \bigl(f(t)^2g_{\lambda(t)}\bigr)^{-1}
 \Ric_{g_{\lambda(t)}}\ge\frac{\mu}{f(t)^2}\Id.
\]
After increasing $C$, we therefore obtain
\begin{align}
 \left.\Ric_{\overline g}\right|_{TN\times TN}
 &\ge\left(\frac{\mu}{f(t)^2}
       -\frac{n-1+C/\ell}{t^2}\right)f(t)^2g_{\lambda(t)}
   \ge\frac n{t^2}f(t)^2g_{\lambda(t)}.
 \label{eq:path-tangential-lower}
\end{align}
Indeed, $f(t)\le at$ gives
\[
 \frac{\mu}{f(t)^2}-\frac{n-1+C/\ell}{t^2}
 \ge\frac1{t^2}
 \left(\frac{\mu}{a^2}-n+1-\frac C\ell\right)
 \ge\frac n{t^2},
\]
where the last inequality follows from
$\mu/a^2>2n$ once $\ell$ is large enough that $C/\ell<1$.
The mixed component has the exact expression
\begin{align}
 \Ric_{\overline g}(\partial_t,\mathord\cdot)
 =\frac{\lambda'(t)}2\Bigl(&
 \divsymb_{g_{\lambda(t)}}
   \left(\left.\partial_sg_s\right|_{s=\lambda(t)}\right)-\dd\tr_{g_{\lambda(t)}}\left(\left.\partial_sg_s\right|_{s=\lambda(t)}\right)\Bigr),
 \label{eq:path-mixed-exact}
\end{align}
because multiplying a slice metric by the spatially constant factor
$f(t)^2$ does not change its Levi-Civita connection.  Consequently, we have
\begin{equation}
 |\Ric_{\overline g}(\partial_t,\mathord\cdot)|_{f(t)^2g_{\lambda(t)}}
 \le\frac{C}{t^2\ell}.
 \label{eq:path-mixed-bound}
\end{equation}

By \eqref{eq:path-tangential-lower}, the inverse of the tangential Ricci
block, viewed as a quadratic form on $T^*N$, is bounded above by
 $\frac{t^2}{n}(f(t)^2g_{\lambda(t)})^{-1}$.  Thus
\begin{align}
 &\Ric_{\overline g}(\partial_t,\partial_t)
 -\left(\left.\Ric_{\overline g}\right|_{TN\times TN}\right)^{-1}
  \bigl(\Ric_{\overline g}(\partial_t,\mathord\cdot),
        \Ric_{\overline g}(\partial_t,\mathord\cdot)\bigr)\ge
 \frac1{8t^2\ell}-\frac{C^2}{nt^2\ell^2}>0
 \label{eq:path-schur-complement}
\end{align}
for sufficiently large $\ell$.  Together with
\eqref{eq:path-tangential-lower}, it implies that
$\Ric_{\overline g}>0$.

Now set $\widehat g=(ae^\ell)^{-2}\overline g$.  Because
$\partial_sg_s=0$ at $s=0,1$, the induced metrics and outward second
fundamental forms at the two boundary components are
\begin{align*}
 \left.\widehat g\right|_{\{e^\ell\}\times N}
 &=g_0,&
 \II^{\widehat g}_{\{e^\ell\}\times N}
 &=-a\left(1-\frac{2\sigma}{\ell}\right)g_0,\\
 \left.\widehat g\right|_{\{e^{2\ell}\}\times N}
 &=\Lambda_{\mathrm{cyl}}^2g_1,&
 \II^{\widehat g}_{\{e^{2\ell}\}\times N}
 &=ae^{-\ell}\left(1-\frac{\sigma}{2\ell}\right)
       \Lambda_{\mathrm{cyl}}^2g_1,
 \qquad \Lambda_{\mathrm{cyl}}=e^{\ell-\sigma}.
\end{align*}
For $\ell>2\sigma$, these satisfy the two required inequalities because
$a<\delta_{\mathrm{cyl}}$.  Pulling back $\widehat g$ by an
orientation-preserving diffeomorphism from $[0,1]$ to
$[e^\ell,e^{2\ell}]$ gives \textup{(i)} and \textup{(ii)}.

Finally, if a compact group of diffeomorphisms preserves every $g_s$,
then it preserves \eqref{eq:path-cylinder-metric} and hence also the
rescaled and pulled-back metric.
\end{proof}

We call the Riemannian cylinder $([0,1]\times N,\widehat g)$ furnished
by Theorem~\ref{thm:path-cylinder-unified} a \emph{Ricci-positive path
cylinder} associated with the path $(g_s)_{s\in[0,1]}$.



\section{The Ricci-positive directed base}
\label{sec:unified-directed-base}

Recall that the directed base is the tuple
$(P_m,\sigma_m;\partial_oP_m,I_+,I_-)$, where $\partial_oP_m$ is the
outgoing boundary and the involution $\sigma_m$ exchanges the two incoming
boundaries $I_+$ and $I_-$.  The purpose of this section is to equip this
fixed topological object with the Ricci-positive metric and boundary
geometry required by the screw quotient defined in (\ref{def:screw-quotient}).
We first specify the model involutions on the round boundary spheres.
We then introduce an equivariant angular-neck construction, whose proof
is deferred to Appendix~\ref{app:perelman-neck}, and use it to normalize
the two incoming boundary components.  Finally, an equivariant
Ricci-positive path cylinder rounds the outgoing boundary.
These ingredients are assembled in
Theorem~\ref{thm:unified-directed-base}.  Throughout the section, fix
\[
m\in\{2,3\},
\qquad
k:=m-1,
\]
and let \(g_{\mathbb S^k}\) denote the unit-round metric on
\(\mathbb S^k\).  The cases \(m=2\) and \(m=3\) produce bases of
dimensions three and four, respectively.  In either case, the
distinguished boundary hypersurfaces have dimension \(m\), while the
auxiliary round factor has dimension \(k\).  The principal ingredients
are the relative equivariant Ricci-positive gluing result,
Lemma~\ref{lem:relative-gluing-unified}, and the path-cylinder
construction, Theorem~\ref{thm:path-cylinder-unified}, established in
Section~\ref{sec:relative-tools}.

Identify $\mathbb R^3=\mathbb C\oplus\mathbb R$ and
$\mathbb R^4=\mathbb C^2$, and define skew-symmetric endomorphisms by
\[
 K_2(z,u)=(\mathrm i z,0),\qquad
 K_3(z_1,z_2)=(\mathrm i z_1,\mathrm i z_2).
\]
For $x\in\mathbb S^m$, the vector $K_mx$ is tangent to $\mathbb S^m$,
and is actually a Killing vector field on $\mathbb S^m$.  Its
time-$\tau$ flow is the restriction to $\mathbb S^m$ of the orthogonal
linear map $\exp(\tau K_m)$, where $\exp$ denotes the matrix
exponential.  In particular,
\begin{equation}
\exp(\pi K_2)(z,u)=(-z,u),\qquad
 \exp(\pi K_3)(z_1,z_2)=(-z_1,-z_2).                            \label{eq:unified-base-3-1}
\end{equation}
Thus $\exp(\pi K_2)$ is rotation through angle $\pi$ about the axis
$\{0\}\oplus\mathbb R\subset\mathbb C\oplus\mathbb R$ and fixes the two points
$(0,\pm1)\in\mathbb S^2$, whereas $\exp(\pi K_3)$ is the antipodal map
of $\mathbb S^3$.



These order-two isometries will be the prescribed actions on the round
outer boundary.  We next construct a local Ricci-positive model carrying
the same symmetry.  The local model is expressed in terms of angular
distance from the collapsed axis.  Its level sets provide the inner
hypersurfaces to which the equivariant necks will be attached.

The auxiliary function \(\mathcal R\) is stated in
Lemma~\ref{lem:concave-warping}; its construction is given in
Appendix~\ref{app:perelman-neck}.  Its curvature inequalities produce
a translation-invariant family of angular \(m\)-spheres with positively
curved induced metrics and controlled outward second fundamental forms.
The following definition specifies these spheres, and
Proposition~\ref{prop:unified-neck} connects each of them to a round
sphere by an equivariant Ricci-positive neck.

\begin{definition}
On
\[\mathcal Q:=\bigl([0,\pi/2)\times\mathbb R\times\mathbb S^k\bigr)/\!\sim,
\]
where each set $\{0\}\times\{z\}\times\mathbb S^k$ is collapsed to one
point,
denote by $o_{z_0}$ the point represented by
$\{0\}\times\{z_0\}\times\mathbb S^k$.  On the strip
$|z-z_0|<\pi$, define
\[
 \varrho_{z_0}(t,z)
 :=\arccos\bigl(\cos t\cos(z-z_0)\bigr).
\]
For $0<r<\pi/2$, the sets
\[
 B_r^{\angle}(o_{z_0}):=\{\varrho_{z_0}\le r\},
 \qquad
 S_r^{\angle}(o_{z_0}):=\{\varrho_{z_0}=r\}
\]
are called, respectively, the angular ball and the angular sphere of
radius $r$ centered at $o_{z_0}$.
\end{definition}
We have the following proposition for constructing the neck metric by Perelman \cite[Assertion]{Perelman1997} and \cite[Section 3]{Perelman1997}. The proof of Proposition~\ref{prop:unified-neck} is given in
Appendix~\ref{app:perelman-neck} for the sake of completeness.
\begin{proposition}
\label{prop:unified-neck}
There is $r_*>0$ with the following property.  Fix
$m\in\{2,3\}$, let $k=m-1$, and choose
\[
             0<r<r_*,\qquad 0<\rho<1,\qquad \pi\rho<r.
\]
For every sufficiently small $w>0$ satisfying
\begin{equation}
w^{m-1}<\rho^m,                            \label{eq:unified-base-3-10}
\end{equation}
there is a function $\mathcal R$ satisfying
Lemma~\ref{lem:concave-warping} for which the following assertions hold
on the subset $\{0\le t\le5r\}$ of $\mathcal Q$.

\begin{enumerate}[label=\textup{(\roman*)}]
\item On a neighborhood of $B_{4r}^{\angle}(o_0)$, the metric
\begin{equation}
g_{\mathrm{amb},m}=\cot^2r\left(
 \dd t^2+\cos^2t\,\dd z^2+\mathcal R(t)^2g_{\mathbb S^k}\right)         \label{eq:unified-base-3-11}
\end{equation}
is smooth and satisfies $\sec_{g_{\mathrm{amb},m}}>0$.

\item For every $z_0\in\mathbb R$, the angular sphere
$S_r^{\angle}(o_{z_0})$ is diffeomorphic to $\mathbb S^m$.  Define its
latitude coordinate $\varphi\in[-\pi/2,\pi/2]$ by
\[
 \cos\varphi=\frac{\mathcal R(t)}{\mathcal R(r)},\qquad
 \sgn\varphi=\sgn(z-z_0).
\]
In this coordinate its induced metric is independent of $z_0$ and has
the form
\begin{equation}
g_{\mathrm{ang},m}=A_{\mathrm{ang}}(\varphi)^2\dd\varphi^2
           +w^2\cos^2\varphi\,g_{\mathbb S^k}.                   \label{eq:unified-base-3-12}
\end{equation}
 Here $A_{\mathrm{ang}}$ is smooth and positive, with
 $A_{\mathrm{ang}}(\pm\pi/2)=w$, and $\sec_{g_{\mathrm{ang},m}}>1$.  With respect to the normal
pointing out of
$B_r^{\angle}(o_{z_0})$, every principal curvature of
$S_r^{\angle}(o_{z_0})$ lies in $(0,1]$.

\item There are $\ell,\lambda>0$ and a Ricci-positive metric
\begin{equation}
g_{\mathrm{neck},m}=\dd s^2+A(s,\varphi)^2\dd\varphi^2
          +B(s)^2\cos^2\varphi\,g_{\mathbb S^k}                 \label{eq:unified-base-3-13}
\end{equation}
on $[0,\ell]\times\mathbb S^m$.  At $s=0$, its induced metric is
$(\rho/\lambda)^2g_{\mathbb S^m}$ and all principal curvatures with
respect to the outward normal $-\partial_s$ equal $-\lambda$.  At
$s=\ell$, its induced metric is $g_{\mathrm{ang},m}$ and all principal curvatures
with respect to $\partial_s$ are greater than one.  Moreover,
$B'(s)>0$ and $B''(s)<0$ for $0\le s\le\ell$.
\end{enumerate}
The construction commutes with translations of the ambient $z$
coordinate and orthogonal transformations of 
$\mathbb S^k$.
\end{proposition}






The neck construction normalizes the two incoming boundary components,
but the metric induced on $\partial_oP_m$ is merely positively curved.
To replace it by an exactly round metric without losing either Ricci
positivity or the involutive symmetry, we use the following equivariant
rounding lemma.

\begin{lemma}
\label{lem:unified-round-path}
Let $m\in\{2,3\}$, let $\gamma$ be a metric of positive sectional
curvature on $\mathbb S^m$, and let a finite group $\Gamma$ act isometrically
on $(\mathbb S^m,\gamma)$.  There is a smooth path
$(\gamma_s)_{s\in[0,1]}$ of $\Gamma$-invariant metrics such that
\[
 \gamma_0=\gamma,\qquad \Ric_{\gamma_s}>0\quad(0\le s\le1),
 \qquad \gamma_1\text{ has constant positive sectional curvature}.
\]
The path can be chosen constant for $s$ near $0$ and near $1$.
\end{lemma}

\begin{proof}
Suppose first that $m=2$.  The main convergence theorem in Hamilton's
paper on the Ricci flow of surfaces \cite{Hamilton1988Surfaces} states
that the normalized Ricci flow starting at $\gamma$ converges smoothly
to a metric $g_*$ of constant positive Gaussian curvature.  Because the
two-dimensional Ricci flow is conformal, $g_*$ lies in the conformal
class of $\gamma$.  The normalized flow is
natural under pullback and has a unique solution; since every element of
$\Gamma$ preserves its initial metric, it preserves the entire flow and its
limit $g_*$.  Write $\gamma=e^{2u}g_*$.  The invariance of $\gamma$ and
$g_*$ implies that $u$ is $\Gamma$-invariant.

For $0\le s\le1$, set $h_s=e^{2su}g_*$.  The conformal curvature
formula and the identity
$e^{2u}\sec_\gamma=\sec_{g_*}-\Delta_{g_*}u$ give
\begin{equation}
\sec_{h_s}=e^{-2su}\bigl((1-s)\sec_{g_*}
                         +s e^{2u}\sec_\gamma\bigr)>0.           \label{eq:unified-base-3-71}
\end{equation}
Thus $h_s$ joins $g_*$ to $\gamma$ through metrics of positive Gaussian
curvature.  Reversing the parameter and composing it with a smooth
nondecreasing function $[0,1]\to[0,1]$ which is constant near both
endpoints gives the required path from $\gamma$ to $g_*$.

Now suppose that $m=3$.  Positive sectional curvature implies positive
Ricci curvature, so the main theorem of Hamilton's work
\cite{Hamilton1982} applies to the volume-normalized Ricci flow starting
at $\gamma$.  It gives a solution $g(t)$ for all $t\ge0$, preserves
positive Ricci curvature, and gives smooth convergence to a metric
$g_\infty$ of constant positive sectional curvature.  As in the
two-dimensional case, naturality and uniqueness imply that every
$g(t)$ and $g_\infty$ are $\Gamma$-invariant.

Choose $T$ so large that $g(T)$ lies in a $\Gamma$-invariant
$C^2$-neighborhood of $g_\infty$ on which Ricci curvature is positive.
The affine segment
\[
 (1-s)g(T)+s g_\infty,\qquad 0\le s\le1,
\]
then consists of $\Gamma$-invariant Ricci-positive metrics.  Concatenate the
flow segment $(g(t))_{0\le t\le T}$ with this affine segment, and
reparameterize both pieces so that they are constant near their
endpoints.  This gives the asserted smooth path.
\end{proof}

We now assemble the preceding ingredients.  The angular region supplies
a ball with two disjoint balls removed, the involution exchanges the two
incoming boundary components, the necks give them the prescribed round
geometry, and the path cylinder rounds the outgoing boundary.

\begin{theorem}
\label{thm:unified-directed-base}
Let $m\in\{2,3\}$, and let
$(P_m;\partial_oP_m,I_+,I_-)$ be the fixed topological model defined in
Section~\ref{sec:double-cable}.  For every $R_*>0$ and
$\varepsilon_*>0$, $P_m$ admits a metric $h_m$ with
$\Ric_{h_m}>0$ and an isometric
involution $\sigma_m:P_m\to P_m$ such that:
\begin{enumerate}[label=\textup{(\roman*)}]
\item $\sigma_m(I_+)=I_-$ and $\sigma_m(I_-)=I_+$;
\item $(I_+,h_m|_{I_+})$ and $(I_-,h_m|_{I_-})$ are round
$m$-spheres of radius $R_*$; if $\II_{I_\pm}$ denotes their second
fundamental form with respect to the outward unit normal of $P_m$, then
\begin{equation}
\II_{I_\pm}>-\varepsilon_*h_m|_{I_\pm};     \label{eq:unified-base-3-73}
\end{equation}
\item $(\partial_oP_m,h_m|_{\partial_oP_m})$ is round, and its second
fundamental form with respect to the outward unit normal is positive
definite;
\item after identifying the round outer boundary with a standard round
$\mathbb S^m$,
\begin{equation}
\sigma_m|_{\partial_oP_m}=\exp(\pi K_m).        \label{eq:unified-base-3-74}
\end{equation}
\end{enumerate}
Here the maps $\exp(\pi K_m)$ are defined in
\eqref{eq:unified-base-3-1}; they are the half-turn of $\mathbb S^2$
about an axis and the antipodal map of $\mathbb S^3$, respectively.
\end{theorem}

\begin{proof}
Choose $0<r<r_*$, where $r_*$ is the constant in
Proposition~\ref{prop:unified-neck}, and then choose
\begin{equation}
0<\rho<\min\{1,r/\pi,R_*\varepsilon_*\}.            \label{eq:unified-base-3-75}
\end{equation}
In particular, $\pi\rho<r$.  Choose $w>0$ sufficiently small that
$w^{m-1}<\rho^m$, and apply Proposition~\ref{prop:unified-neck}.
In the local ambient manifold of that proposition, write
\[
 y=o_0,\qquad y_+=o_{2r},\qquad y_-=o_{-2r}.
\]
Thus $y$ and $y_\pm$ are the points obtained by collapsing, respectively,
$\{0\}\times\{0\}\times\mathbb S^k$ and
$\{0\}\times\{\pm2r\}\times\mathbb S^k$.  Define
\[
 \Omega_r:=B_{4r}^{\angle}(y)\setminus
   \left(\operatorname{int}B_r^{\angle}(y_+)
          \sqcup\operatorname{int}B_r^{\angle}(y_-)\right).
\]
Its three boundary components are
\[
 \Sigma_o=S_{4r}^{\angle}(y),\qquad
 \Sigma_+=S_r^{\angle}(y_+),\qquad
 \Sigma_-=S_r^{\angle}(y_-).
\]

We first verify the topology of $\Omega_r$.  For the unscaled warped
metric $ \dd t^2+\cos^2t\,\dd z^2+\mathcal R(t)^2g_{\mathbb S^k}$ in \eqref{eq:unified-base-3-11},
projection onto the $(t,z)$-factor with metric
$\dd t^2+\cos^2t\,\dd z^2$ is length nonincreasing, while every radial
geodesic in that factor has a horizontal lift.  Hence its distance from
$y=o_0$ is the function $\varrho_0$ in the definition of an angular
ball, where
\begin{equation}
\cos\varrho_0=\cos t\cos z.              \label{eq:unified-base-3-76}
\end{equation}
Replacing $z$ by $z\mp2r$ gives the corresponding distance formula
from $y_\pm$.  The factor $\cot^2r$ in
\eqref{eq:unified-base-3-11} multiplies all these distances by
$\cot r$.  Thus the angular balls of angular radii $r$ and $4r$ are
normal balls for $g_{\mathrm{amb},m}$ of Riemannian radii
$r\cot r$ and $4r\cot r$, respectively.  Since the center distances
are $2r\cot r$, $2r\cot r$, and $4r\cot r$, the two closed inner
balls are disjoint, and each lies in the interior of the outer ball.
Consequently, the boundary-decomposed manifold
$(\Omega_r;\Sigma_o,\Sigma_+,\Sigma_-)$ is diffeomorphic to the fixed
model $(P_m;\partial_oP_m,I_+,I_-)$ from
Section~\ref{sec:double-cable}.

We next record the geometry of the outer boundary before any gluing.
With respect to the outward unit normal of $B_{4r}^{\angle}(y)$, the
meridional principal curvature of $\Sigma_o$ is
\begin{equation}
\frac{\cot(4r)}{\cot r}>0,                \label{eq:unified-base-3-77}
\end{equation}
and every vertical principal curvature is
\begin{equation}
\frac1{\cot r}\frac{\mathcal R'(t)}{\mathcal R(t)}
       \frac{\sin t\cos z}{\sin(4r)}>0.                           \label{eq:unified-base-3-78}
\end{equation}
At the two points where the $\mathbb S^k$-factor collapses, the
expression in \eqref{eq:unified-base-3-78} has limit
\eqref{eq:unified-base-3-77}.  Thus $\Sigma_o$ is smooth and strictly
convex.

The ambient metric is invariant under the involution
\begin{equation}
\iota_2(t,z,(p_1,p_2))=(t,-z,(p_1,-p_2)),\qquad
 \iota_3(t,z,p)=(t,-z,-p),                                      \label{eq:unified-base-3-79}
\end{equation}
where $(p_1,p_2)\in\mathbb S^1$ in the first formula and
$p\in\mathbb S^2$ in the second.  This isometry fixes $y$, preserves
$\Sigma_o$, and exchanges $\Sigma_+$ with $\Sigma_-$.

Take two isometric copies of the neck
$([0,\ell]\times\mathbb S^m,g_{\mathrm{neck},m})$ from
Proposition~\ref{prop:unified-neck}.  Identify their $s=\ell$
boundaries with $\Sigma_+$ and $\Sigma_-$ using the latitude coordinate
in \eqref{eq:unified-base-3-12}, and choose the second identification to
be the image of the first under $\iota_m$.  Along each $\Sigma_\pm$, the
outward second fundamental form of $\Omega_r$ is at least
$-g_{\mathrm{ang},m}$: its outward normal is the negative of the normal
pointing out of the removed angular ball.  At the $s=\ell$ end of each
neck, the outward second fundamental form is strictly greater than
$g_{\mathrm{ang},m}$.  The sum of the two forms is positive definite.

Apply Lemma~\ref{lem:relative-gluing-unified} simultaneously at the two
boundary components.  Attaching the two neck cylinders does not change
the diffeomorphism type, so transport the resulting
Ricci-positive metric to the fixed manifold $P_m$.  We denote this
metric by $h_m$, identify $\Sigma_o$ with $\partial_oP_m$, and identify
the two exposed $s=0$ ends with $I_+$ and $I_-$.  The metric remains
unchanged near all three exposed boundary components.  Equivariance in
the gluing lemma transports $\iota_m$ to an isometric involution
$\sigma_m$ of $P_m$ which exchanges the two necks and hence exchanges
$I_+$ and $I_-$.

At each exposed inner end the radius is $\rho/\lambda$ and the outward
principal curvatures are $-\lambda$.  Multiply the entire metric by
\begin{equation}
c^2,\qquad
                         c=\frac{R_*\lambda}{\rho}.                \label{eq:unified-base-3-80}
\end{equation}
Under this homothety, lengths are multiplied by $c$ and principal
curvatures by $c^{-1}$.  Hence both inner radii become $R_*$, while
every inner principal curvature becomes
\[
 -\frac{\lambda}{c}=-\frac{\rho}{R_*}>-\varepsilon_*
\]
by \eqref{eq:unified-base-3-75}.  This proves
\eqref{eq:unified-base-3-73}.  The outer boundary remains strictly
convex.

It remains to make the outer metric exactly round.  The sectional
curvatures of the ambient metric near $\Sigma_o$ and the principal
curvatures of $\Sigma_o$ are positive.  The Gauss equation shows that
the metric induced on $\partial_oP_m=\Sigma_o$ has positive sectional
curvature.  Apply Lemma~\ref{lem:unified-round-path} to this metric and
the order-two group generated by $\sigma_m|_{\partial_oP_m}$.  This
gives a $\sigma_m$-invariant Ricci-positive path from the induced metric
to a round metric.

Apply Theorem~\ref{thm:path-cylinder-unified} to this path, choosing
$\delta_{\mathrm{cyl}}>0$ smaller than the least principal curvature of
$\partial_oP_m$.  The second fundamental form at the initial boundary
of the resulting cylinder is greater than
$-\delta_{\mathrm{cyl}}g_0$.  Its sum with the outward second fundamental form of
$\partial_oP_m$ is therefore positive definite.  Lemma
\ref{lem:relative-gluing-unified} attaches this cylinder to
$\partial_oP_m$.  Both constructions are equivariant, so $\sigma_m$
extends isometrically across the cylinder.  The gluing leaves the metric
unchanged near $I_+\sqcup I_-$, and its new outer boundary is round and
strictly convex.  Attaching this cylinder does not change the diffeomorphism type of $P_m$.

Finally, we identify the involution on the round outer boundary.  At
the fixed point $y=o_0$, the tangent space splits naturally as
\[
                         T_y\cong\mathbb R_z\oplus\mathbb R^m,
\]
and
\begin{equation}
d\iota_2|_y=\operatorname{diag}(-1,1,-1),\qquad
 d\iota_3|_y=-I_4.                                                \label{eq:unified-base-3-81}
\end{equation}
The first map is conjugate to $(-I_2)\oplus(1)$, and the second is
$-I_4$; by \eqref{eq:unified-base-3-1}, these are the linear models of
$\exp(\pi K_2)$ and $\exp(\pi K_3)$.  Since $\Sigma_o$ is the
$g_{\mathrm{amb},m}$-distance sphere of radius $4r\cot r$ centered at
$y$, naturality of the exponential map gives
\[
 \iota_m(\exp_y v)=\exp_y(d\iota_m|_y v).
\]
Thus the involution on the original outer boundary is conjugate to
$\exp(\pi K_m)$.

The outer cylinder does not change the underlying action on
$\mathbb S^m$, while its terminal metric is round.  The terminal action
is therefore an orthogonal involution with the same fixed-point set:
two points for $m=2$ and the empty set for $m=3$.  Its $+1$-eigenspace
has dimension one in the first case and dimension zero in the second.
After choosing an orthogonal identification of the terminal round sphere
with the unit round sphere, the
action is $(-I_2)\oplus(1)$ for $m=2$ and $-I_4$ for $m=3$.  This proves
\eqref{eq:unified-base-3-74} and completes the construction.
\end{proof}

\begin{remark}
\label{rem:unified-differences}
In this construction, the dependence on
the dimension is confined to the following four points, all determined
by the multiplicity $k=m-1$.
\begin{enumerate}[label=\textup{(\arabic*)}]
\item A two-plane tangent to the $\mathbb S^k$-factor exists only when
$m=3$.  Its sectional curvature in the ambient metric, the angular
sphere, the slice metric, and the neck metric is given, respectively,
by \eqref{eq:unified-base-3-18}, \eqref{eq:unified-base-3-21},
\eqref{eq:unified-base-3-30}, and \eqref{eq:unified-base-3-42}.
\item The waist hypothesis is $w<\rho^2$ for $m=2$ and
$w^2<\rho^3$ for $m=3$.  In both cases it says exactly
$\alpha_0<m$.
\item The mixed Ricci component carries the factor $k$, while the
leading coefficient in \eqref{eq:unified-base-3-61} is
$m-\alpha u\eta/f_u$.  The inequality $\alpha<m$ makes this coefficient
positive in both dimensions and yields the determinant estimate
\eqref{eq:unified-base-3-63}.
\item The involution on the collapsing factor is a reflection of
$\mathbb S^1$ for $m=2$ and the antipodal map of $\mathbb S^2$ for
$m=3$.  The induced maps on the outer sphere are, respectively,
$\exp(\pi K_2)$ and $\exp(\pi K_3)$.  Lemma
\ref{lem:unified-round-path} uses Hamilton's surface Ricci-flow theorem
\cite{Hamilton1988Surfaces} in dimension two and his
three-dimensional convergence theorem \cite{Hamilton1982} in dimension
three.
\end{enumerate}
No other step of the local ambient construction, neck interpolation,
smoothing argument, rescaling, or directed-base argument depends on the
dimension.
\end{remark}



\section{The Ricci-positive screw quotient}
\label{sec:unified-cable}

This section constructs the screw quotient $W_m$ defined in (\ref{def:screw-quotient}) from the directed base of
Theorem~\ref{thm:unified-directed-base}.  A circle warp first supplies
the required boundary second fundamental forms.  Taking the quotient
then preserves the outgoing boundary and combines the two incoming
boundary components into one.  We
then identify both quotient boundaries with
$\mathbb S^1\times\mathbb S^m$ and compute their metrics and circle
classes explicitly.

Fix \(m\in\{2,3\}\), \(R_*>0\), and \(\varepsilon_*>0\).  Let
\((P_m,h_m,\sigma_m)\) be supplied by
Theorem~\ref{thm:unified-directed-base}.  Thus
\[
             \partial P_m=\partial_oP_m\sqcup I_+\sqcup I_-,
\]
the involution $\sigma_m$ exchanges $I_+$ and $I_-$, and the two incoming
boundary components are round $m$-spheres of radius $R_*$ satisfying
\[
 \II_{I_\pm}>-\varepsilon_*h_m|_{I_\pm}
\]
with respect to the outward unit normal.  Denote the radius of the
strictly convex round boundary $\partial_oP_m$ by $R_o$.  Under the
round identification in Theorem~\ref{thm:unified-directed-base},
\[
 \sigma_m|_{\partial_oP_m}=\exp(\pi K_m).
\]

\subsection{The invariant circle warp}

The boundary values of the following warping function determine the
circle lengths in the quotient, while its outward logarithmic
derivatives determine the circle-direction components of the second
fundamental forms.

\begin{lemma}\label{lem:un-warp}
Let \((P_m,h_m,\sigma_m)\) be furnished by
Theorem~\ref{thm:unified-directed-base}.  Given \(F>0\) and
\(\delta_*>0\), there is a positive \(\sigma_m\)-invariant function \(f\)
and a constant \(\alpha>0\)
such that
\[
                         G=h_m+f^2\dd t^2
\]
has \(\Ric_G>0\), and
\begin{align}
 f|_{\partial_oP_m}&=F(1+\alpha),&
 \partial_\nu\log f|_{\partial_oP_m}&>0,
 \label{eq:un-warp-outer}\\
 f|_{I_\pm}&=F,&
 \partial_\nu\log f|_{I_\pm}&>-\delta_*.
 \label{eq:un-warp-inner}
\end{align}
Moreover, if \(X,Y\) are tangent to a boundary component of \(P_m\),
then the second fundamental form of the
corresponding boundary of \((\mathbb R\times P_m,G)\), with respect to
the lift of the outward unit normal \(\nu\), satisfies
\begin{equation}\label{eq:un-warp-second-fundamental-form}
 \II_G(X,Y)=\II_{h_m}(X,Y),\qquad
 \II_G(f^{-1}\partial_t,f^{-1}\partial_t)=\partial_\nu\log f,
 \qquad \II_G(X,f^{-1}\partial_t)=0.
\end{equation}
\end{lemma}

\begin{proof}
Let \(\psi\) be the solution of
\[
 \Delta_{h_m}\psi=0,\quad \psi|_{\partial_oP_m}=1,
 \quad \psi|_{I_+\cup I_-}=0,
\]
and let \(\chi\) be the solution of
\[
 \Delta_{h_m}\chi=-1,\quad \chi|_{\partial P_m}=0.
\]
Since \(\sigma_m\) is an isometry, preserves \(\partial_oP_m\), and
interchanges \(I_+\) and \(I_-\), the functions
\(\psi\circ\sigma_m\) and \(\chi\circ\sigma_m\) solve the same Dirichlet
problems as \(\psi\) and \(\chi\), respectively.  Uniqueness therefore gives
\[
                       \psi\circ\sigma_m=\psi,
                 \qquad \chi\circ\sigma_m=\chi.
\]
The strong maximum principle gives \(0<\psi<1\) and \(\chi>0\) in
\(P_m^\circ\).  Applying the Hopf boundary lemma at the maximum and
minimum boundary components of \(\psi\), and at the boundary minimum of
\(\chi\), gives
\[
 \psi_\nu>0\ \hbox{on }\partial_oP_m,\qquad
 \psi_\nu<0\ \hbox{on }I_\pm,\qquad
 \chi_\nu<0\ \hbox{on }\partial P_m.
\]

For \(\alpha>0\), define
\[
                   f=F(1+\alpha\psi+\alpha^2\chi).
\]
This function is positive and \(\sigma_m\)-invariant.  Since \(\chi=0\) on
\(\partial P_m\), its boundary values are
\[
 f|_{\partial_oP_m}=F(1+\alpha),
 \qquad f|_{I_\pm}=F.
\]
On the outer boundary,
\[
 \partial_\nu\log f
 =\frac{\alpha(\psi_\nu+\alpha\chi_\nu)}{1+\alpha}.
\]
Because \(\psi_\nu\) has a positive minimum and \(\chi_\nu\) is bounded on
the compact boundary \(\partial_oP_m\), this expression is positive for
all sufficiently small \(\alpha>0\).
On either inner boundary,
\[
 \partial_\nu\log f
 =\alpha\psi_\nu+\alpha^2\chi_\nu.
\]
The last expression converges uniformly to zero as \(\alpha\to0\).
Thus, after decreasing \(\alpha\), it is greater than \(-\delta_*\) on
both \(I_+\) and \(I_-\).

It remains to verify the Ricci inequality.  For vectors \(X,Y\) tangent
to \(P_m\), the Ricci tensor of the one-dimensional warped product is
\begin{align*}
 \Ric_G(X,Y)
 &=\Ric_{h_m}(X,Y)-f^{-1}\Hess_{h_m}f(X,Y),\\
 \Ric_G(X,f^{-1}\partial_t)&=0,\\
 \Ric_G(f^{-1}\partial_t,f^{-1}\partial_t)
 &=-f^{-1}\Delta_{h_m}f.
\end{align*}
Here
\[
 \Hess_{h_m}f=F\alpha
       \bigl(\Hess_{h_m}\psi+\alpha\Hess_{h_m}\chi\bigr),
 \qquad
 -\Delta_{h_m}f=F\alpha^2>0.
\]
Because \(P_m\) is compact and \(\Ric_{h_m}>0\), the first displayed
Ricci expression is positive definite for all sufficiently small
\(\alpha>0\).  The vertical expression is strictly positive, and the
mixed expression vanishes.  Hence \(\Ric_G>0\).

Finally, the standard Levi--Civita formulas for a one-dimensional
warped product give
\[
 \nabla_X\nu=\nabla^{h_m}_X\nu,
 \qquad
 \nabla_{f^{-1}\partial_t}\nu
   =(\partial_\nu\log f)f^{-1}\partial_t.
\]
Their tangential components are exactly the three identities in
\eqref{eq:un-warp-second-fundamental-form}.
\end{proof}

\subsection{The screw quotient and its product parametrization of the boundary}

Fix $T>0$, and let $\mathcal T$ and $W_m$ be the isometry and screw
quotient defined in (\ref{def:screw-quotient}).  Equip $W_m$ with the metric
induced by $G$.  The transformation
$\mathcal T$ is an isometry because $h_m$ and $f$ are
$\sigma_m$-invariant.  On $\mathbb R\times\partial_oP_m$ its monodromy
is $\exp(\pi K_m)$.  It exchanges
$\mathbb R\times I_+$ and $\mathbb R\times I_-$, and its first return to
$\mathbb R\times I_+$ is
\[
 \mathcal T^2(t,x)=(t+2T,x).
\]
Consequently the quotient has one outer and one inner boundary
component; the outer circle has period $T$, whereas the inner circle
has period $2T$.  The following product metrics make this distinction
precise.

For \(\ell,\rho>0\) and \(\omega\in\mathbb R\), define a metric
\(g_{\ell,\rho,\omega}\) on
\(\mathbb S^1\times\mathbb S^m\), where \(\theta\) has period \(2\pi\),
by prescribing the squared norm of each tangent vector:
\begin{equation}\label{eq:un-boundary-metric-definition}
 g_{\ell,\rho,\omega}(b\partial_\theta+Y,b\partial_\theta+Y)
 :=\ell^2b^2+\rho^2
   \lvert Y+b\omega K_m(y)\rvert_{g_{\mathbb S^m}}^2,
 \qquad Y\in T_y\mathbb S^m.
\end{equation}
This definition is global even when \(K_m\) has zeros.

\begin{lemma}
\label{lem:un-boundary-components}
The boundary of \(W_m\) is the disjoint union
\[
                         \partial W_m=\partial_oW_m\sqcup\partial_iW_m,
\]
where
\begin{align*}
 \partial_oW_m
 &:=(\mathbb R\times\partial_oP_m)/
   \bigl((t,x)\sim(t+T,\sigma_mx)\bigr),\\
 \partial_iW_m
 &:=(\mathbb R\times(I_+\sqcup I_-))/
   \bigl((t,x)\sim(t+T,\sigma_mx)\bigr).
\end{align*}
Use the round identifications
\(\partial_oP_m\cong\mathbb S^m_{R_o}\) and
\(I_+\cong\mathbb S^m_{R_*}\) fixed above.  Then the maps
\begin{align}
 \Phi_o:\mathbb S^1\times\mathbb S^m&\longrightarrow\partial_oW_m,
 &\Phi_o(\theta,y)
   &:=\left[\frac{T\theta}{2\pi},
       \exp\left(\frac\theta2K_m\right)y\right],
 \label{eq:un-outer-product-map}\\
 \Phi_i:\mathbb S^1\times\mathbb S^m&\longrightarrow\partial_iW_m,
 &\Phi_i(\theta,y)&:=\left[\frac{T\theta}{\pi},y\right]
 \label{eq:un-inner-product-map}
\end{align}
are diffeomorphisms.  They are the boundary product markings of $W_m$.
Moreover,
\begin{align}
 \Phi_o^*(G|_{\partial_oW_m})&=g_{\ell_o,R_o,1/2},&
 \ell_o&=\frac{TF(1+\alpha)}{2\pi},
 \label{eq:un-outer-boundary-metric}\\
 \Phi_i^*(G|_{\partial_iW_m})&=g_{\ell_i,R_*,0},&
 \ell_i&=\frac{2TF}{2\pi}.
 \label{eq:un-inner-boundary-metric}
\end{align}
In particular,
\begin{equation}\label{eq:un-axial-ratio}
                         \ell_o=\frac{1+\alpha}{2}\ell_i.
\end{equation}
Finally, \(\pi_1(W_m)\cong\mathbb Z\).  For any \(y\in\mathbb S^m\),
the loop \(\theta\mapsto\Phi_o(\theta,y)\) represents a generator, while
\(\theta\mapsto\Phi_i(\theta,y)\) represents the square of that
generator.
\end{lemma}

\begin{proof}
The action of \(\mathbb Z\) on \(\mathbb R\times P_m\) generated by
\((t,x)\mapsto(t+T,\sigma_m x)\) is free and properly discontinuous.
A fundamental domain is \([0,T]\times P_m\), with
\((T,x)\) identified with \((0,\sigma_m x)\).  Thus \(W_m\) is compact,
and the Ricci-positive metric \(G\) descends to it.

The action preserves \(\mathbb R\times\partial_oP_m\) and preserves
\(\mathbb R\times(I_+\sqcup I_-)\), while interchanging the two summands
in the latter set.  Taking their quotients gives precisely the two
boundary components stated in the lemma.

We first verify \eqref{eq:un-outer-product-map}.  Replacing \(\theta\) by
\(\theta+2\pi\)
changes the displayed representative to
\[
 \left[\frac{T\theta}{2\pi}+T,
   \sigma_m\exp\left(\frac\theta2K_m\right)y\right],
\]
which represents the same point in the quotient.  Conversely, every
class in \(\partial_oW_m\) has a representative with \(0\le t<T\);
setting \(\theta=2\pi t/T\) and applying
\(\exp(-\theta K_m/2)\) to its sphere coordinate gives the inverse of
\(\Phi_o\).  Hence \(\Phi_o\) is a diffeomorphism.

For tangent vectors \(b\partial_\theta+Y\) and
\(c\partial_\theta+Z\), the differential of
\(\Phi_o\), after applying the inverse of the
sphere isometry \(\exp(\theta K_m/2)\), gives
\[
 \frac{Tb}{2\pi}\partial_t+Y+\frac b2K_m,
 \qquad
 \frac{Tc}{2\pi}\partial_t+Z+\frac c2K_m.
\]
On \(\partial_oP_m\), Lemma~\ref{lem:un-warp} gives
\(f=F(1+\alpha)\), while the sphere metric is
\(R_o^2g_{\mathbb S^m}\).  The pullback of \(G\) is therefore
\[
 \left(\frac{TF(1+\alpha)}{2\pi}\right)^2bc
 +R_o^2g_{\mathbb S^m}
   \left(Y+\frac b2K_m,Z+\frac c2K_m\right),
\]
which proves \eqref{eq:un-outer-boundary-metric}.

For the inner boundary, fix the round identification
\(I_+\cong\mathbb S^m_{R_*}\), and use \(\sigma_m:I_+\to I_-\) for the
corresponding identification of \(I_-\).  Every equivalence class in
the quotient of \(\mathbb R\times(I_+\sqcup I_-)\) has a representative
in \(\mathbb R\times I_+\).  Two representatives in
\(\mathbb R\times I_+\) determine the same class precisely when their
first coordinates differ by an even multiple of \(T\), because one
application of the defining identification moves \(I_+\) to \(I_-\),
whereas two applications return to \(I_+\).  Thus \(\Phi_i\) in
\eqref{eq:un-inner-product-map} is a well-defined
diffeomorphism.  Since \(f=F\) and the sphere metric is
\(R_*^2g_{\mathbb S^m}\) on \(I_+\), its pullback metric is
\[
 \left(\frac{TF}{\pi}\right)^2\dd\theta^2
   +R_*^2g_{\mathbb S^m}=g_{\ell_i,R_*,0},
 \qquad \ell_i=\frac{2TF}{2\pi}.
\]
This proves \eqref{eq:un-inner-boundary-metric}, and
\eqref{eq:un-axial-ratio} follows immediately.

It remains to identify the two circle classes.  By the definition of
the fixed model $P_m$ in Section~\ref{sec:double-cable}, attaching an
\((m+1)\)-disk to each of \(I_+\) and \(I_-\) produces $D^{m+1}$.
Since \(m\ge2\), the attaching spheres \(I_+\) and
\(I_-\) are simply connected.  Applying the van Kampen theorem
at the two attachments therefore shows that
\(\pi_1(P_m)=0\).

The quotient map makes \(W_m\) a fiber bundle over the circle, with
projection
\[
 W_m\longrightarrow\mathbb R/T\mathbb Z,
 \qquad [t,x]\longmapsto[t],
\]
and fiber \(P_m\).  Its homotopy exact sequence gives an isomorphism
\(\pi_1(W_m)\cong\pi_1(\mathbb R/T\mathbb Z)\cong\mathbb Z\).  Under
\(\Phi_o\), the loop obtained by holding \(y\) fixed projects once around
\(\mathbb R/T\mathbb Z\), and hence represents a generator.  Under
\(\Phi_i\), the corresponding loop changes \(t\) by \(2T\), so it projects
twice around the base and represents the square of the first loop.
\end{proof}



\section{The holonomy connector}
\label{sec:holonomy-connector}

Fix \(m\in\{2,3\}\).  We now construct the holonomy connector $\mathcal C$ defined in (\ref{def:holonomy-connector}) and prove that its metric has positive Ricci
curvature.  It joins the twisted outer boundary of the screw quotient to
the untwisted product boundary required by the next block.  By
Lemma~\ref{lem:un-boundary-components}, the outer boundary of
\(W_m\) has metric \(g_{\ell_o,R_o,1/2}\).  The product identification
\(\Phi_o\) reduces the geometric problem to a path in the family
\(g_{\ell,\rho,\omega}\) on
\(\mathbb S^1\times\mathbb S^m\): the functions \(\ell\) and \(\rho\)
specify the circle and sphere scales, while \(\omega\) specifies the
holonomy.

We require the twist to change from \(\omega(0)=1/2\) to
\(\omega(s_*)=0\), while the sphere radius returns to its initial value
at \(s=s_*\).  We first compute the Ricci tensor and the two boundary
second fundamental forms exactly.  We then choose \(\omega\) by a
weighted least-energy principle, which eliminates the axial--sphere
Ricci components, and choose \(\ell\) and \(\rho\) so that the positive
axial and spherical terms dominate the twisting errors.  Finally, we
verify the second fundamental form inequality needed to attach the
left boundary of $\mathcal C$ to \(\partial_oW_m\).

Let \(s_*>0\), let \(\ell,\rho:[0,s_*]\to(0,\infty)\), and let
\(\omega:[0,s_*]\to\mathbb R\) be smooth.  Set
\[
 \mathcal C:=[0,s_*]\times\mathbb S^1\times\mathbb S^m
\]
and define on $\mathcal C$
\begin{equation}\label{eq:un-connector}
 g_{\mathcal C}=\dd s^2+g_{\ell(s),\rho(s),\omega(s)}.
\end{equation}
When
\[
 \omega(0)=\frac12,\qquad \omega(s_*)=0,
 \qquad \rho(s_*)=\rho(0),
\]
we call $(\mathcal C,g_{\mathcal C})$ a \emph{holonomy connector}: its
left boundary has half-turn holonomy, its right boundary is an untwisted
product, and the sphere radii of the two boundary metrics agree.
Thus, for \(Y\in T_y\mathbb S^m\),
\[
 g_{\mathcal C}(b\partial_\theta+Y,b\partial_\theta+Y)
 =\ell^2b^2+\rho^2
   \lvert Y+b\omega K_m(y)\rvert_{g_{\mathbb S^m}}^2.
\]
Define the logarithmic derivatives and the normalized twist rate by
\begin{equation}\label{eq:un-log-derivatives}
 \mathfrak p=\frac{\ell'}{\ell},
 \qquad \mathfrak q=\frac{\rho'}{\rho},
 \qquad \vartheta=\frac{\rho\omega'}{\ell}.
\end{equation}
Let \(U_1,\ldots,U_m\) be a local
\(g_{\mathbb S^m}\)-orthonormal frame.  The coordinate column of the
Killing field \(K_m\) in this frame is
\[
 K_{m,i}=g_{\mathbb S^m}(K_m,U_i),
 \qquad
 \mathbf K_m=(K_{m,1},\ldots,K_{m,m})^{\mathsf T}.
\]
Then
\(\lvert\mathbf K_m\rvert^2
=\lvert K_m\rvert_{g_{\mathbb S^m}}^2\le1\), and the following frame is
\(g_{\mathcal C}\)-orthonormal:
\begin{equation}\label{eq:un-adapted-frame}
 E_0=\partial_s,
 \qquad
 E_1=\ell^{-1}(\partial_\theta-\omega K_m),
 \qquad
 E_{i+1}=\rho^{-1}U_i\quad(1\le i\le m).
\end{equation}

\begin{lemma}
\label{lem:un-ricci}
Write
\(\Ric_{\alpha\beta}:=\Ric_{g_{\mathcal C}}(E_\alpha,E_\beta)\).  In the frame
\eqref{eq:un-adapted-frame}, the potentially nonzero components not
determined by symmetry are
\begin{align}
 \Ric_{00}
  &=-(\mathfrak p'+\mathfrak p^2)-m(\mathfrak q'+\mathfrak q^2)
    -\frac{\vartheta^2}{2}\lvert\mathbf K_m\rvert^2,
  \label{eq:un-ric00}\\
 \Ric_{11}
  &=-\bigl[\mathfrak p'+\mathfrak p(\mathfrak p+m\mathfrak q)\bigr]
    -\frac{\vartheta^2}{2}\lvert\mathbf K_m\rvert^2,
  \label{eq:un-ric11}\\
 \Ric_{i+1,j+1}
  &=\left\{\frac{m-1}{\rho^2}
    -\bigl[\mathfrak q'+\mathfrak q(\mathfrak p+m\mathfrak q)\bigr]\right\}\delta_{ij}
    +\frac{\vartheta^2}{2}K_{m,i}K_{m,j},
  \label{eq:un-ricij}\\
 \Ric_{1,i+1}
  &=-\frac12\bigl[\vartheta'+(m+1)\mathfrak q\vartheta\bigr]K_{m,i}
   =-\frac{K_{m,i}}{2\rho^{m+1}}
       \left(\frac{\rho^{m+2}\omega'}{\ell}\right)'.
  \label{eq:un-ric1i}
\end{align}
Here \(1\le i,j\le m\), and
\[
 \Ric_{0\alpha}=0\qquad(1\le\alpha\le m+1).
\]
With respect to the unit normal \(\partial_s\),
the second fundamental form of the hypersurface
\(\{s\}\times\mathbb S^1\times\mathbb S^m\) satisfies
\begin{align}
 \II_s(E_1,E_1)&=\mathfrak p,\nonumber\\
 \II_s(E_1,E_{i+1})&=\frac{\vartheta}{2}K_{m,i},
 \label{eq:un-shape}\\
 \II_s(E_{i+1},E_{j+1})&=\mathfrak q\delta_{ij}.\nonumber
\end{align}
These identities hold for every local orthonormal frame on
\(\mathbb S^m\), including at the zeros of \(K_m\), which occur when
\(m=2\).
\end{lemma}

\begin{proof}
We fix \(s\).  Since the flow of \(K_m\) consists of isometries of the
unit round sphere, on every coordinate neighborhood the local
diffeomorphism
\[
 (\theta,y)\longmapsto
 \bigl(\theta,\exp(\omega(s)\theta K_m)y\bigr)
\]
pulls \(\ell^2\dd\theta^2+\rho^2g_{\mathbb S^m}\) back to
\(g_{\ell,\rho,\omega}\).  Therefore the intrinsic Ricci tensor of the
\(s\)-slice
satisfies
\begin{equation}\label{eq:un-slice-ricci}
 \Ric_{g_{\ell,\rho,\omega}}(E_1,E_1)=0,
 \qquad
 \Ric_{g_{\ell,\rho,\omega}}(E_{i+1},E_{j+1})
 =\frac{m-1}{\rho^2}\delta_{ij},
\end{equation}
and its axial--sphere components vanish.

In the Gaussian coordinate \(s\), the second fundamental form with
normal \(\partial_s\) is one half of the \(s\)-derivative of the slice
metric.  Evaluating this derivative on the frame
\eqref{eq:un-adapted-frame} gives
\begin{align*}
 \frac12(\partial_s g_{\ell,\rho,\omega})(E_1,E_1)
 &=\frac{\ell'}{\ell}=\mathfrak p,\\
 \frac12(\partial_s g_{\ell,\rho,\omega})(E_1,E_{i+1})
 =\frac{\rho\omega'}{2\ell}K_{m,i}
 &=\frac{\vartheta}{2}K_{m,i},\\
 \frac12(\partial_s g_{\ell,\rho,\omega})(E_{i+1},E_{j+1})
 =\frac{\rho'}{\rho}\delta_{ij}
 &=\mathfrak q\delta_{ij},
\end{align*}
which proves \eqref{eq:un-shape}.  The Gaussian-coordinate identity
\(\nabla_{\partial_s}X=\partial_sX+\II_s^\sharp X\) for tangential
vector fields will now be applied to the frame
\eqref{eq:un-adapted-frame}.  Extend \(U_i\) independently of \(s\).
Then
\[
 \partial_sE_1=-\mathfrak p E_1-\vartheta\sum_i K_{m,i}E_{i+1},
 \qquad
 \partial_sE_{i+1}=-\mathfrak q E_{i+1},
\]
whereas \eqref{eq:un-shape} gives
\[
 \II_s^\sharp E_1=\mathfrak p E_1+\frac{\vartheta}{2}
   \sum_i K_{m,i}E_{i+1},
 \qquad
 \II_s^\sharp E_{i+1}
 =\frac{\vartheta}{2}K_{m,i}E_1+\mathfrak q E_{i+1}.
\]
Adding the corresponding terms yields
\begin{equation}\label{eq:un-normal-rotation}
 \nabla_{E_0}E_1=-\frac{\vartheta}{2}\sum_i K_{m,i}E_{i+1},
 \qquad
 \nabla_{E_0}E_{i+1}=\frac{\vartheta}{2}K_{m,i}E_1.
\end{equation}
Consequently,
\begin{align*}
 (\nabla_{E_0}\II_s)(E_1,E_1)
 &=\mathfrak p'+\frac{\vartheta^2}{2}\lvert\mathbf K_m\rvert^2,\\
 (\nabla_{E_0}\II_s)(E_{i+1},E_{j+1})
 &=\mathfrak q'\delta_{ij}-\frac{\vartheta^2}{2}K_{m,i}K_{m,j},\\
 (\nabla_{E_0}\II_s)(E_1,E_{i+1})
 &=\frac12\bigl[\vartheta'+(\mathfrak q-\mathfrak p)\vartheta\bigr]K_{m,i}.
\end{align*}
Indeed, the definition of the covariant derivative of a symmetric
two-tensor and \eqref{eq:un-normal-rotation} give
\begin{align*}
 (\nabla_{E_0}\II_s)(E_1,E_1)
 &=\mathfrak p'
   -2\II_s\left(-\frac{\vartheta}{2}
       \sum_r K_{m,r}E_{r+1},E_1\right)
 =\mathfrak p'+\frac{\vartheta^2}{2}\lvert\mathbf K_m\rvert^2,\\
 (\nabla_{E_0}\II_s)(E_{i+1},E_{j+1})
 &=\mathfrak q'\delta_{ij}
   -\II_s\left(\frac{\vartheta}{2}K_{m,i}E_1,E_{j+1}\right)
   -\II_s\left(E_{i+1},\frac{\vartheta}{2}K_{m,j}E_1\right)\\
 &=\mathfrak q'\delta_{ij}-\frac{\vartheta^2}{2}K_{m,i}K_{m,j},\\
 (\nabla_{E_0}\II_s)(E_1,E_{i+1})
 &=\frac{\vartheta'}2 K_{m,i}
   -\II_s\left(-\frac{\vartheta}{2}
       \sum_r K_{m,r}E_{r+1},E_{i+1}\right)
   -\II_s\left(E_1,\frac{\vartheta}{2}K_{m,i}E_1\right)\\
 &=\frac12\bigl[\vartheta'+(\mathfrak q-\mathfrak p)\vartheta\bigr]K_{m,i}.
\end{align*}
Moreover,
\[
 \operatorname{tr}_{g_{\ell,\rho,\omega}}\II_s=\mathfrak p+m\mathfrak q,
 \qquad
 |\II_s|^2=\mathfrak p^2+m\mathfrak q^2
   +\frac{\vartheta^2}{2}\lvert\mathbf K_m\rvert^2.
\]
The coefficient \(\vartheta^2/2\) is the sum of the two equal contributions
from the symmetric axial--sphere entries of \(\II_s\).

For vectors \(X,Y\) tangent to the \(s\)-slice, the
identities in Lemma~\ref{lem:gaussian-ricci-unified} can be written as
\begin{align*}
 \Ric_{g_{\mathcal C}}(E_0,E_0)
 &=-\partial_s
   (\operatorname{tr}_{g_{\ell,\rho,\omega}}\II_s)-|\II_s|^2,\\
 \Ric_{g_{\mathcal C}}(X,Y)
 &=\Ric_{g_{\ell,\rho,\omega}}(X,Y)
  -(\nabla_{E_0}\II_s)(X,Y)
  -(\operatorname{tr}_{g_{\ell,\rho,\omega}}\II_s)\II_s(X,Y).
\end{align*}
For the normal component this gives
\[
 \Ric_{00}
 =-(\mathfrak p'+m\mathfrak q')
  -\left(\mathfrak p^2+m\mathfrak q^2
          +\frac{\vartheta^2}{2}\lvert\mathbf K_m\rvert^2\right),
\]
which is \eqref{eq:un-ric00}.  For the tangential components,
\eqref{eq:un-slice-ricci} and
the three preceding identities for \(\nabla_{E_0}\II_s\)
give
\begin{align*}
 \Ric_{11}
 &=0-\left(\mathfrak p'
       +\frac{\vartheta^2}{2}\lvert\mathbf K_m\rvert^2\right)
   -(\mathfrak p+m\mathfrak q)\mathfrak p,\\
 \Ric_{i+1,j+1}
 &=\frac{m-1}{\rho^2}\delta_{ij}
   -\left(\mathfrak q'\delta_{ij}
          -\frac{\vartheta^2}{2}K_{m,i}K_{m,j}\right)
   -(\mathfrak p+m\mathfrak q)\mathfrak q\delta_{ij},\\
 \Ric_{1,i+1}
 &=0-\frac12\bigl[\vartheta'+(\mathfrak q-\mathfrak p)\vartheta\bigr]K_{m,i}
   -\frac12(\mathfrak p+m\mathfrak q)\vartheta K_{m,i}\\
 &=-\frac12\bigl[\vartheta'+(m+1)\mathfrak q\vartheta\bigr]K_{m,i}.
\end{align*}
These are \eqref{eq:un-ric11}--\eqref{eq:un-ric1i}.

It remains to compute the components containing one \(E_0\).  Under the
local diffeomorphism used above, the slice metric becomes the product
metric on \(\mathbb S^1_\ell\times\mathbb S^m_\rho\), and the second
fundamental form becomes
\[
 \mathfrak p\ell^2\dd\theta^2+\mathfrak q\rho^2g_{\mathbb S^m}
 +\frac{\rho^2\omega'}2\left(
   \dd\theta\otimes g_{\mathbb S^m}(K_m,\mathord\cdot)
   +g_{\mathbb S^m}(K_m,\mathord\cdot)\otimes\dd\theta\right).
\]
The first two summands are parallel on the product.  The divergence of
the last summand is zero: \(\dd\theta\) is parallel, \(K_m\) is
independent of \(\theta\), and
\(\operatorname{div}_{g_{\mathbb S^m}}K_m=0\) because \(K_m\) is a
Killing field.  Since
\(\operatorname{tr}_{g_{\ell,\rho,\omega}}\II_s
=\mathfrak p+m\mathfrak q\) is constant on each slice,
the contracted Codazzi identity
\[
 \Ric_{g_{\mathcal C}}(E_0,\mathord\cdot)
 =\operatorname{div}_{g_{\ell,\rho,\omega}}\II_s
  -\dd(\operatorname{tr}_{g_{\ell,\rho,\omega}}\II_s)
\]
implies \(\Ric_{g_{\mathcal C}}(E_0,E_\alpha)=0\) for every \(\alpha\ge1\).
Finally,
\[
 \frac1{\rho^{m+1}}
 \left(\frac{\rho^{m+2}\omega'}{\ell}\right)'
 =\frac1{\rho^{m+1}}(\rho^{m+1}\vartheta)'
 =\vartheta'+(m+1)\mathfrak q\vartheta,
\]
which proves the second expression in \eqref{eq:un-ric1i}.
\end{proof}

Equation~\eqref{eq:un-ric1i} isolates the only Ricci components coupling
the circle and sphere directions.  They vanish identically precisely
when \(\rho^{m+2}\omega'/\ell\) is constant.  This condition is the
Euler--Lagrange equation of the natural weighted twisting energy below.
The variational characterization determines \(\omega\) uniquely once
\(\ell,\rho\) and the two endpoint values of \(\omega\) have been
prescribed.

\begin{lemma}\label{lem:least-energy-twist}
Let $\ell,\rho\in C^\infty([0,s_*],(0,\infty))$ and let
$\omega_-,\omega_+\in\mathbb R$.  Among all
$\omega\in H^1([0,s_*])$ with $\omega(0)=\omega_-$ and
$\omega(s_*)=\omega_+$, the weighted Dirichlet energy functional
\begin{equation}\label{eq:weighted Dirichlet energy}
\mathcal E(\omega)=\int_0^{s_*}
 \frac{\rho(s)^{m+2}}{\ell(s)}\bigl(\omega'(s)\bigr)^2\,\dd s
\end{equation}
has the unique minimizer
\[
 \omega_*(s)=\omega_-+(\omega_+-\omega_-)
 \frac{\displaystyle\int_0^s
        \frac{\ell(t)}{\rho(t)^{m+2}}\,\dd t}
      {\displaystyle\int_0^{s_*}
        \frac{\ell(t)}{\rho(t)^{m+2}}\,\dd t}.
\]
Equivalently, it is the unique solution of
\[
 \left(\frac{\rho^{m+2}\omega'}{\ell}\right)'=0.
\]
\end{lemma}

\begin{proof}
For every $\eta\in H_0^1([0,s_*])$, the first variation at $\omega$ in
the direction $\eta$ is
\[
 D\mathcal E_\omega[\eta]
 =2\int_0^{s_*}\frac{\rho^{m+2}}\ell\omega'\eta'\,\dd s
 =-2\int_0^{s_*}
   \left(\frac{\rho^{m+2}\omega'}\ell\right)'\eta\,\dd s.
\]
Thus the Euler--Lagrange equation is the displayed differential
equation in the statement.  Integrating it and imposing the two
endpoint values gives the stated formula for $\omega_*$.  For every
$\eta\in H_0^1([0,s_*])$,
\[
 \mathcal E(\omega_*+\eta)
 =\mathcal E(\omega_*)+
  \int_0^{s_*}\frac{\rho^{m+2}}\ell(\eta')^2\,\dd s,
\]
because the cross term vanishes.  Hence $\omega_*$ is the unique
minimizer.  By \eqref{eq:un-ric1i}, the same equation makes every
axial--sphere Ricci component vanish.
\end{proof}

With the mixed Ricci components eliminated, it remains to choose the two
scale functions.  The next proposition chooses \(\ell\) so that the
axial Ricci terms have a uniform positive lower bound, and chooses
\(\rho\) to return to its initial value while having small logarithmic
derivatives of opposite signs at the two boundary components.  The
least-energy condition then determines \(\omega\) uniquely.

\begin{proposition}
\label{prop:un-connector}
Let \(\ell_0,\rho_0>0\).  Choose
\[
 0<\mathfrak p_0<\frac1{\rho_0},
 \qquad
 \Lambda_{\mathrm{con}}\ge
 \max\left\{2,\frac{40\rho_0}{\ell_0}\right\},
\]
and put
\begin{equation}\label{eq:un-connector-parameters}
 \mathfrak p_1=\frac{\mathfrak p_0}{\Lambda_{\mathrm{con}}^2},
 \qquad
 s_*=\frac{\Lambda_{\mathrm{con}}^2-1}{2\mathfrak p_0}.
\end{equation}
Define, for \(0\le s\le s_*\),
\begin{align}
 \ell(s)&=\ell_0\sqrt{1+2\mathfrak p_0s},
 \label{eq:un-ell-profile}\\
 \rho(s)&=\rho_0\exp\left(
 -\frac{\mathfrak p_1}{100s_*}s(s_*-s)\right),
 \label{eq:un-rho-profile}\\
\omega(s)&=\frac12\left(
  1-\frac{\displaystyle\int_0^s
               \frac{\ell(t)}{\rho(t)^{m+2}}\dd t}
             {\displaystyle\int_0^{s_*}
               \frac{\ell(t)}{\rho(t)^{m+2}}\dd t}
 \right).
 \label{eq:un-omega-profile}
\end{align}
Then the metric \eqref{eq:un-connector} has \(\Ric>0\), and its
boundary values are
\[
\begin{array}{c|c|c|c}
 &\ell&\rho&\omega\\ \hline
 s=0&\ell_0&\rho_0&1/2\\
 s=s_*&\Lambda_{\mathrm{con}}\ell_0&\rho_0&0.
\end{array}
\]
With respect to the outward normals of $\mathcal C$ and the frame
\eqref{eq:un-adapted-frame}, the matrices of the boundary second
fundamental forms are
\begin{align}
 \II^{\mathrm{out}}_0&=
 \begin{pmatrix}
  -\mathfrak p_0&-\frac{\vartheta(0)}2\mathbf K_m^{\mathsf T}\\[2pt]
  -\frac{\vartheta(0)}2\mathbf K_m&
     \frac{\mathfrak p_1}{100}\operatorname{Id}_m
 \end{pmatrix},
 \label{eq:un-II-left}\\
 \II^{\mathrm{out}}_{s_*}&=
 \begin{pmatrix}
  \mathfrak p_1&\frac{\vartheta(s_*)}2\mathbf K_m^{\mathsf T}\\[2pt]
  \frac{\vartheta(s_*)}2\mathbf K_m&
     \frac{\mathfrak p_1}{100}\operatorname{Id}_m
 \end{pmatrix}.
 \label{eq:un-II-right}
\end{align}
In particular, the right boundary is strictly convex:
\[
 \II^{\mathrm{out}}_{s_*}>0.
\]
\end{proposition}

\begin{proof}
We first record the contribution of the axial profile.  From
\eqref{eq:un-ell-profile},
\[
 \mathfrak p(s)=\frac{\mathfrak p_0}{1+2\mathfrak p_0s},
 \qquad
 \mathfrak p'=-2\mathfrak p^2,
 \qquad
 \mathfrak p_1\le\mathfrak p(s)\le\mathfrak p_0.
\]
Indeed, \(1+2\mathfrak p_0s_*=\Lambda_{\mathrm{con}}^2\), so
\[
 \ell(s_*)=\Lambda_{\mathrm{con}}\ell_0,
 \qquad
 \mathfrak p(s_*)=
 \frac{\mathfrak p_0}{\Lambda_{\mathrm{con}}^2}=\mathfrak p_1.
\]
In particular,
\[
 -\bigl(\mathfrak p'+\mathfrak p^2\bigr)=\mathfrak p^2
 \ge\mathfrak p_1^2,
\]
which will dominate the error terms below.

We next consider the spherical radius.  Equation
\eqref{eq:un-rho-profile} gives
\[
 \mathfrak q(s)=\frac{\mathfrak p_1}{100}
 \left(\frac{2s}{s_*}-1\right).
\]
Its integral over \([0,s_*]\) is zero.  Thus \(\rho(s_*)=\rho_0\),
while \(s(s_*-s)\le s_*^2/4\) and
\eqref{eq:un-connector-parameters} give
\begin{equation}\label{eq:un-rho-bound}
 \rho(s_*)=\rho_0,
 \qquad
 e^{-(\Lambda_{\mathrm{con}}^2-1)/
       (800\Lambda_{\mathrm{con}}^2)}\rho_0
 \le\rho(s)\le\rho_0.
\end{equation}
Moreover,
\[
 |\mathfrak q|\le\frac{\mathfrak p_1}{100},
 \qquad
 \mathfrak q'=
 \frac{\Lambda_{\mathrm{con}}^2}
 {25(\Lambda_{\mathrm{con}}^2-1)}\mathfrak p_1^2
\le\frac{4}{75}p_1^2
 <\frac{3}{50}p_1^2.
\]
When \(\mathfrak q\ge0\), necessarily
\(s\in[s_*/2,s_*]\), and direct substitution yields
\[
 \frac{\mathfrak p\mathfrak q}{\mathfrak p_1^2/100}
 =\frac{\Lambda_{\mathrm{con}}^2(2s/s_*-1)}
 {1+(\Lambda_{\mathrm{con}}^2-1)s/s_*}\le1.
\]
When \(\mathfrak q\le0\), we instead use
\[
 \mathfrak p+m\mathfrak q
 \ge\left(1-\frac3{100}\right)\mathfrak p_1>0,
 \qquad
 \mathfrak q(\mathfrak p+m\mathfrak q)\le0.
\]

After \(\ell\) and \(\rho\) have been fixed, the twist is determined by
\eqref{eq:un-omega-profile}.  Differentiating that formula gives
\[
 \frac{\rho^{m+2}\omega'}{\ell}
 =-\frac12\left(\int_0^{s_*}
       \frac{\ell(t)}{\rho(t)^{m+2}}\dd t\right)^{-1}.
\]
The right-hand side is independent of \(s\); hence
\begin{equation}\label{eq:un-mixed-cancellation}
 \left(\frac{\rho^{m+2}\omega'}{\ell}\right)'=0.
\end{equation}
It follows from \eqref{eq:un-ric1i} that
\(\Ric_{1,i+1}=0\) for every \(1\le i\le m\).  The normalization in
\eqref{eq:un-omega-profile} also gives
\(\omega(0)=1/2\) and \(\omega(s_*)=0\).

Since \(\rho\le\rho_0\),
\[
 \int_0^{s_*}\frac{\ell(s)}{\rho(s)^{m+2}}\dd s
 \ge\frac1{\rho_0^{m+2}}\int_0^{s_*}\ell(s)\dd s
 =\frac{\ell_0(\Lambda_{\mathrm{con}}^3-1)}
 {3\mathfrak p_0\rho_0^{m+2}}.
\]
By \eqref{eq:un-log-derivatives} and
\eqref{eq:un-omega-profile},
\[
 \vartheta(s)=-\frac1{2\rho(s)^{m+1}}
 \left(\int_0^{s_*}\frac{\ell(t)}{\rho(t)^{m+2}}\dd t\right)^{-1}.
\]
Since \(m+1\le4\), the lower bound for \(\rho\) in
\eqref{eq:un-rho-bound} and the choice
\(\Lambda_{\mathrm{con}}\ge40\rho_0/\ell_0\) imply
\begin{align}
 \frac{|\vartheta|}{\mathfrak p_1}
 &\le\frac{3e^{1/200}}2
       \frac{\rho_0}{\ell_0}
       \frac{\Lambda_{\mathrm{con}}^2}
            {\Lambda_{\mathrm{con}}^3-1}\le\frac{3e^{1/200}}{80}
       \frac{\Lambda_{\mathrm{con}}^3}
            {\Lambda_{\mathrm{con}}^3-1}\le\frac{3e^{1/200}}{70}<\frac{1}{10}..
 \label{eq:un-twist-rate-bound}
\end{align}
Here we used \(\Lambda_{\mathrm{con}}\ge2\), for which
\(\Lambda_{\mathrm{con}}^3/
(\Lambda_{\mathrm{con}}^3-1)\le8/7\).

We now substitute these estimates into the formulas of
Lemma~\ref{lem:un-ricci}.  The remaining terms are measured relative to
the positive scale \(\mathfrak p_1^2\).  Substituting
\(\mathfrak p'=-2\mathfrak p^2\) into
\eqref{eq:un-ric11} gives
\[
 \Ric_{11}=\mathfrak p^2-m\mathfrak p\mathfrak q
            -\frac{\vartheta^2}{2}|\mathbf K_m|^2.
\]
If \(\mathfrak q\le0\), its middle term is nonnegative.  If
\(\mathfrak q\ge0\), we use the estimate
\(\mathfrak p\mathfrak q\le\mathfrak p_1^2/100\) proved above.
Since \(m\le3\) and \(|K_m|\le1\),
\begin{align*}
 \Ric_{11}
 &\ge\left(
  1-\frac{3}{100}-\frac{1}{200}
 \right)p_1^2
 >\frac{19}{20}p_1^2.
\end{align*}
Here and below the displayed numerical constants are uniform for
\(m=2,3\).  Similarly, \eqref{eq:un-ric00} becomes
\begin{align*}
 \Ric_{00}
 &=\mathfrak p^2-m\mathfrak q'-m\mathfrak q^2
   -\frac{\vartheta^2\lvert\mathbf K_m\rvert^2}{2}\ge
 \left(
  1-\frac{9}{50}-\frac{3}{10^4}-\frac{1}{200}
 \right)p_1^2
 >
 \frac45p_1^2.
\end{align*}
For the restriction of \(\Ric\) to
\(\operatorname{span}\{E_2,\ldots,E_{m+1}\}\), first suppose that
\(\mathfrak q\le0\).  The final inequality in the spherical-profile
calculation shows that
\[
 \mathfrak q'+\mathfrak q(\mathfrak p+m\mathfrak q)
 \le\mathfrak q'.
\]
If \(\mathfrak q\ge0\), then the preceding bounds on
\(\mathfrak p\mathfrak q\) and \(|\mathfrak q|\) give
\[
 \mathfrak q(\mathfrak p+m\mathfrak q)
 \le\left(\frac1{100}+\frac3{10^4}\right)\mathfrak p_1^2.
\]
Consequently, in both cases,
\[
q'+q(p+mq)
 <
 \frac{3}{40}p_1^2
 \le
 \frac{3}{640}p_0^2
 <
 \frac{1}{200\rho_0^2}.
\]
Here we used \(\Lambda_{\mathrm{con}}\ge2\) and
\(\mathfrak p_0<1/\rho_0\).  Since \(m-1\ge1\)
and \(\rho\le\rho_0\), the scalar multiple of the identity in
\eqref{eq:un-ricij} is therefore larger than
\(\frac{199}{200\rho_0^2}\); the remaining term
\((\vartheta^2/2)\mathbf K_m\otimes\mathbf K_m\) is nonnegative.

Lemma~\ref{lem:un-ricci} gives
\(\Ric(E_0,E_\alpha)=0\) for \(1\le\alpha\le m+1\), while
\eqref{eq:un-mixed-cancellation} gives
\(\Ric(E_1,E_{i+1})=0\) for \(1\le i\le m\).  Thus \(\Ric\) is block
diagonal with respect to
\[
 \mathbb RE_0\oplus\mathbb RE_1
 \oplus\operatorname{span}\{E_2,\ldots,E_{m+1}\},
\]
and the preceding estimates prove \(\Ric_{g_{\mathcal C}}>0\).

Finally,
\(\mathfrak q(0)=-\mathfrak p_1/100\) and
\(\mathfrak q(s_*)=\mathfrak p_1/100\).  Apply
\eqref{eq:un-shape} with outward normal \(-\partial_s\) on the left and
\(+\partial_s\) on the right to obtain
\eqref{eq:un-II-left}--\eqref{eq:un-II-right}.  The spherical block in
\eqref{eq:un-II-right} is
\((\mathfrak p_1/100)\operatorname{Id}_m>0\).  Its Schur complement is
bounded below by
\[
 p_1-\frac{100}{p_1}\frac{\vartheta(s_*)^2}{4}
 >
 \left(1-\frac14\right)p_1
 =\frac34p_1>0,
\]
where we used \eqref{eq:un-twist-rate-bound} and
\(\lvert\mathbf K_m\rvert\le1\).  Thus
\(\II^{\mathrm{out}}_{s_*}>0\).
\end{proof}

It remains to attach the left boundary of $\mathcal C$ to the twisted
outer boundary of \(W_m\).  Using \(\Phi_o\) on \(\partial_oW_m\) and
the standard product coordinates on
\(\{0\}\times\mathbb S^1\times\mathbb S^m\), both boundaries carry the
metric \(g_{\ell_0,\rho_0,1/2}\).
Lemma~\ref{lem:relative-gluing-unified} therefore reduces the gluing
problem to proving that the sum of their outward second fundamental
forms is positive definite.  At the left boundary, the outward normal of $\mathcal C$
contributes a negative axial term, a positive spherical term, and an
axial--sphere term controlled by \eqref{eq:un-twist-rate-bound}.  We use
the following elementary form of the Schur-complement criterion.  For a
symmetric quadratic form on
\[
 \mathbb RE_1\oplus
 \operatorname{span}\{E_2,\ldots,E_{m+1}\}
\]
with positive lower bounds on its two diagonal restrictions, positivity
follows if the square of half the mixed coefficient is smaller than the
product of those two lower bounds.

\begin{lemma}
\label{lem:un-left-boundary-gluing}
Identify the outer boundary of a quotient \(W_m\), using \(\Phi_o\),
with the boundary
\(\{0\}\times\mathbb S^1\times\mathbb S^m\) of $\mathcal C$ in
Proposition~\ref{prop:un-connector}.  Suppose that the common induced
metric is \(g_{\ell_0,\rho_0,1/2}\).  With respect to the corresponding
axial--sphere splitting, assume that the outward second fundamental
form of \(W_m\) satisfies
\[
 \II_{W_m}(E_1,E_1)\ge\delta_{\mathrm{ax}}>0,\qquad
 \II_{W_m}|_{T\mathbb S^m}>0,\qquad
 \II_{W_m}(E_1,Y)=0
 \quad\text{for all }Y\in T\mathbb S^m.
\]
If
\[
 0<\mathfrak p_0<
 \min\left\{\frac{\delta_{\mathrm{ax}}}{2},
             \frac1{\rho_0}\right\},
\]
then the sum of the two outward second fundamental forms is positive
definite.
\end{lemma}

\begin{proof}
At \(s=0\), the outward normal of $\mathcal C$ is \(-\partial_s\).
For \(X=xE_1+Y\), where \(Y\) is tangent to the sphere factor,
\eqref{eq:un-II-left}, the hypotheses on \(W_m\), and
\(|\mathbf K_m|\le1\) give
\begin{equation}
 \bigl(\II_{W_m}+\II^{\mathrm{out}}_0\bigr)(X,X)
 \ge
 (\delta_{\mathrm{ax}}-\mathfrak p_0)x^2
 -|\vartheta(0)|\,|x|\,|Y|
 +\frac{\mathfrak p_1}{100}|Y|^2.
 \label{eq:un-left-quadratic-lower}
\end{equation}
The parameter assumptions imply
\[
 \delta_{\mathrm{ax}}>2\mathfrak p_0
 =2\Lambda_{\mathrm{con}}^2\mathfrak p_1
 \ge8\mathfrak p_1.
\]
Consequently,
\[
 (\delta_{\mathrm{ax}}-\mathfrak p_0)
 \frac{\mathfrak p_1}{100}
 >\frac{1}{25}\mathfrak p_1^2.
\]
On the other hand, \eqref{eq:un-twist-rate-bound} gives
\[
 \frac{\vartheta(0)^2}{4}
 <
 \frac{p_1^2}{400}.
\]
Thus the square of half the mixed coefficient in
\eqref{eq:un-left-quadratic-lower} is smaller than the product of its
two diagonal coefficients.  The Schur-complement criterion proves that
the right-hand side of \eqref{eq:un-left-quadratic-lower}, and hence the
sum of the two second fundamental forms, is positive definite.
\end{proof}


\section{The building block and infinite telescope}
\label{sec:infinite-telescope}

Recall that a building block in our construction is a compact cobordism with boundary
product markings whose incoming and outgoing product circles satisfy
$[\lambda_{\mathrm{in}}]=[\lambda_{\mathrm{out}}]^2$.  The purpose of
this section is to pass from the local geometric
constructions of the preceding sections to the complete noncompact
manifolds in Theorem~\ref{thm:main}.  We first combine the screw quotient
with a holonomy connector to obtain such a compact Ricci-positive block with untwisted
product boundaries.  The second fundamental form of its incoming
boundary is bounded below by an arbitrarily small negative multiple of
the induced metric, its outgoing boundary is strictly convex, and the
incoming product circle represents the square of the corresponding
outgoing generator.  We then attach the Ricci-positive initial cap and
attach the blocks successively along a ray.  Central subcylinders that
are disjoint from all smoothing collars ensure local finiteness and completeness of
the resulting metric, while the square relation at each stage
determines the fundamental group.

The product parametrization of the boundary in the next proposition are furnished by
Lemma~\ref{lem:un-boundary-components}\hspace{0.25em}and
Proposition~\ref{prop:un-connector}.

\begin{proposition}
\label{prop:square-root-block}
Fix \(m\in\{2,3\}\).  Given
\[
 \ell_I>0,\qquad \rho_I>0,\qquad \delta>0,
 \qquad s_{\min}>0,
\]
there are a compact connected \((m+2)\)-manifold \(B_m\), a metric
\(g_B\) with \(\Ric_{g_B}>0\), and boundary product markings
\[
 \Phi_I:\mathbb S^1\times\mathbb S^m\longrightarrow\partial_iB_m,
 \qquad
 \Phi_O:\mathbb S^1\times\mathbb S^m\longrightarrow\partial_oB_m
\]
with the following properties.

\begin{enumerate}[label=\textup{(\roman*)}]
\item The boundary has exactly the two components
\[
 \partial B_m=\partial_iB_m\sqcup\partial_oB_m,
\]
and
\[
 \Phi_I^*(g_B|_{\partial_iB_m})=g_{\ell_I,\rho_I,0},
 \qquad
 \Phi_O^*(g_B|_{\partial_oB_m})=g_{\ell_O,\rho_O,0}
\]
for some \(\ell_O,\rho_O>0\).

\item With respect to the outward normals,
\[
 \II_{\partial_iB_m}>-\delta\,g_B|_{\partial_iB_m},
 \qquad
 \II_{\partial_oB_m}>0.
\]

\item
\[
 \pi_1(B_m)\cong\mathbb Z=\langle a\rangle.
\]
For any \(y\in\mathbb S^m\), the outgoing product circle
\(\theta\mapsto\Phi_O(\theta,y)\) represents \(a\), whereas the incoming
product circle \(\theta\mapsto\Phi_I(\theta,y)\) represents \(a^2\).
After the loops are transported to one basepoint, these elements are
independent of the transporting paths because \(\pi_1(B_m)\) is
abelian.

\item The block contains a holonomy connector
\[
 \mathcal C=[0,s_*]\times\mathbb S^1\times\mathbb S^m,
 \qquad s_*\ge s_{\min},
\]
whose metric is exactly the metric of
Proposition~\ref{prop:un-connector} on
\[
 [s_*/4,s_*]\times\mathbb S^1\times\mathbb S^m.
\]
In particular, the central subcylinder
\begin{equation}\label{def: central-subcylinder}
 [s_*/4,3s_*/4]\times\mathbb S^1\times\mathbb S^m
\end{equation}
is disjoint from the support of every gluing modification.
\end{enumerate}
\end{proposition}

\begin{proof}
Apply Theorem~\ref{thm:unified-directed-base} with prescribed inner
radius \(R_*=\rho_I\) and with
\[
 0<\varepsilon_*<\delta.
\]
Apply Lemma~\ref{lem:un-warp} with a number
\(0<\delta_*<\delta\), obtaining
\(f=F(1+\alpha\psi+\alpha^2\chi)\).
Choose \(T,F>0\) so that
\[
 \frac{TF}{\pi}=\ell_I.
\]
Form the quotient \(W_m\).  Lemma~\ref{lem:un-boundary-components}
gives
\[
 \Phi_i^*(G|_{\partial_iW_m})=g_{\ell_I,\rho_I,0},
 \qquad
 \Phi_o^*(G|_{\partial_oW_m})
 =g_{\ell_{\mathrm{tw}},\rho_{\mathrm{tw}},1/2},
\]
where
\[
 \ell_{\mathrm{tw}}=\frac{TF(1+\alpha)}{2\pi},
 \qquad \rho_{\mathrm{tw}}=R_o.
\]
On \(\partial_iW_m\), the sphere restriction of the outward second
fundamental form is greater than \(-\varepsilon_*g\), its axial entry is
greater than \(-\delta_*\), and its axial--sphere block is zero.  Hence
\[
 \II_{\partial_iW_m}>-\delta\,G|_{\partial_iW_m}.
\]
Indeed, the form is block diagonal in the axial--sphere splitting, and
each diagonal block has all eigenvalues greater than \(-\delta\).
On \(\partial_oW_m\), the sphere restriction is positive definite, the
axial--sphere block is zero, and the axial principal curvature is
positive.  Compactness therefore gives \(\delta_{\mathrm{ax}}>0\) such
that the latter is at least \(\delta_{\mathrm{ax}}\).

Choose
\[
 0<\mathfrak p_0<
 \min\left\{\frac{\delta_{\mathrm{ax}}}{2},
                  \frac1{\rho_{\mathrm{tw}}}\right\}.
\]
Then choose \(\Lambda_{\mathrm{con}}\) so large that
\[
 \Lambda_{\mathrm{con}}\ge
 \max\left\{2,\frac{40\rho_{\mathrm{tw}}}{\ell_{\mathrm{tw}}}\right\},
 \qquad
 s_*=\frac{\Lambda_{\mathrm{con}}^2-1}{2\mathfrak p_0}
 \ge s_{\min}.
\]
Proposition~\ref{prop:un-connector} gives a Ricci-positive holonomy connector
whose left metric is
\(g_{\ell_{\mathrm{tw}},\rho_{\mathrm{tw}},1/2}\) and whose right metric is
\[
 g_{\Lambda_{\mathrm{con}}\ell_{\mathrm{tw}},\rho_{\mathrm{tw}},0}.
\]
Lemma~\ref{lem:un-left-boundary-gluing} verifies the strict second
fundamental form inequality at \(\partial_oW_m\).  Apply
Lemma~\ref{lem:relative-gluing-unified}, with its support in the collar
$[0,s_*/4)\times\mathbb S^1\times\mathbb S^m$ of $\mathcal C$.  The resulting manifold is
\[
 B_m=W_m\cup_{\Phi_o\circ\iota_0^{-1}}\mathcal C,
 \qquad
 \iota_0(\theta,y)=(0,\theta,y),
\]
with incoming boundary \(\partial_iW_m\), outgoing boundary
\(\{s_*\}\times\mathbb S^1\times\mathbb S^m\), and
\[
 \ell_O=\Lambda_{\mathrm{con}}\ell_{\mathrm{tw}},
 \qquad \rho_O=\rho_{\mathrm{tw}}.
\]
Define the boundary product markings in the statement by
\[
 \Phi_I:=\Phi_i,\qquad
 \Phi_O:=\iota_{s_*},\qquad
 \iota_{s_*}(\theta,y)=(s_*,\theta,y).
\]
The final assertion of Proposition~\ref{prop:un-connector} gives
\(\II_{\partial_oB_m}>0\).

Collapsing the holonomy connector in the \(s\)-direction gives a deformation
retraction \(B_m\to W_m\).  Lemma~\ref{lem:un-boundary-components}
therefore gives
\[
 \pi_1(B_m)\cong\pi_1(W_m)\cong\mathbb Z,
\]
and identifies the incoming circle with \(a^2\).  Under the same
retraction the outgoing circle becomes the outer circle of \(W_m\), so
it represents \(a\).  If two paths are used to move a boundary loop to
a fixed basepoint, the resulting elements are conjugate and hence equal
because the group is abelian.  The support choice above proves the
final assertion.
\end{proof}

Proposition~\ref{prop:square-root-block} provides the block used at each
stage of the construction.  To begin the successive attachments, we need
a compact
Ricci-positive manifold whose boundary has the same untwisted product
form, is strictly convex, and whose circle factor bounds a disk in the
interior.  The following elementary doubly warped cap has precisely
these properties.

\begin{lemma}
\label{lem:un-cap}
For \(m\in\{2,3\}\) and arbitrary \(\ell,\rho>0\), the manifold
\(D^2\times\mathbb S^m\) carries a Ricci-positive metric whose boundary is the
product \(\mathbb S^1_\ell\times\mathbb S^m_\rho\) and is strictly convex.
\end{lemma}

\begin{proof}
Set
\[
 s_{\mathrm{cap}}=\frac{\pi\ell}{3},
 \qquad
 \ell_{\mathrm{cap}}(s)=2\ell\sin\left(\frac{s}{2\ell}\right),
\]
and, for \(\gamma_{\mathrm{cap}}>0\), put
\[
 \rho_{\mathrm{cap}}(s)
 =\rho\exp\left(\gamma_{\mathrm{cap}}
       (s^2-s_{\mathrm{cap}}^2)\right).
\]
On
\([0,s_{\mathrm{cap}}]\times\mathbb S^1\times\mathbb S^m\), consider
\begin{equation}\label{eq:un-cap-metric}
 g_{\mathrm{cap}}(\gamma_{\mathrm{cap}})
 =\dd s^2+\ell_{\mathrm{cap}}(s)^2\dd\theta^2
       +\rho_{\mathrm{cap}}(s)^2g_{\mathbb S^m}.
\end{equation}
At \(s=0\),
\begin{align*}
 \ell_{\mathrm{cap}}(s)
 &=s-\frac{s^3}{24\ell^2}+O(s^5),\\
 \rho_{\mathrm{cap}}(s)
 &=\rho_{\mathrm{cap}}(0)
   \left(1+\gamma_{\mathrm{cap}}s^2+O(s^4)\right).
\end{align*}
Thus the circle-radius function has the required odd expansion with
derivative one, and the sphere-radius function has an even expansion.
Consequently, \eqref{eq:un-cap-metric} extends smoothly across the
center of the \(D^2\)-factor.

The Ricci curvatures on a unit radial vector, a unit circle vector, and
a unit vector tangent to \(\mathbb S^m\) are, respectively,
\begin{align*}
 &-\frac{\ell_{\mathrm{cap}}''}{\ell_{\mathrm{cap}}}
   -m\frac{\rho_{\mathrm{cap}}''}{\rho_{\mathrm{cap}}},\qquad
 -\frac{\ell_{\mathrm{cap}}''}{\ell_{\mathrm{cap}}}
   -m\frac{\ell_{\mathrm{cap}}'\rho_{\mathrm{cap}}'}
           {\ell_{\mathrm{cap}}\rho_{\mathrm{cap}}},\\
 &\frac{m-1}{\rho_{\mathrm{cap}}^2}
       -\frac{\rho_{\mathrm{cap}}''}{\rho_{\mathrm{cap}}}
       -(m-1)\left(
          \frac{\rho_{\mathrm{cap}}'}{\rho_{\mathrm{cap}}}\right)^2
       -\frac{\ell_{\mathrm{cap}}'\rho_{\mathrm{cap}}'}
              {\ell_{\mathrm{cap}}\rho_{\mathrm{cap}}}.
\end{align*}
When \(\gamma_{\mathrm{cap}}=0\), these three quantities are
\[
 \frac1{4\ell^2},\qquad
 \frac1{4\ell^2},\qquad
 \frac{m-1}{\rho^2}.
\]
They are strictly positive because \(m\ge2\).  As
\(\gamma_{\mathrm{cap}}\to0\), the metrics
\eqref{eq:un-cap-metric} converge smoothly on the compact manifold
\(D^2\times\mathbb S^m\) to the product metric obtained at
\(\gamma_{\mathrm{cap}}=0\).  Hence
\(\Ric_{g_{\mathrm{cap}}(\gamma_{\mathrm{cap}})}>0\) for every
sufficiently small \(\gamma_{\mathrm{cap}}>0\).

At \(s=s_{\mathrm{cap}}\),
\[
 \ell_{\mathrm{cap}}(s_{\mathrm{cap}})=\ell,
 \qquad
 \rho_{\mathrm{cap}}(s_{\mathrm{cap}})=\rho.
\]
Thus the induced boundary metric is
\(\ell^2\dd\theta^2+\rho^2g_{\mathbb S^m}\).  With respect to the
outward unit normal \(+\partial_s\), the boundary principal curvatures
in the circle and sphere directions are
\[
 \frac{\ell_{\mathrm{cap}}'(s_{\mathrm{cap}})}
      {\ell_{\mathrm{cap}}(s_{\mathrm{cap}})}
 =\frac{\sqrt3}{2\ell}>0,
 \qquad
 \frac{\rho_{\mathrm{cap}}'(s_{\mathrm{cap}})}
      {\rho_{\mathrm{cap}}(s_{\mathrm{cap}})}
 =2\gamma_{\mathrm{cap}}s_{\mathrm{cap}}>0,
\]
where the second value has multiplicity \(m\).  Therefore the boundary
is strictly convex.
\end{proof}

We now start our construction. First,
we fix a sufficiently small $\gamma_{\mathrm{cap}}>0$ as in the proof and
write $g_{\mathrm{cap}}:=g_{\mathrm{cap}}(\gamma_{\mathrm{cap}})$.  The
Ricci-positive manifold
$(D^2\times\mathbb S^m,g_{\mathrm{cap}})$ is the initial cap introduced
in the introduction.  For every $y\in\mathbb S^m$, its boundary product
circle bounds the disk $D^2\times\{y\}$.

We now use the initial cap to begin the construction and attach the building
blocks inductively.  At each stage, we choose the parameter in the
estimate \(\II_{\partial_iB_m}>-\delta g\) to be smaller than half the
least outward principal curvature of the preceding block.  The relative
gluing lemma therefore applies at every interface.
Meanwhile, the holonomy-connector lengths are chosen to tend rapidly to
infinity.  In each block, the central subcylinder defined in (\ref{def: central-subcylinder}) between the levels
$s_j/4$ and $3s_j/4$ is disjoint from every smoothing collar.

Fix \(\ell_0,\rho_0>0\), and take the initial cap from
Lemma~\ref{lem:un-cap} with boundary metric \(g_{\ell_0,\rho_0,0}\).
Let \(\lambda_{\mathrm{cap}}>0\) be its least outward boundary principal
curvature.  Apply Proposition~\ref{prop:square-root-block} with this
incoming metric and with
\[
 0<\delta_0<\frac{\lambda_{\mathrm{cap}}}{2},
 \qquad s_{\min}=4,
\]
and denote the resulting block by \(B_0\).  Under the boundary product
markings, the sum of the two outward second fundamental forms is
bounded below by
\((\lambda_{\mathrm{cap}}-\delta_0)g>0\).
Lemma~\ref{lem:relative-gluing-unified} therefore joins the initial cap to
\(B_0\).

Inductively, suppose the outgoing boundary of \(B_j\) has product metric
\[
 g_{\ell_{j+1},\rho_{j+1},0}
\]
and least outward principal curvature \(\lambda_j>0\).  Construct
\(B_{j+1}\) with this incoming metric and with
\[
 0<\delta_{j+1}<\frac{\lambda_j}{2},
 \qquad s_{\min}=2^{j+3}.
\]
The identity in the product coordinates is an isometry of the boundary
metrics, and the sum of the outward second fundamental forms is bounded
below by \((\lambda_j-\delta_{j+1})g>0\).  The relative gluing lemma
therefore applies.  Notice that no uniform positive lower bound for
\(\lambda_j\) is needed.

Let
\begin{equation}\label{def:telescope}
 M=(D^2\times\mathbb S^m)\cup B_0\cup B_1\cup B_2\cup\cdots
\end{equation}
denote the manifold obtained after all these identifications.  We call
this union the \emph{infinite telescope}.  For each
\(j\ge0\), write
\[
 \mathcal C_j=[0,s_j]\times\mathbb S^1\times\mathbb S^m
 \subset B_j
\]
for the holonomy connector supplied by
Proposition~\ref{prop:square-root-block}.  The preceding choices give
\[
 s_j\ge2^{j+2}.
\]
The internal smoothing in \(B_j\) is supported to the left of
\(\{s_j/4\}\times\mathbb S^1\times\mathbb S^m\).  At the outgoing
boundary of \(B_j\), choose the support of the next gluing inside the
collar
\[
 (3s_j/4,s_j]\times\mathbb S^1\times\mathbb S^m.
\]
Consequently, we may denote the central subcylinder of $B_j$ by
\begin{equation}\label{eq:un-central-slab}
 \mathcal P_j
 :=[s_j/4,3s_j/4]\times\mathbb S^1\times\mathbb S^m
 \subset\mathcal C_j.
\end{equation}
The final metric on $\mathcal P_j$ is exactly the holonomy-connector
metric of Proposition~\ref{prop:un-connector}.

Set
\[
 \Sigma_j=\{s_j/2\}\times\mathbb S^1\times\mathbb S^m
 \subset\mathcal C_j,
\]
and denote the two closed subcylinders cut out by $\Sigma_j$ by
\[
 \mathcal C_j^-
 :=[0,s_j/2]\times\mathbb S^1\times\mathbb S^m,
 \qquad
\mathcal C_j^+
 :=[s_j/2,s_j]\times\mathbb S^1\times\mathbb S^m.
\]
Let \(M_N\) be the closure of the component of
\(M\setminus\Sigma_N\) containing the initial cap.  The dual graph of the
initial cap and the blocks is a ray, so \(\Sigma_N\) separates a finite initial
segment from the connected infinite tail.  Hence
\begin{equation}\label{eq:telescope-exhaustion}
 \partial M_N=\Sigma_N,\qquad
 M_N\subset\operatorname{int}M_{N+1},\qquad
 \bigcup_{N\ge0}\operatorname{int}M_N=M,
\end{equation}
and \(M\setminus M_N\) is connected.  Since \(M_N\) is contained in the
union of the initial cap and finitely many compact blocks, it is compact.

The open sets \(\operatorname{int}M_N\) form a nested cover of \(M\).
Every compact subset is therefore contained in one of them and meets
only finitely many gluing collars.  Thus the smoothing operations are
locally finite and define a smooth metric on \(M\) with
\(\Ric_M>0\).

This completes the geometric assembly, and then we check the topology
of \(M\).  The hypersurfaces \(\Sigma_N\) separate the infinite
telescope into compact initial segments with connected complements,
which identifies the end structure.  Applying van Kampen's theorem to
the same exhaustion then records the square relation introduced by each
successive block.
The complement \(M\setminus M_N\) is the connected tail beginning at
the midpoint of \(\mathcal C_N\).  The exhaustion
\eqref{eq:telescope-exhaustion} consequently proves that \(M\) has one
end.

For \(j<N\), let \(a_j\) be the generator represented by the outgoing
product circle of \(B_j\).  Let \(a_N\) be represented by a product
circle in \(\Sigma_N\).  The product structure of the holonomy connector
\(\mathcal C_N\) identifies this
class with the outgoing generator of \(B_N\).
For \(0\le j<N\), Proposition~\ref{prop:square-root-block} says that
the incoming product circle of \(B_{j+1}\) represents \(a_{j+1}^2\).
Therefore the boundary identification between \(B_j\) and \(B_{j+1}\)
gives
\[
 a_j=a_{j+1}^2.
\]
The initial cap is simply connected, and its boundary product circle bounds
the \(D^2\)-factor.  Since this circle is glued to the incoming circle
of \(B_0\), it gives \(a_0^2=1\).  The subcylinder
$\mathcal C_N^-$ deformation retracts to its left boundary.  Repeated
 applications of the Seifert--van Kampen theorem therefore give
\begin{align}
 \pi_1(M_N)
 &=\left\langle a_0,\ldots,a_N\ \middle|\
   a_0^2=1,\ a_j=a_{j+1}^2\ (0\le j<N)\right\rangle\notag\\
 &\cong\langle a_N\mid a_N^{2^{N+1}}=1\rangle
 \cong\mathbb Z/2^{N+1}\mathbb Z.
 \label{eq:un-finite-pi1}
\end{align}
Under these identifications, inclusion sends
\[
 \pi_1(M_N)\longrightarrow\pi_1(M_{N+1}),
 \qquad [k]\longmapsto[2k].
\]
This homomorphism is injective.  Every loop and every homotopy of loops
has compact image and is contained in some \(M_N\).  Hence
\begin{equation}\label{eq:un-prufer}
 \pi_1(M)
 =\varinjlim_N\bigl(\mathbb Z/2^{N+1}\mathbb Z,[k]\mapsto[2k]\bigr)
 =\CPrufer.
\end{equation}

It remains to verify that the resulting Ricci-positive metric is
complete.  This is where the quantitative lower bounds on the
holonomy-connector lengths enter: every curve escaping through the
telescope must cross $\mathcal P_j$ for every sufficiently large $j$.
The sum of the corresponding lower bounds diverges.  We conclude by
recording orientability and completing the proof of the main theorem.

Let \(\zeta\) be a piecewise smooth curve joining
\(D^2\times\mathbb S^m\) to \(M\setminus M_N\).  It crosses
$\mathcal P_j$ for every $0\le j<N$.  On each such slab the metric
contains the orthogonal summand \(\dd s^2\).  Parameterizing a crossing
by \(t\), its \(s\)-coordinate traverses
\([s_j/4,3s_j/4]\), so the portion of \(\zeta\) in that slab has length
at least
\[
 \int|s'(t)|\,\dd t\ge\frac{s_j}{2}.
\]
The slabs $\mathcal P_j$ are pairwise disjoint.  Summing these contributions
and then taking the infimum over \(\zeta\) gives
\[
 d(D^2\times\mathbb S^m,M\setminus M_N)
 \ge\frac12\sum_{j=0}^{N-1}s_j
 \ge\frac12\sum_{j=0}^{N-1}2^{j+2}\longrightarrow\infty.
\]
Fix \(x\in M\) and \(R>0\), and choose
\(p\in D^2\times\mathbb S^m\).  For all sufficiently large \(N\),
\[
 d(D^2\times\mathbb S^m,M\setminus M_N)>R+d(p,x).
\]
For \(y\in M\setminus M_N\), the triangle inequality gives
\[
 d(x,y)\ge d(p,y)-d(p,x)>R.
\]
Consequently, \(\overline{B}_R(x)\subset M_N\).  Since \(M_N\) is
compact, every closed bounded metric ball is compact.  Thus the length
metric is proper, and the Hopf--Rinow theorem implies geodesic
completeness.

Finally, \(P_m\) is connected and orientable.  The orientation sign of
the diffeomorphism \(\sigma_m\) is therefore constant, so it is enough
to compute the sign of its differential at one fixed point retained in
the directed-base construction.  At that point,
\[
 d\sigma_2=\operatorname{diag}(-1,1,-1),
 \qquad d\sigma_3=-\operatorname{Id}_{\mathbb R^4},
\]
and both determinants are \(+1\).  Thus each screw quotient, holonomy
connector, and initial cap is orientable.  Since the dual gluing graph is
a ray, their
orientations can be chosen successively so that every boundary
identification reverses boundary orientations.  Hence \(M\) is
orientable.

\begin{proof}[Proof of Theorem~\ref{thm:main}]
For \(m\in\{2,3\}\), the preceding locally finite construction gives a
smooth one-ended \((m+2)\)-manifold with \(\Ric>0\).  Equation
\eqref{eq:un-prufer} identifies its fundamental group with \(\CPrufer\),
and the estimates on the slabs $\mathcal P_j$ prove completeness.  The
cases \(m=2\)
and \(m=3\) have dimensions four and five, respectively.
\end{proof}


\section{Further properties of the constructed manifolds}
\label{sec:further-properties}

Fix \(m\in\{2,3\}\), set \(n=m+2\), and let
\((M^n,g)\) be one of the complete Ricci-positive manifolds constructed
in Section~7.  This section describes both the topology of its
universal cover and the large-scale geometry arising from the free
connector parameters.  Writing \(\widehat M_N\) for the distinguished
lift of the \(N\)-th compact manifold, we first prove that
\[
 \widehat M_N
 \cong
 (D^2\times\mathbb S^m)
 \#^{\,2^{N+1}-1}
 (\mathbb S^2\times\mathbb S^m),
\]
and pass to the limit to obtain
\[
 \widetilde M
 \cong
 \mathbb R^{m+2}\#_{\infty}
 (\mathbb S^2\times\mathbb S^m).
\]
We then show that, for every prescribed sequence
\(e_j\in[1,3/2]\), the connector parameters can be chosen so that
there exist sequences \(r_j,\widetilde r_j\to\infty\) satisfying
\[
 \frac{\log\operatorname{Vol}_M(B_{r_j}(p))}{\log r_j}
 =e_j+o(1),
 \qquad
 \frac{\log\operatorname{Vol}_{\widetilde M}
 (B_{\widetilde r_j}(\widetilde p))}
 {\log\widetilde r_j}
 =e_j+o(1).
\]
Finally, we prove that every admissible metric has
\[
 \operatorname{AVR}(M)=0,
 \qquad
 \limsup_{r\to\infty}
 \frac{\operatorname{Vol}_M(B_r(p))}{r}
 =\infty.
\]
For the universal cover, the asymptotic volume ratio vanishes for every
parameter choice when \(m=2\), and for the alternating parameter
choices when \(m=2\) or \(m=3\).  For arbitrary parameter choices when
\(m=3\), its value is not determined by the present argument.


\subsection{Topology of the universal cover}
\label{subs:topology-of-the-universal-cover}


Let \(q:\widetilde M\to M\) be the universal covering map.  Fix a point
\(p\) in the initial cap and a point \(\widetilde p\in q^{-1}(p)\).  For
\(N\ge0\), let \(\widehat M_N\) be the component of \(q^{-1}(M_N)\)
containing \(\widetilde p\), where \(M_N\) is defined in
\eqref{eq:telescope-exhaustion}.  We also write \(\widehat M_{-1}\) for
the component of the preimage of the initial cap containing
\(\widetilde p\); thus \(\widehat M_{-1}\cong D^2\times\mathbb S^m\).

We first make the boundary parametrizations used below explicit.  Fix
the round identifications
\(\iota_o:\mathbb S^m\to\partial_oP_m\) and
\(\iota_+:\mathbb S^m\to I_+\) used in
Lemma~\ref{lem:un-boundary-components}, and put
\(\iota_-:=\sigma_m\circ\iota_+\).  By
Theorem~\ref{thm:unified-directed-base},
\begin{equation}\label{eq:universal-boundary-parametrizations}
 \sigma_m\circ\iota_o
 =\iota_o\circ\exp(\pi K_m),
 \qquad
 \sigma_m\circ\iota_+=\iota_-.
\end{equation}
If \(T_j\) is the period of the screw quotient in \(B_j\), its outer
and inner boundary product parametrizations are therefore
\begin{equation}\label{eq:universal-screw-boundary-maps}
 \begin{aligned}
  \Phi_{o,j}(e^{\mathrm i\theta},y)
  &=\left[\frac{T_j\theta}{2\pi},
    \iota_o\!\left(\exp\left(\frac\theta2K_m\right)y\right)\right],\\
  \Phi_{i,j}(e^{\mathrm i\theta},y)
  &=\left[\frac{T_j\theta}{\pi},\iota_+(y)\right].
 \end{aligned}
\end{equation}
The first identity in \eqref{eq:universal-boundary-parametrizations}
makes \(\Phi_{o,j}\) well defined, and the second shows that the first
return to \(I_+\) occurs after translation by \(2T_j\).  The holonomy
connector is a cylinder with the displayed product coordinate, so it
transports the circle and sphere coordinates in
\eqref{eq:universal-screw-boundary-maps} to every intermediate slice,
including \(\Sigma_j\), and to the outgoing boundary of \(B_j\).

The following elementary calculation is the only surgery result needed
in the proof.

\begin{lemma}\label{lem:universal-two-core}
Let
\[
 Y_\pm\cong
 (D^2\times\mathbb S^m)
 \mathbin{\#}^{\,r_\pm}(\mathbb S^2\times\mathbb S^m),
 \qquad r_\pm\ge0,
\]
where the connected sums are disjoint from fixed product collars of
\(\partial Y_\pm\).  In the fixed product coordinates, suppose the maps
from \(\partial Y_\pm\) to
\(\mathbb S^1\times I_\pm\subset\mathbb S^1\times P_m\), after
composition with \(\theta\)-independent product diffeomorphisms extending
over \(Y_\pm\), have the form
\begin{equation}\label{eq:universal-attaching-maps}
 \varphi_\pm(e^{\mathrm i\theta},y)
 =\bigl(e^{\mathrm i\theta},\alpha_\pm(\theta)y\bigr),
 \qquad
 \alpha_\pm:\mathbb S^1\longrightarrow SO(m+1).
\end{equation}
Then
\begin{equation}\label{eq:universal-two-core-result}
 Y_+\cup_{\varphi_+}(\mathbb S^1\times P_m)
 \cup_{\varphi_-}Y_-
 \cong
 (D^2\times\mathbb S^m)
 \mathbin{\#}^{\,r_++r_-}(\mathbb S^2\times\mathbb S^m)
 \mathbin{\#}S(\xi),
\end{equation}
where \(S(\xi)\to\mathbb S^2\) is the oriented
\(\mathbb S^m\)-bundle determined by the boundary identification
\begin{equation}\label{eq:universal-clutching-map}
 (e^{\mathrm i\theta},y)
 \longmapsto
 \bigl(e^{-\mathrm i\theta},
  \alpha_-(\theta)\alpha_+(\theta)^{-1}y\bigr).
\end{equation}
In particular, if \([\alpha_+]=[\alpha_-]\) in
\(\pi_1SO(m+1)\), then \(S(\xi)\cong\mathbb S^2\times\mathbb S^m\).
If \(\alpha_+=\alpha_-\) as maps, their common loop is transported to
the boundary of the \(D^2\times\mathbb S^m\) factor on the right-hand
side of \eqref{eq:universal-two-core-result}.
Moreover, if a diffeomorphism on \(Y_+\) has already been fixed and is
a product on its boundary collar, the diffeomorphism in
\eqref{eq:universal-two-core-result} can be chosen to agree with it on
all of \(Y_+\).
\end{lemma}

\begin{proof}
Relative to the boundary components \(I_+\sqcup I_-\), the manifold
\(P_m\) is obtained from
\[
 (I_+\sqcup I_-)\times[0,1]
\]
by attaching a single \(1\)-handle \(D^1\times D^m\) along
\(\mathbb S^0\times D^m\).  Consequently, the nonproduct part of
\(\mathbb S^1\times P_m\) is the round handle
\[
 \mathbb S^1\times D^1\times D^m.
\]
Equipping \(\mathbb S^1\) with one \(0\)-cell and one \(1\)-cell
decomposes this round handle into an ordinary index-one handle
\(D^1\times D^{m+1}\) and an ordinary index-two handle
\(D^2\times D^m\).  The index-one handle joins \(Y_+\) to \(Y_-\),
while the index-two handle is attached along the boundary of its
standard core disk.

Equivalently, fill the two deleted disks in \(P_m\) and regard the two
attachments as framed surgeries on the parallel circles
\[
 \mathbb S^1\times\{p_+\},
 \qquad
 \mathbb S^1\times\{p_-\}
 \subset \mathbb S^1\times D^{m+1}.
\]
First perform the surgery corresponding to \(Y_+\).  The product of
\(\mathbb S^1\) with a straight arc from \(p_-\) to \(p_+\), capped
off along its \(p_+\)-boundary by a meridional disk introduced by the
first surgery, is precisely the standard core disk of the index-two
handle above.  Transporting the stabilized normal framing along the
product annulus gives \(\alpha_+\), whereas the framing prescribed for
the second surgery is \(\alpha_-\).  The relative framing of the second
surgery is therefore
\[
 \alpha_-\alpha_+^{-1}.
\]

A regular neighborhood of this core disk, together with the second
surgery piece, is diffeomorphic to the complement of an
\((m+2)\)-disk in the sphere bundle obtained by gluing two copies of
\(D^2\times\mathbb S^m\) via
\eqref{eq:universal-clutching-map}.  Reintroducing the connected-sum
factors, whose supports were chosen away from the fixed collars, gives
\eqref{eq:universal-two-core-result}.  If
\([\alpha_+]=[\alpha_-]\) in \(\pi_1SO(m+1)\), then the relative
clutching loop is null-homotopic, and hence the resulting sphere bundle
is diffeomorphic to \(\mathbb S^2\times\mathbb S^m\).

All the preceding handle modifications are supported in
\(\mathbb S^1\times P_m\) and on the outward sides of the two fixed
collars.  They can therefore be carried out without changing any
prescribed product diffeomorphism on \(Y_+\), proving the final
assertion.  The first surgery also shows that, when
\(\alpha_+=\alpha_-\), their common framing is carried to the
remaining boundary, whereas the closed sphere-bundle summand depends
only on the relative loop \(\alpha_-\alpha_+^{-1}\).  This proves the
preceding assertion.
\end{proof}

\begin{theorem}\label{thm:universal-cover-topology}
For every \(N\ge0\),
\[
 \widehat M_N\cong
 (D^2\times\mathbb S^m)
 \mathbin{\#}^{\,2^{N+1}-1}
 (\mathbb S^2\times\mathbb S^m).
\]
Consequently,
\[
 \widetilde M\cong
 \mathbb R^{m+2}\mathbin{\#}_{\infty}
 (\mathbb S^2\times\mathbb S^m).
\]
In particular, \(\widetilde M\) has one end and is simply connected at
infinity.
\end{theorem}

\begin{proof}
We first determine the finite lifted regions and their boundary
parametrizations.  By \eqref{eq:un-finite-pi1},
\(\pi_1(M_N)\cong\mathbb Z/2^{N+1}\mathbb Z\), and its inclusion into
\(\pi_1(M)\) is injective.  The restriction of \(q\) is therefore the
universal finite covering
\begin{equation}\label{eq:universal-finite-cover}
 q:\widehat M_N\longrightarrow M_N,
 \qquad \deg(q|_{\widehat M_N})=2^{N+1}.
\end{equation}
In particular, \(\widehat M_N\) is compact and simply connected.

Consider the lift of the screw quotient in \(B_N\) which occurs in
\(\widehat M_N\).  Its product identification is
\begin{equation}\label{eq:universal-cyclic-lift}
 \frac{\mathbb R\times P_m}
 {(t,x)\sim(t+2^{N+1}T_N,x)}
 \cong\mathbb S^1\times P_m,
 \qquad
 (e^{\mathrm i\theta},x)
 \longmapsto
 \left[\frac{2^{N+1}T_N\theta}{2\pi},x\right].
\end{equation}
The two incoming boundary components have no
\(\theta\)-dependent sphere rotation.  Substituting
\(2^{N+1}\theta\) for the circle parameter in the first formula of
\eqref{eq:universal-screw-boundary-maps} shows that the outgoing
boundary has the exact rotation
\begin{equation}\label{eq:universal-lifted-twist}
 \theta\longmapsto\exp(2^N\theta K_m).
\end{equation}

Set
\begin{equation}\label{eq:universal-boundary-framing}
 \alpha_{-1}(\theta)=\operatorname{Id},
 \qquad
 \alpha_N(\theta)
 =\exp\bigl((2^{N+1}-1)\theta K_m\bigr)
 \quad(N\ge0).
\end{equation}
We prove by induction, beginning with
\(\widehat M_{-1}=D^2\times\mathbb S^m\), that
\begin{equation}\label{eq:universal-finite-core}
 \widehat M_N\cong
 (D^2\times\mathbb S^m)
 \mathbin{\#}^{\,2^{N+1}-1}
 (\mathbb S^2\times\mathbb S^m),
\end{equation}
and that its outgoing boundary is parametrized by
\((e^{\mathrm i\theta},y)\mapsto
(e^{\mathrm i\theta},\alpha_N(\theta)y)\) relative to the residual
\(D^2\times\mathbb S^m\) factor.

At level \(N\), the preimage of \(M_{N-1}\) in
\(\widehat M_N\) has two components, each diffeomorphic to
\(\widehat M_{N-1}\).  They are attached to the two incoming
boundaries of the product in \eqref{eq:universal-cyclic-lift}.  On the
first component use the circle parameter \(\theta\).  The natural
parameter on the second differs by
\(\theta\mapsto\theta-\pi/2^N\).  Since \(\alpha_{N-1}\) lies in the
one-parameter subgroup generated by \(K_m\), this translation multiplies
it by a constant orthogonal transformation.  Compose only the second
component with the corresponding \(\theta\)-independent product
diffeomorphism, which extends over its residual
\(D^2\times\mathbb S^m\) factor.  We leave the first component and the
product identification \eqref{eq:universal-cyclic-lift} fixed.  Thus
both maps have exactly the form \eqref{eq:universal-attaching-maps} with
\(\alpha_+=\alpha_-=\alpha_{N-1}\), and no constant factor is introduced
into the outgoing parametrization.

Lemma~\ref{lem:universal-two-core} now gives one new untwisted
\(\mathbb S^2\times\mathbb S^m\) factor.  The two copies of
\(\widehat M_{N-1}\) contribute \(2(2^N-1)\) factors, so their total
number is \(2^{N+1}-1\).  The common incoming framing passes to the
unfilled boundary, while the relative framing alone determines the
closed factor.  Combining it with \eqref{eq:universal-lifted-twist}
gives the exact recurrence
\begin{equation}\label{eq:universal-framing-recurrence}
 \alpha_N(\theta)
 =\exp(2^N\theta K_m)\alpha_{N-1}(\theta)
 =\exp\bigl((2^{N+1}-1)\theta K_m\bigr).
\end{equation}
This proves \eqref{eq:universal-finite-core} with its stated boundary
parametrization.  Notice that for \(m=2\) the common loop
\(\alpha_N\) represents the nonzero element of \(\pi_1SO(3)\), because
\(2^{N+1}-1\) is odd.  This does not create a twisted sphere bundle:
the clutching loop in \eqref{eq:universal-clutching-map} is the quotient
of two identical loops.  For \(m=3\), the simultaneous rotations in the
two orthogonal planes give a null-homotopic loop in \(SO(4)\).

We next pass to the full universal cover.  The inclusions
\(M_N\subset\operatorname{int}M_{N+1}\) imply
\(\widehat M_N\subset\operatorname{int}\widehat M_{N+1}\).  Moreover,
\begin{equation}\label{eq:universal-cover-exhaustion}
 \widetilde M=\bigcup_{N\ge-1}\widehat M_N.
\end{equation}
Indeed, any \(\widetilde x\in\widetilde M\) lies in a translate
\(\gamma\widehat M_j\) for some \(j\).  Equation
\eqref{eq:un-prufer} gives
\(\pi_1(M)=\bigcup_{N\ge0}\langle a_N\rangle\), so for some
\(N\ge j\) the deck transformation \(\gamma\) belongs to
\(\langle a_N\rangle\).  This subgroup is precisely the stabilizer of
\(\widehat M_N\).  Hence
\(\gamma\widehat M_j\subset\gamma\widehat M_N=\widehat M_N\), which
proves \eqref{eq:universal-cover-exhaustion}.

It remains to identify this nested union.  Split each copy of
\(\mathbb S^2\times\mathbb S^m\) into two copies of
\(D^2\times\mathbb S^m\) along
\(\mathbb S^1\times\mathbb S^m\).  Starting with
\(Y_0=D^2\times\mathbb S^m\), complete its boundary with one punctured
half of the first copy and attach the other punctured half of the next
copy along the resulting spherical boundary.  Repeating this operation
gives compact manifolds
\[
 Y_r\cong
 (D^2\times\mathbb S^m)
 \mathbin{\#}^{\,r}(\mathbb S^2\times\mathbb S^m),
 \qquad
 Y_r\subset\operatorname{int}Y_{r+1}.
\]
Give the complementary half which completes each sphere product the
same boundary parametrization as the preceding residual half.  Its
relative clutching loop is then trivial.  The boundary parametrization
of the next residual half is independent of this attachment; at
\(r=2^{N+1}-1\), choose it to be the one determined by \(\alpha_N\).

The relative assertion in Lemma~\ref{lem:universal-two-core} now gives,
inductively, diffeomorphisms
\[
 F_N:\widehat M_N\longrightarrow Y_{2^{N+1}-1},
 \qquad
 F_N|_{\widehat M_{N-1}}=F_{N-1},
\]
which are product maps on the boundary collars.  Consequently the maps
\(F_N\) define a smooth diffeomorphism from \(\widetilde M\) onto
\(\bigcup_{r\ge0}Y_r\).

To identify the latter union with the stated proper connected sum,
choose pairwise disjoint balls tending to infinity along a proper ray
in \(\mathbb R^{m+2}\), and replace the \(r\)-th ball by a punctured
copy of \(\mathbb S^2\times\mathbb S^m\).  Splitting every inserted
copy into its two \(D^2\times\mathbb S^m\) halves identifies this
usual exhaustion, piece by piece and relative to the boundary collars,
with \((Y_r)\).  Therefore
\[
 \bigcup_{r\ge0}Y_r
 \cong
 \mathbb R^{m+2}\mathbin{\#}_{\infty}
 (\mathbb S^2\times\mathbb S^m).
\]
The complementary regions of this exhaustion are connected.  They are
also simply connected: each punctured half is simply connected for
\(m\ge2\), and van Kampen's theorem applies successively across the
common \(\mathbb S^1\times\mathbb S^m\) and
\(\mathbb S^{m+1}\) boundaries.  Thus the union has one end and is
simply connected at infinity.
\end{proof}

\subsection{Flexible volume growth}
For the $j$-th holonomy connector
\[
 \mathcal C_j=[0,s_j]\times\mathbb S^1\times\mathbb S^m,
\]
let $\ell_j,\rho_j,\omega_j$ denote the profiles appearing in the metric
\eqref{eq:un-connector} on $\mathcal C_j$.  Denote by
$\mathfrak p_{0,j}$ and $\Lambda_j$ the parameters corresponding to
$\mathfrak p_0$ and $\Lambda_{\mathrm{con}}$ in
Proposition~\ref{prop:un-connector}, and set
\[
 \mathfrak p_{1,j}
 =\frac{\mathfrak p_{0,j}}{\Lambda_j^2},
 \qquad
 s_j
 =\frac{\Lambda_j^2-1}{2\mathfrak p_{0,j}}.
 \tag{8.13}\label{eq:indexed-connector-parameters}
\]
Write
\[
 g_{\mathcal C_j}
 :=
 \dd s^2+g_{\ell_j(s),\rho_j(s),\omega_j(s)}
\]
for the model metric on $\mathcal C_j$.  Its radius profiles are
\[
 \ell_j(s)
 =\ell_j(0)\sqrt{1+2\mathfrak p_{0,j}s},
 \qquad
 \rho_j(s)
 =\rho_j(0)
 \exp\left(
   -\frac{\mathfrak p_{1,j}}{100s_j}s(s_j-s)
 \right).
 \tag{8.14}\label{eq:indexed-connector-functions}
\]
Recall from \eqref{eq:un-central-slab} that
$\mathcal P_j=[s_j/4,3s_j/4]\times\mathbb S^1\times\mathbb S^m$ is the
protected central slab of $B_j$.  The final metric on $\mathcal P_j$
agrees exactly with $g_{\mathcal C_j}$.  For
$0\le a<b\le1$, abbreviate
\[
 \mathcal C_j[a,b]
 :=
 [as_j,bs_j]\times\mathbb S^1\times\mathbb S^m.
\]

\begin{lemma}
\label{lem:protected-connector-volume}
For $\Lambda\ge2$ and $0\le x\le1$, define
\[
 \mathcal I_m(\Lambda,x)
 :=
 \int_0^x
 \sqrt{1+(\Lambda^2-1)t}\,
 \exp\left(
   -\frac{m(\Lambda^2-1)}{200\Lambda^2}t(1-t)
 \right)\dd t.
 \tag{8.15}\label{eq:dimensionless-volume-function}
\]
For every $0\le a<b\le1$, the model metric $g_{\mathcal C_j}$ satisfies
\[
 \begin{aligned}
 \operatorname{Vol}_{g_{\mathcal C_j}}
 \bigl(\mathcal C_j[a,b]\bigr)
 &=
 2\pi\operatorname{Vol}(\mathbb S^m)\,
 \ell_j(0)\rho_j(0)^m s_j \times
 \bigl(
   \mathcal I_m(\Lambda_j,b)
   -\mathcal I_m(\Lambda_j,a)
 \bigr).
 \end{aligned}
 \tag{8.16}\label{eq:protected-connector-volume}
\]
The same identity holds for the final metric whenever
$1/4\le a<b\le3/4$.

For fixed $0\le a<b\le1$,
\[
 \begin{aligned}
 \lim_{\Lambda\to\infty}&
 \frac{
   \mathcal I_m(\Lambda,b)-\mathcal I_m(\Lambda,a)
 }{\Lambda}
 =
 C_{m,a,b},\\
 C_{m,a,b}
 &:=
 \int_a^b
 \sqrt t\,
 \exp\left(-\frac{m}{200}t(1-t)\right)\dd t
 \in(0,\infty).
 \end{aligned}
 \tag{8.17}\label{eq:dimensionless-volume-asymptotic}
\]
Equivalently,
\[
 \mathcal I_m(\Lambda,b)-\mathcal I_m(\Lambda,a)
 =
 \Lambda\bigl(C_{m,a,b}+o(1)\bigr).
\]
Consequently, if $1/4\le a<b\le3/4$ are fixed and
$\Lambda_j\to\infty$, then
\[
 \operatorname{Vol}
 \bigl(\mathcal C_j[a,b]\bigr)
 \sim
 \pi\operatorname{Vol}(\mathbb S^m)\,C_{m,a,b}\,
 \ell_j(0)\rho_j(0)^m
 \frac{\Lambda_j^3}{\mathfrak p_{0,j}}.
 \tag{8.18}\label{eq:connector-volume-scale}
\]
In particular, this volume is comparable, with constants depending
only on $m,a,b$, to
\[
 \ell_j(0)\rho_j(0)^m
 \frac{\Lambda_j^3}{\mathfrak p_{0,j}}.
\]
\end{lemma}

\begin{proof}
Relative to the circle direction and a frame tangent to
$\mathbb S^m$, the twist in
$g_{\ell_j(s),\rho_j(s),\omega_j(s)}$ is represented by a triangular
change of basis with determinant one.  It therefore does not affect the
volume density, and
\[
 \dd V_{g_{\mathcal C_j}}
 =
 \ell_j(s)\rho_j(s)^m\,
 \dd s\,\dd\theta\,\dd V_{g_{\mathbb S^m}}.
\]
Substituting $s=ts_j$ and using
\eqref{eq:indexed-connector-parameters} and
\eqref{eq:indexed-connector-functions}, we obtain
\[
 \begin{aligned}
 \operatorname{Vol}_{g_{\mathcal C_j}}
 \bigl(\mathcal C_j[a,b]\bigr)
 &=
 2\pi\operatorname{Vol}(\mathbb S^m)\,
 \ell_j(0)\rho_j(0)^m s_j\\
 &\quad\times
 \int_a^b
 \sqrt{1+(\Lambda_j^2-1)t}\,
 \exp\left(
   -\frac{m(\Lambda_j^2-1)}{200\Lambda_j^2}t(1-t)
 \right)\dd t,
 \end{aligned}
\]
which is \eqref{eq:protected-connector-volume}.

To prove \eqref{eq:dimensionless-volume-asymptotic}, divide the
integrand in \eqref{eq:dimensionless-volume-function} by $\Lambda$.
It becomes
\[
 \sqrt{
   t+\frac{1-t}{\Lambda^2}
 }\,
 \exp\left(
   -\frac{m}{200}
   \left(1-\frac1{\Lambda^2}\right)t(1-t)
 \right).
\]
This converges pointwise on $[0,1]$ to
\[
 \sqrt t\,
 \exp\left(-\frac{m}{200}t(1-t)\right).
\]
Moreover, the normalized integrand is bounded above by one for every
$\Lambda\ge2$ and $0\le t\le1$.  The dominated convergence theorem
therefore gives
\eqref{eq:dimensionless-volume-asymptotic}.

Finally, recall that
\[
 s_j
 =
 \frac{\Lambda_j^2}{2\mathfrak p_{0,j}}
 \left(1-\frac1{\Lambda_j^2}\right).
\]
Combining this identity with
\eqref{eq:protected-connector-volume} and
\eqref{eq:dimensionless-volume-asymptotic} yields
\eqref{eq:connector-volume-scale}.  Since the final metric agrees with
$g_{\mathcal C_j}$ on $\mathcal P_j$, the same formulas
remain valid there without any smoothing error.
\end{proof}

The two free parameters of the holonomy connector govern complementary geometric
scales.  The parameter \(\Lambda_j\) prescribes the expansion factor of
the circle radius, while, once \(\Lambda_j\) has been fixed,
\(\mathfrak p_{0,j}\) determines the longitudinal length.  By choosing
these parameters inductively, one can prescribe logarithmic
volume-growth exponents along suitable sequences of radii.

\begin{proposition}
\label{prop:volume-growth-flexibility}
Let \((\mathfrak e_j)_{j\ge1}\) be any sequence with
\[
 \mathfrak e_j\in\left[1,\frac32\right].
\]
The holonomy-connector parameters may be chosen so that there exist strictly
increasing sequences of radii \(r_j,\widetilde r_j\to\infty\) satisfying
\[
 \frac{\log\operatorname{Vol}_M(B_{r_j}(p))}{\log r_j}
 =\mathfrak e_j+o(1),
 \qquad
 \frac{\log\operatorname{Vol}_{\widetilde M}
 (B_{\widetilde r_j}(\widetilde p))}
 {\log\widetilde r_j}
 =\mathfrak e_j+o(1).
 \tag{8.19}\label{eq:prescribed-volume-exponents}
\]
In particular, by taking
\[
 \mathfrak e_{2j}=1,
 \qquad
 \mathfrak e_{2j+1}=\frac32,
\]
one obtains
\[
 \liminf_{r\to\infty}
 \frac{\log\operatorname{Vol}_M(B_r(p))}{\log r}=1,
 \qquad
 \limsup_{r\to\infty}
 \frac{\log\operatorname{Vol}_M(B_r(p))}{\log r}
 \ge\frac32,
 \tag{8.20}\label{eq:base-growth-oscillation}
\]
and
\[
 \liminf_{r\to\infty}
 \frac{\log\operatorname{Vol}_{\widetilde M}
 (B_r(\widetilde p))}{\log r}=1,
 \qquad
 \limsup_{r\to\infty}
 \frac{\log\operatorname{Vol}_{\widetilde M}
 (B_r(\widetilde p))}{\log r}
 \ge\frac32.
 \tag{8.21}\label{eq:cover-growth-oscillation}
\]
\end{proposition}

\begin{proof}
Before choosing the parameters of the \(j\)-th holonomy connector, let
\(\mathcal X_j\subset M\) be the compact region ending at its left
boundary, and let \(\widehat{\mathcal X}_j\) be the component of
\(q^{-1}(\mathcal X_j)\) containing
\(\widetilde p\).  At this stage,
\[
 \operatorname{Vol}(\mathcal X_j),\quad
 \operatorname{diam}(\mathcal X_j),\quad
 \operatorname{Vol}(\widehat{\mathcal X}_j),\quad
 \operatorname{diam}(\widehat{\mathcal X}_j),
 \quad \ell_j(0),\quad \rho_j(0)
\]
are already fixed.  Proposition~\ref{prop:un-connector} imposes only an
upper bound on \(\mathfrak p_{0,j}\) and a lower bound on \(\Lambda_j\).

Suppose first that \(1<\mathfrak e_j\le3/2\).  Choose an admissible
constant
\(\varepsilon_j>0\) and set
\[
 \mathfrak p_{0,j}
 =
 \varepsilon_j
 \Lambda_j^{-(3-2\mathfrak e_j)/(\mathfrak e_j-1)}.
 \tag{8.22}\label{eq:parameter-choice-q}
\]
The definitions of \(s_j\) and the volume scale
\eqref{eq:connector-volume-scale} then give the
exact identities
\[
 \begin{aligned}
 s_j
 &=
 \frac{1-\Lambda_j^{-2}}{2\varepsilon_j}\,
 \Lambda_j^{1/(\mathfrak e_j-1)},\\
 \ell_j(0)\rho_j(0)^m
 \frac{\Lambda_j^3}{\mathfrak p_{0,j}}
 &=
 \frac{\ell_j(0)\rho_j(0)^m}{\varepsilon_j}\,
 \Lambda_j^{\mathfrak e_j/(\mathfrak e_j-1)}.
 \end{aligned}
 \tag{8.23}\label{eq:parameter-scaling-q}
\]
Thus, with \(j\) fixed and \(\Lambda_j\to\infty\),
\[
 \log s_j
 =
 \frac{1}{\mathfrak e_j-1}\log\Lambda_j+O_j(1),
 \qquad
 \log\left(
 \ell_j(0)\rho_j(0)^m
 \frac{\Lambda_j^3}{\mathfrak p_{0,j}}
 \right)
 =
 \frac{\mathfrak e_j}{\mathfrak e_j-1}\log\Lambda_j+O_j(1).
\]
We now choose \(\Lambda_j\) sufficiently large that all previously fixed
volume, diameter, and boundary-scale terms contribute at most
\(\log s_j/(j+1)\) on the logarithmic scale.  We also require
\[
 \frac{
 \operatorname{diam}(\mathcal X_j)
 +\operatorname{diam}(\widehat{\mathcal X}_j)
 +2^{j+1}\bigl(\Lambda_j\ell_j(0)+\rho_j(0)\bigr)}
 {s_j}
 \le\frac1{j+1},
 \qquad
 \frac{j+1}{\log s_j}\le\frac1{j+1}.
 \tag{8.24}\label{eq:scale-separation}
\]
These requirements are compatible because
\[
 \frac{1}{\mathfrak e_j-1}\ge2,
\]
so \(s_j\) grows at least quadratically in \(\Lambda_j\).

Let
\[
 \Sigma_j
 =
 \left\{\frac{s_j}{2}\right\}
 \times\mathbb S^1\times\mathbb S^m
\]
be the midpoint slice of the \(j\)-th holonomy connector, and let \(M_j\) be the
compact region bounded by \(\Sigma_j\) that contains the initial cap.
Define
\[
 d_j^-=\min_{y\in\Sigma_j}d(p,y),
 \qquad
 d_j^+=\max_{y\in\Sigma_j}d(p,y),
 \qquad
 r_j=d_j^++\frac{s_j}{8}.
\]
The intrinsic diameter of \(\Sigma_j\) satisfies
\[
 d_j^+-d_j^-
 \le
 2\pi\bigl(\Lambda_j\ell_j(0)+\rho_j(0)\bigr).
 \tag{8.25}\label{eq:slice-distance-spread}
\]
Moreover, throughout the subcylinder
$\mathcal C_j[1/2,3/4]\subset\mathcal P_j$, the $s$-coordinate is a unit-speed distance
coordinate.  Hence every point \(y\) on the slice at parameter \(s\)
satisfies
\[
 d_j^-+s-\frac{s_j}{2}
 \le d(p,y)
 \le d_j^++s-\frac{s_j}{2}.
 \tag{8.26}\label{eq:connector-distance-comparison}
\]
It follows that
\[
 \left\{
 \frac{s_j}{2}<s<\frac{5s_j}{8}
 \right\}
 \subset B_{r_j}(p)
 \subset
 M_j\cup
 \left\{
 \frac{s_j}{2}<s<
 \left(
 \frac58+\frac{d_j^+-d_j^-}{s_j}
 \right)s_j
 \right\}.
 \tag{8.27}\label{eq:ball-connector-sandwich}
\]
By \eqref{eq:scale-separation}, the right-hand cutoff lies in
$\mathcal P_j$ for all sufficiently large \(j\).

Since
\[
 \frac{s_j}{4}
 \le d_j^-
 \le \operatorname{diam}(\mathcal X_j)+\frac{s_j}{2},
\]
equations \eqref{eq:scale-separation} and
\eqref{eq:slice-distance-spread} imply
\[
 \log r_j=\log s_j+o(\log s_j).
\]
Lemma~\ref{lem:protected-connector-volume}, applied to the subcylinder
on the left-hand side of
\eqref{eq:ball-connector-sandwich}, gives the corresponding lower volume
bound.  The volume of the entire unsmoothed holonomy connector is bounded above
by a uniform multiple of
\[
 \ell_j(0)\rho_j(0)^m
 \frac{\Lambda_j^3}{\mathfrak p_{0,j}}.
\]
The smoothing regions may be chosen to contribute less than this
quantity.  Consequently, the fixed volume of \(\mathcal X_j\), the full
holonomy-connector volume, and the two subcylinders in
\eqref{eq:ball-connector-sandwich} yield
\[
 \log\operatorname{Vol}_M(B_{r_j}(p))
 =
 \mathfrak e_j\log s_j+o(\log s_j).
\]
Since \(\log r_j=\log s_j+o(\log s_j)\), we obtain
\[
 \frac{\log\operatorname{Vol}_M(B_{r_j}(p))}
 {\log r_j}
 =\mathfrak e_j+o(1).
\]

We next treat the endpoint \(\mathfrak e_j=1\).  Fix any admissible value
\(\Lambda_j\ge2\) and let \(\mathfrak p_{0,j}\downarrow0\).  Since
\(\Lambda_j\) is fixed, we have the exact identity
\[
 s_j
 =
 \frac{\Lambda_j^2-1}{2}\,
 \mathfrak p_{0,j}^{-1}.
\]
Moreover, for every fixed $1/4\le a<b\le3/4$,
Lemma~\ref{lem:protected-connector-volume} gives
\[
 \operatorname{Vol}\bigl(
 \{a s_j\le s\le b s_j\}\subset\mathcal C_j
 \bigr)
 =
 C_j(a,b)\,\mathfrak p_{0,j}^{-1},
\]
where
\[
 \begin{aligned}
 C_j(a,b)
 &:=
 \pi\operatorname{Vol}(\mathbb S^m)\,
 \ell_j(0)\rho_j(0)^m(\Lambda_j^2-1)\cdot
 \bigl(
 \mathcal I_m(\Lambda_j,b)
 -
 \mathcal I_m(\Lambda_j,a)
 \bigr)>0
\end{aligned}
\]
is independent of \(\mathfrak p_{0,j}\).  Consequently,
\[
 \log s_j
 =
 -\log\mathfrak p_{0,j}+O_j(1),
 \qquad
 \log\operatorname{Vol}\bigl(
 \{a s_j\le s\le b s_j\}
 \bigr)
 =
 -\log\mathfrak p_{0,j}+O_j(1).
\]
After choosing \(\mathfrak p_{0,j}\) sufficiently small, the same
scale-separation, distance, and volume comparisons therefore yield
\[
 \frac{\log\operatorname{Vol}_M(B_{r_j}(p))}
 {\log r_j}
 =1+o(1).
\]


It remains to verify the corresponding estimates on the universal
cover.  The component of the preimage of the \(j\)-th holonomy connector
adjacent to \(\widehat{\mathcal X}_j\) is the \(2^{j+1}\)-fold cyclic
cover in the circle direction.  Let \(\widetilde\Sigma_j\) be the lift
of the midpoint
slice bounding the component that contains \(\widetilde p\), and define
\[
 \widetilde d_j^-=
 \min_{y\in\widetilde\Sigma_j}d(\widetilde p,y),
 \qquad
 \widetilde d_j^+=
 \max_{y\in\widetilde\Sigma_j}d(\widetilde p,y),
 \qquad
 \widetilde r_j=
 \widetilde d_j^++\frac{s_j}{8}.
\]
The protected volume is multiplied by \(2^{j+1}\), and the analogue of
\eqref{eq:slice-distance-spread} acquires the same factor in its circle
term.  The smoothing-volume estimates lift with the same degree.
Because
\[
 \log 2^{j+1}=(j+1)\log2=o(\log s_j)
\]
by \eqref{eq:scale-separation}, this finite covering degree is
negligible on the logarithmic scale.  Repeating the preceding ball
comparison therefore gives
\[
 \frac{\log\operatorname{Vol}_{\widetilde M}
 (B_{\widetilde r_j}(\widetilde p))}
 {\log\widetilde r_j}
 =\mathfrak e_j+o(1).
\]
By increasing the successive scales if necessary, both
\((r_j)\) and \((\widetilde r_j)\) may be taken to be strictly
increasing.  The alternating choice
\(\mathfrak e_{2j}=1\) and \(\mathfrak e_{2j+1}=3/2\) now yields
\eqref{eq:base-growth-oscillation} and
\eqref{eq:cover-growth-oscillation}.
\end{proof}

We conclude with two consequences for the asymptotic volume ratio.  For
a complete \(n\)-manifold \(X\) with nonnegative Ricci curvature, set
\[
 \operatorname{AVR}(X)
 :=
 \lim_{r\to\infty}
 \frac{\operatorname{Vol}(B_r(x))}
 {\operatorname{Vol}(B_r^{\mathbb R^n})}.
\]
By the Bishop--Gromov theorem, this limit exists and is independent of
the base point \(x\) \cite[Chapter~9]{Petersen2016}.

\begin{proposition}
\label{prop:volume-consequences}
For every admissible choice of the holonomy-connector parameters,
\[
 \operatorname{AVR}(M)=0,
 \qquad
 \limsup_{r\to\infty}
 \frac{\operatorname{Vol}_M(B_r(p))}{r}
 =\infty.
 \tag{8.28}\label{eq:base-avr-linear}
\]
If \(m=2\), then every admissible choice also satisfies
\[
 \operatorname{AVR}(\widetilde M)=0.
 \tag{8.29}\label{eq:cover-avr-dim4}
\]
For either \(m=2\) or \(m=3\), the alternating construction in
Proposition~\ref{prop:volume-growth-flexibility} satisfies
\[
 \operatorname{AVR}(M)
 =
 \operatorname{AVR}(\widetilde M)
 =0.
 \tag{8.30}\label{eq:alternating-avr}
\]
For arbitrary choices of the holonomy-connector parameters when \(m=3\), the
value of \(\operatorname{AVR}(\widetilde M)\) is not determined by the
present argument.
\end{proposition}

\begin{proof}
A theorem of Anderson and Li states that a complete manifold with
nonnegative Ricci curvature and positive asymptotic volume ratio has
finite fundamental group \cite{AndersonTopology1990,Li1986}.  Since
\(\pi_1(M)\cong\CPrufer\) is infinite, it follows immediately that
\(\operatorname{AVR}(M)=0\).

Navarro, Pan, and Zhu proved that a complete manifold with positive
Ricci curvature and at most linear volume growth has finitely generated
fundamental group \cite{NavarroPanZhu2024}.  Since \(\pi_1(M)\) is not
finitely generated,
\(\operatorname{Vol}_M(B_r(p))\) cannot be bounded above by \(Cr\) for
any constant \(C\).  Equivalently,
\[
 \limsup_{r\to\infty}
 \frac{\operatorname{Vol}_M(B_r(p))}{r}
 =\infty.
\]

When \(m=2\), the universal cover is four-dimensional.  A theorem of
Huang and Huang asserts that if the universal cover of a complete
four-manifold with nonnegative Ricci curvature has Euclidean volume
growth, then its fundamental group is finitely generated.  This is
impossible here \cite{HuangHuang2025}, and hence
\(\operatorname{AVR}(\widetilde M)=0\).

Finally, for the alternating choice of parameters, the even subsequence
in Proposition~\ref{prop:volume-growth-flexibility} satisfies
\[
 \operatorname{Vol}_M(B_{r_{2j}}(p))
 =r_{2j}^{\,1+o(1)},
 \qquad
 \operatorname{Vol}_{\widetilde M}
 (B_{\widetilde r_{2j}}(\widetilde p))
 =\widetilde r_{2j}^{\,1+o(1)}.
\]
Since \(\dim M=m+2\ge4\), the corresponding Bishop--Gromov volume ratios
tend to zero along this subsequence
\cite[Chapter~9]{Petersen2016}.  Their monotonicity then implies
\[
 \operatorname{AVR}(M)
 =
 \operatorname{AVR}(\widetilde M)
 =0.
\]
\end{proof}


\appendix
\section{Perelman's neck construction}\label{app:perelman-neck}
In this appendix we prove Proposition~\ref{prop:unified-neck}.  The
argument follows Perelman's neck construction
\cite[Assertion and Section~3]{Perelman1997}, with the quantitative
normalization and equivariance required in
Section~\ref{sec:unified-directed-base} recorded explicitly.  The proof
has three steps.  Lemmas~\ref{lem:concave-warping}--
\ref{lem:angular-latitude-form} construct the positively curved angular
spheres and determine their boundary geometry.
Lemma~\ref{lem:angular-slice-interpolation} and
Lemma~\ref{lem:variable-angular-ricci} introduce a path of slice metrics
and compute the curvature of the associated neck.  Finally,
Lemma~\ref{lem:slow-logarithmic-neck} and
Lemma~\ref{lem:neck-endpoint-normalization} choose a sufficiently slow
interpolation, prove positivity of its Ricci tensor, and normalize the
two ends.

We begin with the auxiliary function \(\mathcal R\).  Its construction
is a quantitative version of \cite[Section~3]{Perelman1997}.

\begin{lemma}\label{lem:concave-warping}
There is $r_0>0$ such that, for every $0<r<r_0$, there is $w_0(r)>0$
with the following property.  For any $0<w<w_0(r)$, there is a smooth
function $\mathcal R:[0,5r]\to\mathbb R$ satisfying
\begin{align}
 \mathcal R(0)&=0,& \mathcal R'(0)&=1,&
 \mathcal R^{(2j)}(0)&=0\quad(j\ge1),                              \label{eq:unified-base-3-2}\\
 \mathcal R&>0,&0<\mathcal R'&<1,&\mathcal R''&<0
                         \quad\text{on }(0,5r],                    \label{eq:unified-base-3-3}\\
 -\frac{\mathcal R''}{\mathcal R}&\ge \frac2{r^2}
                         &&\text{on }(0,r^2/2],&
 -\frac{\mathcal R''}{\mathcal R}&\ge1
                         \quad\text{on }(0,5r],                    \label{eq:unified-base-3-4}\\
 \cot r\,\mathcal R(r)&=w.                                        \label{eq:unified-base-3-5}
\end{align}
Moreover,
\begin{equation}
\mathcal R(t)=\frac{w\tan r}{\sin(r+r^4/4)}
                  \sin(t+r^4/4)                                  \label{eq:unified-base-3-6}
\end{equation}
on a neighborhood of $[r,5r]$ in $[0,5r]$.
\end{lemma}

\begin{proof}
Set
\[
 a=\frac{r^2}{2},\qquad
 b=\frac{r^4}{4},\qquad
 \kappa=\frac{\sqrt2}{r}.
\]
We first construct a positive solution $u_0$ with $u_0=\sin(t+b)$ on $[a,5r]$. Consider the piecewise constant coefficient
\[
 q_0(t)=
 \begin{cases}
  \kappa^2,&0\le t\le a,\\
  1,&a<t\le5r,
 \end{cases}
\]
and let \(u_0\) be the \(C^1\) solution of
\[
 u_0''+q_0u_0=0
\]
which equals \(\sin(t+b)\) on \([a,5r]\).  We will show in the following that $u_0>0$ on $[0,5r]$. Matching its value and
derivative at \(t=a\) gives
\begin{align*}
 u_0(0)
 &=\sin(a+b)\cos(\kappa a)
   -\kappa^{-1}\cos(a+b)\sin(\kappa a)
   =\frac{r^4}{6}+O(r^6),\\
 u_0'(0)
 &=\kappa\sin(a+b)\sin(\kappa a)
   +\cos(a+b)\cos(\kappa a)
   =1+\frac{r^2}{4}+O(r^4).
\end{align*}
After decreasing \(r_0\), we have
\[
 u_0(0)>0,\qquad u_0'(0)>0,\qquad
 \kappa^2\ge1,\qquad a<r,\qquad
 \kappa a=\frac r{\sqrt2}<\frac{\pi}{2},
 \qquad 5r+b<\frac{\pi}{2}.
\]
On \([0,a]\),
\[
 u_0(t)=u_0(0)\cos(\kappa t)
        +\kappa^{-1}u_0'(0)\sin(\kappa t),
\]
which yields \(u_0>0\).  Hence \(u_0''=-\kappa^2u_0<0\), and
\(u_0'(a)=\cos(a+b)>0\), which also implies \(u_0'>0\) on \([0,a]\).
On \([a,5r]\), positivity of \(u_0\) and \(u_0'\) follows from
\(5r+b<\pi/2\).

Choose \(T\in(a,r)\) sufficiently close to \(a\) and 
\(q\in C^\infty([0,5r])\), so that
\[
 q=\kappa^2\ \text{on }[0,a],\qquad
 1\le q\le\kappa^2\ \text{on }[a,T],\qquad
 q=1\ \text{on }[T,5r].
\]
The $T$ and function $q$ can be chosen so that
\[
 \|q-q_0\|_{L^1([0,5r])}
 \le \kappa^2(T-a)
\]
is arbitrarily small.  Let \(u\) solve
\[
 u''+qu=0,\qquad
 u(T)=\sin(T+b),\qquad u'(T)=\cos(T+b).
\]
Writing the equation as a first-order system, and applying Grönwall's
inequality, we obtain
\[
 \|(u,u')-(u_0,u_0')\|_{C^0([0,T])}
 \le C(r)\|q-q_0\|_{L^1([0,T])}.
\]
Thus, after taking \(T-a\) sufficiently small, we have
\begin{equation}
 u>0,\qquad u'>0\quad\text{on }[0,5r],
 \qquad
 u(t)=\sin(t+b)\quad\text{on }[T,5r].
 \label{eq:warping-background-solution}
\end{equation}

Let \(v\) be the solution
\[
 v''+qv=0,\qquad v(0)=0,\qquad v'(0)=1.
\]
Since \(q=\kappa^2\) on \([0,a]\),
\begin{equation}
 v(t)=\kappa^{-1}\sin(\kappa t)
 \qquad(0\le t\le a).
 \label{eq:warping-pole-solution}
\end{equation}
In particular, \(v(t),v'(t)>0\) on \((0,a]\).

We next match the values of \(v\) with a small multiple of \(u\).
The Wronskian
\[
 W(t)=v'(t)u(t)-v(t)u'(t)
\]
is constant, this is because
\[
 W'=v''u-vu''=(-qv)u-v(-qu)=0.
\]
Consequently,
\[
 W(t)=W(0)=u(0)>0,\qquad
 \left(\frac vu\right)'
 =\frac{v'u-vu'}{u^2}
 =\frac{u(0)}{u^2}>0.
\]
Set
\[
 \omega_*=\frac{v(a/4)}{u(a/4)}>0,\qquad
 w_0(r)=\omega_*\frac{\sin(r+b)}{\tan r}>0.
\]
For \(0<w<w_0(r)\), denote
\[
 \omega=\frac{w\tan r}{\sin(r+b)}\in(0,\omega_*).
\]
The strict monotonicity of \(v/u\) gives a unique
\(\xi\in(0,a/4)\) satisfying
\begin{equation}
 v(\xi)=\omega u(\xi).
 \label{eq:warping-value-match}
\end{equation}
At point $\xi$, 
\begin{align}
 D
 &=v'(\xi)-\omega u'(\xi)\notag\\
 &=\frac{v'(\xi)u(\xi)-v(\xi)u'(\xi)}{u(\xi)}
 =\frac{u(0)}{u(\xi)}>0.
 \label{eq:warping-derivative-drop}
\end{align}

Define the continuous piecewise-smooth function
\[
 \mathcal R_{\mathrm{pt}}(t):=
 \begin{cases}
  v(t),&0\le t\le\xi,\\
  \omega u(t),&\xi\le t\le5r.
 \end{cases}
\]
By definition, it is positive function on \((0,5r]\), and both one-sided derivatives are
positive.  Its derivative at $\xi$ has the jump
\[
 \mathcal R_{\mathrm{pt}}'(\xi+)
 -\mathcal R_{\mathrm{pt}}'(\xi-)=-D.
\]
Since \(\mathcal R_{\mathrm{pt}}\) is continuous, its distributional
first derivative has no Dirac mass.  Hence, on the constant-coefficient
interval \((0,a)\),
\begin{equation}
 \left(\frac{\dd^2}{\dd t^2}+\kappa^2\right)
 \mathcal R_{\mathrm{pt}}
 =-D\,\delta_\xi,
 \label{eq:warping-corner-distribution}
\end{equation}
in the distributional sense.
Choose an even function
\[
 \varphi\in C_c^\infty((-1,1)),\qquad
 \varphi\ge0,\qquad \int_{\mathbb R}\varphi\,\dd s=1,
\]
and let \(\varphi_h(s)=h^{-1}\varphi(s/h)\).  Choose \(h>0\) small enough so
that
\[
 [\xi-3h,\xi+3h]\subset(0,a),
 \qquad \kappa h<\frac{\pi}{2}.
\]
Define
\[
 c_h=\int_{\mathbb R}
       \varphi_h(s)\cos(\kappa s)\,\dd s>0.
\]
If \(y''+\kappa^2y=0\), then
\[
 y(t-s)=y(t)\cos(\kappa s)
        -\kappa^{-1}y'(t)\sin(\kappa s),
\]
and the evenness of \(\varphi_h\) gives
\begin{equation}
 \int_{\mathbb R}\varphi_h(s)y(t-s)\,\dd s=c_hy(t).
 \label{eq:warping-convolution-eigenfunction}
\end{equation}

On the open cover
\[
 \{|t-\xi|<2h\}\cup\{|t-\xi|>h\}
\]
of \([0,5r]\), define
\[
 \mathcal R(t)=
 \begin{cases}
 c_h^{-1}\displaystyle\int_{\mathbb R}
       \varphi_h(s)\mathcal R_{\mathrm{pt}}(t-s)\,\dd s,
       &|t-\xi|<2h,\\[2mm]
 \mathcal R_{\mathrm{pt}}(t),&|t-\xi|>h.
 \end{cases}
\]
If \(|t-\xi|<2h\) and \(|s|<h\), then
\[
 t-s\in(\xi-3h,\xi+3h)\subset(0,a).
\]
On the overlap annulus
\(h<|t-\xi|<2h\), \eqref{eq:warping-convolution-eigenfunction} shows that the two
definitions agree.  Therefore \(\mathcal R\) is smooth.
\eqref{eq:warping-corner-distribution} then yields, on $ \{|t-\xi|<2h\}$,
\begin{equation}
 \mathcal R''+\kappa^2\mathcal R
 =-\frac{D}{c_h}\varphi_h(t-\xi)\le0.
 \label{eq:warping-smoothed-corner}
\end{equation}
Moreover, since \(\mathcal R_{\mathrm{pt}}\) is continuous, its first
distributional derivative is the positive piecewise
function
\[
 \mathcal R_{\mathrm{pt}}'
 =
 \begin{cases}
  v',&t<\xi,\\
  \omega u',&t>\xi.
 \end{cases}
\]
The nonnegativity of \(\varphi_h\) therefore gives
\[
 \mathcal R>0,\qquad \mathcal R'>0
 \quad\text{ on } \{|t-\xi|<2h\}.
\]
The same positivity statements hold when
$\{|t-\xi|\geq 2h\}$ by the definition of $\mathcal R$.

On $ \{|t-\xi|\geq 2h\}$, 
\[
 -\frac{\mathcal R''}{\mathcal R}=q.
\]
In $ \{|t-\xi|<2h\}$, \eqref{eq:warping-smoothed-corner} gives
\[
 -\frac{\mathcal R''}{\mathcal R}
 =\kappa^2+
   \frac{D\,\varphi_h(t-\xi)}{c_h\mathcal R(t)}
 \ge\kappa^2=\frac2{r^2}.
\]
Consequently,
\[
 -\frac{\mathcal R''}{\mathcal R}\ge\frac2{r^2}
 \quad\text{on }(0,r^2/2],
 \qquad
 -\frac{\mathcal R''}{\mathcal R}\ge1
 \quad\text{on }(0,5r].
\]
Since \(\mathcal R>0\), the above inequalities imply
\(\mathcal R''<0\).  Hence \(\mathcal R'\) is strictly decreasing, and
\(\mathcal R'(0)=1\) gives
\[
 0<\mathcal R'(t)<1,\qquad0<t\le5r.
\]

Since the smoothing procedure is supported away from \(0\),
\eqref{eq:warping-pole-solution} gives
\[
 \mathcal R(t)=\frac r{\sqrt2}
 \sin\left(\frac{\sqrt2}{r}t\right)
\]
near \(t=0\), which proves all the identities in
\eqref{eq:unified-base-3-2}.  It is also supported to the left of
\(T<r\).  By \eqref{eq:warping-background-solution},
\[
 \mathcal R(t)
 =\omega\sin(t+b)
 =\frac{w\tan r}{\sin(r+r^4/4)}
   \sin(t+r^4/4)
\]
on \([T,5r]\), proving \eqref{eq:unified-base-3-6}.  Evaluating at
\(t=r\) gives
\[
 \cot r\,\mathcal R(r)=w,
\]
which is \eqref{eq:unified-base-3-5}.  This completes the proof.
\end{proof}
Next, we prove the ambient metric $g_{\mathrm{amb},m}$ is smooth and has positive sectional curvature near
\(B_{4r}^{\angle}(o_0)\).
\begin{lemma}
\label{lem:neck-ambient-curvature}
There exists a constant \(r_{\mathrm{amb}}>0\) with the following
property.  Let \(m\in\{2,3\}\), \(k=m-1\), and
\(0<r<r_{\mathrm{amb}}\), and let \(\mathcal R\) be the function
constructed in Lemma~\ref{lem:concave-warping}.  Then the metric
\eqref{eq:unified-base-3-11} is smooth and has positive sectional curvature near
\(B_{4r}^{\angle}(o_0)\).
\end{lemma}

\begin{proof}
Choose $r_{\mathrm{amb}}>0$ small enough so that $r_{\mathrm{amb}}<\min\{r_0, \pi/10\}$ for $r_0$ defined in Lemma~\ref{lem:concave-warping}, and whenever
 $0<r<r_{\mathrm{amb}}$,
\begin{equation}
\cot r\tan(r^2/2)<\frac1{\sqrt2},\qquad
 \frac{\tan(t+r^4/4)}{\tan t}<1+\tan^2r
 \quad\text{for }r^2/2\le t\le r.                                  \label{eq:unified-base-3-14}
\end{equation}
Such a choice is possible: \(\cot r\tan(r^2/2)\) tends to zero as
\(r\to0\), while, uniformly for \(r^2/2\le t\le r\),
\[
\frac{\tan(t+r^4/4)}{\tan t}
 =1+\frac{r^4}{4}\frac{\sec^2t}{\tan t}
  +O\!\left(\frac{r^8}{t^2}\right)
 \le1+\frac{r^2}{2}+O(r^4)
 \le1+\tan^2r
\]
for all sufficiently small \(r\).  Let \(\mathcal R\) be the auxiliary function constructed in
Lemma~\ref{lem:concave-warping}, and consider the unscaled
metric
\[
 \widehat g_{\mathrm{amb},m}
 =\dd t^2+\cos^2t\,\dd z^2+\mathcal R(t)^2g_{\mathbb S^k}.
\]
For orthonormal vertical vectors $V,W\in T\mathbb S^k$, the possible
sectional curvatures are
\begin{align}
 \sec_{\widehat g_{\mathrm{amb},m}}
   \bigl(\partial_t,(\cos t)^{-1}\partial_z\bigr)&=1,             \label{eq:unified-base-3-15}\\
 \sec_{\widehat g_{\mathrm{amb},m}}(\partial_t,V)
   &=-\frac{\mathcal R''}{\mathcal R},                            \label{eq:unified-base-3-16}\\
 \sec_{\widehat g_{\mathrm{amb},m}}
   \bigl((\cos t)^{-1}\partial_z,V\bigr)
 &=\frac{\mathcal R'}{\mathcal R}\tan t,                         \label{eq:unified-base-3-17}\\
 \sec_{\widehat g_{\mathrm{amb},m}}(V,W)
   &=\frac{1-(\mathcal R')^2}{\mathcal R^2}.                     \label{eq:unified-base-3-18}
\end{align}
The \eqref{eq:unified-base-3-18} occurs only when $m=3$.  
\eqref{eq:unified-base-3-3}--\eqref{eq:unified-base-3-4} proved in Lemma \ref{lem:concave-warping} make all these sectional curvatures positive, and
the conditions in \eqref{eq:unified-base-3-2} are exactly the smoothness conditions
for collapsing the unit-round $\mathbb S^k$ factor at $t=0$.
Since
\[
 \varrho_0(t,z)=\arccos(\cos t\cos z)\geq t,
\]
one has
\[
 B_{4r}^{\angle}(o_0)\subset\{0\leq t\leq4r\}
 \Subset\{0\leq t<5r\}.
\]
Hence, the metric is defined and smooth on a neighborhood of $B_{4r}^{\angle}(o_0)$.
Finally, scaling the metric by the positive constant \(\cot^2r\)
divides all sectional curvatures by \(\cot^2r\), without changing their
signs.  This proves the lemma.
\end{proof}
The next lemma gives the required sectional-curvature lower bound on
the angular spheres and computes their principal curvatures.
\begin{lemma}
\label{lem:angular-sphere-geometry}
Under the hypotheses of Lemma~\ref{lem:neck-ambient-curvature}, for
every \(z_0\in\mathbb R\), the hypersurface
\(S_r^{\angle}(o_{z_0})\) is smooth. The induced metric
\(g_{\mathrm{ang},m}\) satisfies
\[
                         \sec_{g_{\mathrm{ang},m}}>1 .
\]
All principal curvatures for the normal pointing out of
\(B_r^{\angle}(o_{z_0})\) belong to \((0,1]\), and the principal curvature of the meridional curve, i.e. the curve $(t, z, \theta_0)$ in $S_r^{\angle}(o_{z_0})$ with fixed $\theta_0\in \mathbb S^k$, is exactly \(1\).
\end{lemma}

\begin{proof}
Fix $z_0\in\mathbb R$ and consider $S_r^{\angle}(o_{z_0})$.
The projection to the
round metric $\dd t^2+\cos^2t\,\dd z^2$ is length nonincreasing, and its
radial geodesics lift with constant $\mathbb S^k$-coordinate.  Hence, on
$|z-z_0|<\pi$, $\varrho_{z_0}$ is also the distance from $o_{z_0}$ for
$\widehat g_{\mathrm{amb},m}$.  Since the angular sphere $S_r^{\angle}(o_{z_0})$ is the level set
$\{\varrho_{z_0}=r\}$, its outward unit normal is
$\nu=\nabla\varrho_{z_0}$.  Fix $z$, differentiate the identity
$\cos\varrho_{z_0}(t,z)=\cos t\cos(z-z_0)$ with respect to the ambient
coordinate $t$. This gives
\[
 -\sin\varrho_{z_0}\,\partial_t\varrho_{z_0}
   =-\sin t\cos(z-z_0).
\]
On $S_r^{\angle}(o_{z_0})$ we have $\varrho_{z_0}=r$ and
$\cos(z-z_0)=\cos r/\cos t$.  Because $\partial_t$ is a unit vector,
$\partial_t\varrho_{z_0}
=\langle\nabla\varrho_{z_0},\partial_t\rangle$; hence
\begin{equation}
\langle\nu,\partial_t\rangle=\cot r\tan t.                       \label{eq:unified-base-3-19}
\end{equation}
We now specify the tangent vectors used below.  Away from the two poles
of the angular sphere, let $E$ be a unit vector spanning
\[
 TS_r^{\angle}(o_{z_0})\cap\vspan\{\partial_t,\partial_z\},
\]
and let $V,W$ be orthonormal vectors tangent to the $\mathbb S^k$
factor, with orthonormality measured using
$\widehat g_{\mathrm{amb},m}$.  Since $\varrho_{z_0}$ is independent of
the $\mathbb S^k$-coordinate, these vertical vectors are also tangent to
the angular sphere.  Thus the orthogonal splitting is
\[
 TS_r^{\angle}(o_{z_0})
   =\vspan\{E\}\oplus T\mathbb S^k.
\]
In particular, $E,V,W$ are all unit tangent vectors to the same angular
sphere: $E$ is the meridional direction, whereas $V,W$ are vertical
directions.  The vector $W$ is needed only when $k=2$, equivalently
$m=3$.  The formulas at the two poles are understood by smooth
extension.
The Gauss equation gives the sectional curvatures of the metric
induced by $\widehat g_{\mathrm{amb},m}$:
\begin{align}
 \sec_{\widehat g_{\mathrm{amb},m}|_{S_r^{\angle}(o_{z_0})}}(E,V)
 &=-\frac{\mathcal R''}{\mathcal R}
       \bigl(1-\cot^2r\tan^2t\bigr)
   +\frac{\mathcal R'}{\mathcal R}\tan t\,\cot^2r\tan^2t
   +\cot^2r\frac{\mathcal R'}{\mathcal R}\tan t,                 \label{eq:unified-base-3-20}\\
 \sec_{\widehat g_{\mathrm{amb},m}|_{S_r^{\angle}(o_{z_0})}}(V,W)
 &=\mathcal R^{-2}\left(
      1-(\mathcal R')^2\bigl(1-\cot^2r\tan^2t\bigr)\right),       \label{eq:unified-base-3-21}
\end{align}
where \eqref{eq:unified-base-3-21} occurs only for $m=3$.  Let
\(q=-\mathcal R''/\mathcal R\geq1\).  Then
\[
 \bigl(\mathcal R'\sin t-\mathcal R\cos t\bigr)'
   =(1-q)\mathcal R\sin t\leq0 .
\]
This expression tends to zero as \(t\downarrow0\).  Hence
\[
 \left(\frac{\mathcal R}{\sin t}\right)'
 =\frac{\mathcal R'\sin t-\mathcal R\cos t}{\sin^2t}\leq0,
\]
which gives
\begin{equation}
\mathcal R\le\sin t,\qquad
 \frac{\mathcal R'}{\mathcal R}\le\cot t,\qquad
 \mathcal R'\le\cos t.                                          \label{eq:unified-base-3-22}
\end{equation}
For $t\le r^2/2$, equation \eqref{eq:unified-base-3-19} and the first inequality in
\eqref{eq:unified-base-3-14} give $\langle\nu,\partial_t\rangle<1/\sqrt2$.
Therefore
\[
 -\frac{\mathcal R''}{\mathcal R}
       \bigl(1-\langle\nu,\partial_t\rangle^2\bigr)
 \ge\frac1{r^2}>\cot^2r,
\]
so the first term of \eqref{eq:unified-base-3-20} already has the required lower
bound.  On \([r^2/2,r]\), set
\[
 h(t):=\frac{\mathcal R'(t)}{\mathcal R(t)}
             -\cot(t+r^4/4).
\]
Then
\[
 h'(t)+\left(
 \frac{\mathcal R'(t)}{\mathcal R(t)}
             +\cot(t+r^4/4)\right)h(t)
 =\frac{\mathcal R''(t)}{\mathcal R(t)}+1\leq0.
\]
Equation \eqref{eq:unified-base-3-6} gives \(h(r)=0\).
After multiplication by the positive integrating factor, integration
from \(t\) to \(r\) gives \(h(t)\geq0\).  Therefore
\[
 \frac{\mathcal R'}{\mathcal R}\ge\cot(t+r^4/4).
\]
Moreover,
$-\mathcal R''/\mathcal R\ge1\ge
(\mathcal R'/\mathcal R)\tan t$ by \eqref{eq:unified-base-3-22}.  Substitution in
\eqref{eq:unified-base-3-20} now gives, for $r^2/2\le t\le r$,
\begin{equation}
\sec_{\widehat g_{\mathrm{amb},m}|_{S_r^{\angle}(o_{z_0})}}(E,V)
 \ge\cot(t+r^4/4)\tan t(1+\cot^2r)>\cot^2r,                      \label{eq:unified-base-3-23}
\end{equation}
where the strict inequality is the second inequality in
\eqref{eq:unified-base-3-14}.
When $m=3$, \eqref{eq:unified-base-3-19} and \eqref{eq:unified-base-3-22} also give
\begin{equation}
\sec_{\widehat g_{\mathrm{amb},m}|_{S_r^{\angle}(o_{z_0})}}(V,W)
 \ge\frac1{\sin^2t}
      -\cot^2t\bigl(1-\cot^2r\tan^2t\bigr)
 =1+\cot^2r>\cot^2r.                                              \label{eq:unified-base-3-24}
\end{equation}
Scaling by $\cot^2r$ divides these curvatures by $\cot^2r$, proving
$\sec_{g_{\mathrm{ang},m}}>1$.

We next compute the principal curvatures.  Denote
\[
 e_z=(\cos t)^{-1}\partial_z.
\]
Differentiating the defining equation of $\varrho_{z_0}$ with respect
to $z$, with $t$ fixed, and using \eqref{eq:unified-base-3-19} gives,
on $S_r^{\angle}(o_{z_0})$,
\[
 \nu=(\cot r\tan t)\partial_t
      +\frac{\sin(z-z_0)}{\sin r}\,e_z.
\]
The two coefficients have squared sum one.  Thus, up to sign, $E$ has the expression
\[
 E=\frac{\sin(z-z_0)}{\sin r}\,\partial_t
      -(\cot r\tan t)e_z.
\]
The metric $\dd t^2+\cos^2t\,\dd z^2$ has constant curvature one.  Its radial
distance function therefore satisfies
\[
 \Hess\varrho_{z_0}
   =\cot\varrho_{z_0}
      \bigl(\dd t^2+\cos^2t\,\dd z^2
            -\dd\varrho_{z_0}^{2}\bigr).
\]
Since $\dd\varrho_{z_0}(E)=0$ and $|E|=1$, the second fundamental form of
$E$ direction is
\[
 \II_{\widehat g_{\mathrm{amb},m}}(E,E)
   =\langle\nabla_E\nu,E\rangle
   =\Hess\varrho_{z_0}(E,E)=\cot r.
\]
For a unit vertical vector $V\in\mathbb S^k$, the warped-product connection formula
and the fact that $\mathcal R$ depends only on $t$ give
\[
 \nabla_V\nu=\nu(\log\mathcal R)V,
 \qquad
 \II_{\widehat g_{\mathrm{amb},m}}(V,V)
   =\nu(\log\mathcal R)
   =\frac{\mathcal R'}{\mathcal R}
      \langle\nu,\partial_t\rangle
   =\frac{\mathcal R'}{\mathcal R}\cot r\tan t.
\]
More generally, for vertical vectors $V,W$,
\[
 \II_{\widehat g_{\mathrm{amb},m}}(V,W)
   =\nu(\log\mathcal R)
      \widehat g_{\mathrm{amb},m}(V,W),
 \qquad
 \II_{\widehat g_{\mathrm{amb},m}}(E,V)=0.
\]
Thus the splitting
$TS_r^{\angle}(o_{z_0})=\vspan\{E\}\oplus T\mathbb S^k$
diagonalizes the second fundamental form.  The principal
curvature for $E$ direction is $\cot r$, while the vertical principal curvature is
$(\mathcal R'/\mathcal R)\cot r\tan t$ with multiplicity $k=m-1$.

Finally, set $g_{\mathrm{amb},m}=\cot^2r\,
\widehat g_{\mathrm{amb},m}$.  Under a constant rescaling
$g\mapsto c^2g$, the Levi--Civita connection is unchanged, while unit
vectors are divided by $c$; consequently every principal curvature is
divided by $c$.  Here $c=\cot r$.  The principal curvatures with respect
to $g_{\mathrm{amb},m}$ are therefore
\begin{equation}
\tan r\,\cot r=1,
 \qquad
 \tan r\,\frac{\mathcal R'}{\mathcal R}\cot r\tan t
   =\frac{\mathcal R'}{\mathcal R}\tan t\le1.                    \label{eq:unified-base-3-25}
\end{equation}
The vertical value is positive by \eqref{eq:unified-base-3-3}; at the
two collapsed poles it is understood by its positive smooth limit.
This proves the principal-curvature bounds. The lemma then follows.
\end{proof}


Next, we prove the induced metric $g_{\mathrm{ang},m}$ on \(S_r^{\angle}(o_{z_0})\) has the form \eqref{eq:unified-base-3-12} under latitude coordinate.
\begin{lemma}
\label{lem:angular-latitude-form}
Under the hypotheses of
Lemma~\ref{lem:angular-sphere-geometry},
\(S_r^{\angle}(o_{z_0})\cong\mathbb S^m\).  The relation
\begin{equation}\label{eq:latitude-form}
\cos\varphi=\frac{\mathcal R(t)}{\mathcal R(r)},\qquad
 \sgn\varphi=\sgn(z-z_0)
\end{equation}
define a global latitude coordinate
\(\varphi\in[-\pi/2,\pi/2]\).  In this coordinate, the induced metric $g_{\mathrm{ang},m}$
has the form \eqref{eq:unified-base-3-12}, where
\[
 A_{\mathrm{ang}}\in
 C^\infty([-\pi/2,\pi/2],(0,\infty)),\qquad
 A_{\mathrm{ang}}(\pm\pi/2)=w,
\]
\[
 \frac1\pi\int_{-\pi/2}^{\pi/2}
 A_{\mathrm{ang}}(\varphi)\,\dd\varphi=\cos r .
\]
The function \(A_{\mathrm{ang}}\) is independent of \(z_0\), and at
both poles
\begin{equation}\label{eq:pole expansion}
\frac{A_{\mathrm{ang}}(\varphi)}w
      =1+O(\cos^2\varphi).
\end{equation}

\end{lemma}

\begin{proof}
 Since
 $\mathcal R'>0$ on $[0,r]$, the relation \eqref{eq:latitude-form}
 determines
 \begin{equation}\label{eq:def of t varphi}
    t(\varphi)=\mathcal R^{-1}
       \bigl(\mathcal R(r)\cos\varphi\bigr),
  \qquad -\frac\pi2\le\varphi\le\frac\pi2.
 \end{equation}
 The defining identity
 \[
  \cos r
  =\cos t(\varphi)\cos\bigl(z(\varphi)-z_0\bigr)
 \]
 then determines its \(z\)-coordinate by
 \begin{equation}\label{eq:def of z varphi}
     z(0)=z_0,
  \qquad
  z(\varphi)=z_0+\sgn(\varphi)
       \arccos\!\left(\frac{\cos r}{\cos t(\varphi)}\right)
       \quad (\varphi\ne0).
 \end{equation}
 Thus $(\varphi,p)\mapsto(t(\varphi),z(\varphi),p)$ parametrizes
 $S_r^{\angle}(o_{z_0})$, with the $\mathbb S^k$ factor collapsed at
 $\varphi=\pm\pi/2$.  Pulling back the ambient metric
 \eqref{eq:unified-base-3-11} gives
 \begin{equation}\label{eq:def of A ang}
    A_{\mathrm{ang}}(\varphi)
  :=\cot r\sqrt{t'(\varphi)^2
                +\cos^2t(\varphi)\,z'(\varphi)^2},
 \end{equation}
 and, by \eqref{eq:unified-base-3-5},
 \[
  \cot^2r\,\mathcal R(t(\varphi))^2g_{\mathbb S^k}
  =\cot^2r\,\mathcal R(r)^2\cos^2\varphi\,g_{\mathbb S^k}
  =w^2\cos^2\varphi\,g_{\mathbb S^k}.
 \]
 This proves the metric formula \eqref{eq:unified-base-3-12} and also
 shows directly that $A_{\mathrm{ang}}$ is independent of $z_0$.

 Differentiating \eqref{eq:def of t varphi} and
 \eqref{eq:def of z varphi} gives, when
 $0<|\varphi|<\pi/2$,
 \[
  t'(\varphi)
    =-\frac{\mathcal R(r)\sin\varphi}
            {\mathcal R'(t(\varphi))},
  \qquad
  z'(\varphi)
    =-\tan t(\varphi)\cot\bigl(z(\varphi)-z_0\bigr)t'(\varphi).
 \]
 Consequently,
 \[
  A_{\mathrm{ang}}(\varphi)
   =\cot r\,
     \frac{\mathcal R(r)|\sin\varphi|\cos t(\varphi)\sin r}
     {\mathcal R'(t(\varphi))
      \sqrt{\cos^2t(\varphi)-\cos^2r}}.
 \]
 The apparent singularity at $\varphi=0$ is removable.  Indeed, \eqref{eq:def of t varphi} and \eqref{eq:def of z varphi} give
 \[
  t(\varphi)
   =r-\frac{\mathcal R(r)}{2\mathcal R'(r)}\varphi^2+O(\varphi^4),
  \qquad
  z(\varphi)-z_0
   =\sqrt{\frac{\mathcal R(r)\tan r}{\mathcal R'(r)}}\,\varphi
      +O(\varphi^3).
 \]
 Hence $A_{\mathrm{ang}}$, defined by \eqref{eq:def of A ang}, is
 smooth and nondegenerate at $\varphi=0$.  At either pole,
 $\mathcal R(t)=t+O(t^3)$ gives
 $|t'(\varphi)|\longrightarrow\mathcal R(r)$, while the same defining
 identity gives $z'(\varphi)\longrightarrow0$.  Hence
 \[
  A_{\mathrm{ang}}(\pm\pi/2)=\cot r\,\mathcal R(r)=w.
 \]
 Smoothness of the rotationally symmetric induced metric at either
 collapsed pole gives the even expansion
 \[
   \frac{A_{\mathrm{ang}}(\varphi)}w
        =1+O(\cos^2\varphi).
 \]
 Finally, the meridional curve projects to a semicircle of radius
 $\sin r$ in the unit round two-sphere.  Its length is therefore
 $\pi\sin r$ before scaling and $\pi\cos r$ after scaling by
 $\cot^2r$.  Thus
\begin{equation}
\frac1\pi\int_{-\pi/2}^{\pi/2}A_{\mathrm{ang}}(\varphi)\,\dd\varphi=\cos r.    \label{eq:unified-base-3-26}
\end{equation}
\end{proof}
We next interpolate between $g_{\mathrm{ang},m}$ and a round metric
while retaining the required sectional-curvature lower bound.
\begin{lemma}\label{lem:angular-slice-interpolation}
Let \(m\in\{2,3\}\), \(k=m-1\), and suppose that
\[
 0<\rho<1,\qquad \pi\rho<r,\qquad
 w^{m-1}<\rho^m,
\]
and \(g_{\mathrm{ang},m}\) is the angular spherical metric in
Lemmas~\ref{lem:angular-sphere-geometry}
and~\ref{lem:angular-latitude-form}.  Define \(u_*\), \(\eta(\varphi)\),
\(\widetilde g_{u,v}\), $\alpha_0$, and \(v_0(u)\) by
\begin{equation}
u_*=\max_\varphi\frac{A_{\mathrm{ang}}(\varphi)}w,
 \qquad
 \eta(\varphi)
   =\frac{A_{\mathrm{ang}}(\varphi)/w-1}{u_*-1}.                  \label{eq:unified-base-3-27}
\end{equation}
\begin{equation}
\widetilde g_{u,v}
  =v^2\left(
       \bigl(1+(u-1)\eta(\varphi)\bigr)^2\dd\varphi^2
       +\cos^2\varphi\,g_{\mathbb S^k}\right),               \label{eq:unified-base-3-28}
\end{equation}
for $1\le u\le u_*$ and $v>0$, and
\begin{equation}
\alpha_0=\frac{\log u_*}{\log\rho-\log w},
 \qquad
 v_0(u)=\rho u^{-1/\alpha_0},\qquad 1\le u\le u_*.              \label{eq:unified-base-3-32}
\end{equation}

Then,
\[
 u_*>\frac{\rho}{w}>1,\qquad u_*w<1,\qquad
 1<\alpha_0<m,
\]
\[
 \widetilde g_{1,v_0(1)}=\rho^2g_{\mathbb S^m},\qquad
 \widetilde g_{u_*,v_0(u_*)}=g_{\mathrm{ang},m},
\]
and
\[
 \sec_{\widetilde g_{u,v_0(u)}}>1
 \qquad \text{ for } 1\leq u\leq u_*.
\]
\end{lemma}

\begin{proof}
By \eqref{eq:unified-base-3-26},
$u_*\ge \cos r/w>\rho/w>1$.  Moreover,
$A_{\mathrm{ang}}>0$ implies
$-1/(u_*-1)<\eta\le1$.  Consequently,
$1+(u-1)\eta(\varphi)>0$ whenever $1\le u\le u_*$.
The pole expansion \eqref{eq:pole expansion} in
Lemma~\ref{lem:angular-latitude-form} gives
\begin{equation}\label{eq:expansion of eta}
 \eta=O(\cos^2\varphi),\qquad
 \eta'=O(\cos\varphi),\qquad
 \eta'\tan\varphi=O(1),\qquad
 \eta\tan\varphi=O(\cos\varphi).
\end{equation}



Thus $\widetilde g_{1,v}=v^2g_{\mathbb S^m}$, whereas
$\widetilde g_{u_*,w}=g_{\mathrm{ang},m}$.  Denote
\[
 f_u(\varphi)=1+(u-1)\eta(\varphi).
\]
The preceding pole expansions \eqref{eq:pole expansion} and \eqref{eq:expansion of eta} show that
\(\widetilde g_{u,v}\) is a smooth metric on \(\mathbb S^m\), for
every \((u,v)\in[1,u_*]\times(0,\infty)\), and depending
smoothly in \((u,v)\).
Let $E=(vf_u)^{-1}\partial_\varphi$, and $V,W$ be orthonormal
vectors tangent to the $\mathbb S^k$ factor.  The warped-product
curvature formulas give
\begin{align}
 \sec_{\widetilde g_{u,v}}(E,V)
 &=\frac1{v^2}\left(
      \frac1{f_u^2}
      -\frac{(u-1)\eta'\tan\varphi}{f_u^3}\right),                \label{eq:unified-base-3-29}\\
 \sec_{\widetilde g_{u,v}}(V,W)
 &=\frac1{v^2}\left(
      \frac1{\cos^2\varphi}
      -\frac{\tan^2\varphi}{f_u^2}\right),                       \label{eq:unified-base-3-30}
\end{align}
where \eqref{eq:unified-base-3-30} occurs only for $m=3$.  At the
poles, the expressions \eqref{eq:unified-base-3-29} and \eqref{eq:unified-base-3-30} are understood by their smooth limits.

If $\eta(\varphi)=1$, then
$A_{\mathrm{ang}}(\varphi)$ is maximal.  Since
$A_{\mathrm{ang}}(\pm\pi/2)=w<u_*w$, hence such a point is interior and
$\eta'(\varphi)=0$.  Formula \eqref{eq:unified-base-3-29} at
$(u,v)=(u_*,w)$ gives the curvature $\sec_{\widetilde g_{u,v}}(E,V)=(u_*w)^{-2}$, which should greater
than one by Lemma~\ref{lem:angular-sphere-geometry}.  Hence
\begin{equation}
u_*w<1.                              \label{eq:unified-base-3-31}
\end{equation}
The inequality $u_*>\rho/w$ gives $\alpha_0>1$, while
\eqref{eq:unified-base-3-10} and \eqref{eq:unified-base-3-31} give
\[
 \log u_*<-\log w<m(\log\rho-\log w).
\]
Thus
\begin{equation}
1<\alpha_0<m,                        \label{eq:unified-base-3-33}
\end{equation}
and $v_0(1)=\rho$, $v_0(u_*)=w$.  We claim that every sectional
curvature of $\widetilde g_{u,v_0(u)}$ is greater than one.

First observe that the numerator
\[
 1+(u-1)\bigl(\eta-\eta'\tan\varphi\bigr)
\]
in \eqref{eq:unified-base-3-29} is positive at $u=1$ and at $u=u_*$:
at the first endpoint it equals one, and at the second its positivity
follows from $\sec_{g_{\mathrm{ang},m}}>1$.  Since it is affine in
$u$, it is positive throughout $[1,u_*]$.  When $m=3$, the numerator
$f_u^2-\sin^2\varphi$ in \eqref{eq:unified-base-3-30} is also
positive.  If $\eta<0$, it decreases with $u$ and is positive at
$u=u_*$; if $\eta\ge0$, then $f_u\ge1$.

Differentiating \eqref{eq:unified-base-3-29} and \eqref{eq:unified-base-3-30} with respect to $u$ gives
\begin{align}
 \frac{\dd}{\dd u}\log\sec_{\widetilde g_{u,v_0(u)}}(E,V)
 &=\frac1{\alpha_0u}\left(
  2+\frac{\alpha_0u(\eta-\eta'\tan\varphi)}
           {1+(u-1)(\eta-\eta'\tan\varphi)}
   -\frac{3\alpha_0u\eta}{f_u}\right),                           \label{eq:unified-base-3-34}\\
 \frac{\dd}{\dd u}\log\sec_{\widetilde g_{u,v_0(u)}}(V,W)
 &=\frac2{\alpha_0u}\left(
  1+\frac{\alpha_0u\eta\sin^2\varphi}
          {f_u(f_u^2-\sin^2\varphi)}\right).                     \label{eq:unified-base-3-35}
\end{align}
For \eqref{eq:unified-base-3-35}, if $\eta\ge0$ the derivative is
nonnegative.  If $\eta<0$, its bracket is decreasing because
\[
 \frac{\dd}{\dd u}\frac{u}{f_u(f_u^2-\sin^2\varphi)}>0.
\]
Thus its sign can change only from positive to negative, so the
curvature $\sec_{\widetilde g_{u,v_0(u)}}(V,W)$ has no interior minimum.

For \eqref{eq:unified-base-3-34}, if
$\eta'\tan\varphi<0$, then $f_u\le u$ and
\begin{equation}
\sec_{\widetilde g_{u,v_0(u)}}(E,V)
 >\frac1{v_0(u)^2f_u^2}
 \ge\frac1{u^2v_0(u)^2}>1.                                      \label{eq:unified-base-3-36}
\end{equation}
Indeed, $uv_0(u)$ is increasing because $\alpha_0>1$, and its terminal
value is $u_*w<1$.  It remains to consider the case when
$\eta'\tan\varphi\ge0$.  If $\eta\ge0$, the right side of
\eqref{eq:unified-base-3-34} equals
\begin{equation}
\frac1{\alpha_0uf_u}\left(
 2(1-\eta)+2\eta u(1-\alpha_0)
  -\frac{\alpha_0u\eta'\tan\varphi}
         {1+(u-1)(\eta-\eta'\tan\varphi)}\right).                \label{eq:unified-base-3-37}
\end{equation}
The bracket is decreasing because
\[
 \frac{\dd}{\dd u}\left(
   \frac u{1+(u-1)(\eta-\eta'\tan\varphi)}\right)
 =\frac{1-\eta+\eta'\tan\varphi}
        {[1+(u-1)(\eta-\eta'\tan\varphi)]^2}\ge0.
\]
If $\eta<0$, \eqref{eq:unified-base-3-34} becomes
\begin{equation}
\frac1{\alpha_0u}\left(
 2-\frac{\alpha_0u}{f_u}\left(
  2\eta+\frac{\eta'\tan\varphi}
    {1+(u-1)(\eta-\eta'\tan\varphi)}\right)\right).             \label{eq:unified-base-3-38}
\end{equation}
At every zero of its bracket the last parenthesis is positive.  Also,
\[
 \left(\frac u{f_u}\right)'=\frac{1-\eta}{f_u^2}>0
\]
and
\[
 \frac{\dd}{\dd u}\left(
   \frac{\eta'\tan\varphi}
        {1+(u-1)(\eta-\eta'\tan\varphi)}\right)
 =\frac{\eta'\tan\varphi\,(\eta'\tan\varphi-\eta)}
        {[1+(u-1)(\eta-\eta'\tan\varphi)]^2}\ge0.
\]
Hence this bracket can cross zero only from positive to negative, so
the curvature $\sec_{\widetilde g_{u,v_0(u)}}(E,V)$ also has no interior minimum.  At  endpoints, both curvatures are greater than one: at $u=1$ the metric is the round
sphere of radius $\rho<1$, and at $u=u_*$ it is
$g_{\mathrm{ang},m}$.  This proves the claim.
\end{proof}


We next define the neck metric on $I\times \mathbb S^m$ and calculate its curvature.
\begin{lemma}\label{lem:variable-angular-ricci}
Let \(u_*\), \(\eta\), \(f_u=1+(u-1)\eta\), and
\(\widetilde g_{u,v}\) be the one introduced in
Lemma~\ref{lem:angular-slice-interpolation}.  Let
\(I\subset(0,\infty)\) be an interval and the functions
\[
 u:I\longrightarrow[1,u_*],\qquad v:I\longrightarrow(0,\infty)
\]
be smooth.  On \(I\times\mathbb S^m\), define the neck metric
\begin{equation}
\widehat g_{\mathrm{neck},m}=\dd t^2+t^2v(t)^2\left(
 f_{u(t)}(\varphi)^2\dd\varphi^2
 +\cos^2\varphi\,g_{\mathbb S^k}\right).                         \label{eq:unified-base-3-39}
\end{equation}
Then the following identities for sectional curvature and Ricci curvature
\eqref{eq:unified-base-3-40}--\eqref{eq:unified-base-3-47} hold.
All off-diagonal Ricci components vanish except possibly the
\((\partial_t,E)\)-component in
\eqref{eq:unified-base-3-45}.
\end{lemma}

\begin{proof}
Let $E=(tvf_u)^{-1}\partial_\varphi$, and let $V,W$ be orthonormal
vectors tangent to the $\mathbb S^k$ factor.  Here and throughout this
calculation, we write $u=u(t)$, $v=v(t)$, and
$f_u=f_{u(t)}(\varphi)$ for brevity.
The principal curvatures of the $t$-slices with respect to
$\partial_t$ are
\begin{equation}
\II_t(E,E)=\frac1t+\frac{v'}v+\frac{u'\eta}{f_u},
 \qquad
 \II_t(V,V)=\frac1t+\frac{v'}v.                                 \label{eq:unified-base-3-40}
\end{equation}
Since the metric induced on the $t$-slice is
$t^2\widetilde g_{u(t),v(t)}$, the Gauss equation then gives
\begin{align}
 \sec_{\widehat g_{\mathrm{neck},m}}(E,V)
   &=\frac1{t^2}\sec_{\widetilde g_{u,v}}(E,V)
           -\II_t(E,E)\II_t(V,V),                                      \label{eq:unified-base-3-41}\\
 \sec_{\widehat g_{\mathrm{neck},m}}(V,W)
   &=\frac1{t^2}\sec_{\widetilde g_{u,v}}(V,W)
           -\II_t(V,V)^2.                                               \label{eq:unified-base-3-42}
\end{align}
The second formula \eqref{eq:unified-base-3-42} occurs only for $m=3$.  The normal variation of the
slice metric gives
\begin{align}
 \sec_{\widehat g_{\mathrm{neck},m}}(\partial_t,V)
   &=-\frac{(tv)''}{tv},                                                 \label{eq:unified-base-3-43}\\
 \sec_{\widehat g_{\mathrm{neck},m}}(\partial_t,E)
   &=-\frac{\partial_t^2(tv f_u)}{tv f_u}\notag\\
   &=-\frac{(tv)''}{tv}
      -2\left(\frac1t+\frac{v'}v\right)\frac{u'\eta}{f_u}
      -\frac{u''\eta}{f_u}.                                           \label{eq:unified-base-3-44}
\end{align}
The only possibly nonzero mixed component of the Ricci tensor is
\begin{equation}
\Ric_{\widehat g_{\mathrm{neck},m}}(\partial_t,E)
 =-\frac{k\eta(\varphi)u'(t)\tan\varphi}
         {tv(t)f_{u(t)}(\varphi)^2}.                              \label{eq:unified-base-3-45}
\end{equation}
Smoothness at the poles gives
$\eta(\varphi)\tan\varphi=O(\cos\varphi)$, so this expression
extends smoothly to $\varphi=\pm\pi/2$.  Tracing the sectional
curvatures yields
\begin{equation}
\begin{aligned}
 \Ric_{\widehat g_{\mathrm{neck},m}}(\partial_t,\partial_t)
   &=\sec_{\widehat g_{\mathrm{neck},m}}(\partial_t,E)
     +k\sec_{\widehat g_{\mathrm{neck},m}}(\partial_t,V),\\
 \Ric_{\widehat g_{\mathrm{neck},m}}(E,E)
   &=\sec_{\widehat g_{\mathrm{neck},m}}(\partial_t,E)
     +k\sec_{\widehat g_{\mathrm{neck},m}}(E,V),\\
 \Ric_{\widehat g_{\mathrm{neck},m}}(V,V)
   &=\sec_{\widehat g_{\mathrm{neck},m}}(\partial_t,V)
     +\sec_{\widehat g_{\mathrm{neck},m}}(E,V)
     +(k-1)\sec_{\widehat g_{\mathrm{neck},m}}(V,W).
 \end{aligned}                                                     \label{eq:unified-base-3-46}
\end{equation}
The last identity is absent when $k=1$, and all off-diagonal components
other than \eqref{eq:unified-base-3-45} vanish.  In particular,
\eqref{eq:unified-base-3-43}--\eqref{eq:unified-base-3-44} imply
\begin{equation}
\Ric_{\widehat g_{\mathrm{neck},m}}(\partial_t,\partial_t)
 =-m\frac{(tv)''}{tv}
   -\frac{\eta}{f_u}\left(
       u''+2\left(\frac1t+\frac{v'}v\right)u'\right).           \label{eq:unified-base-3-47}
\end{equation}
\end{proof}
We then construct the family of neck metric $\widehat g_{\mathrm{neck},m}$ with positive Ricci curvature.
\begin{lemma}
\label{lem:slow-logarithmic-neck}
Under the hypotheses of
Lemma~\ref{lem:angular-slice-interpolation}, there is a family of metric,
parametrized by all sufficiently small \(\sigma>0\), of numbers
\(t_0<t_1(\sigma)\) and smooth functions
\[
 u_\sigma:[t_0,t_1(\sigma)]\longrightarrow[1,u_*],\qquad
 v_\sigma:[t_0,t_1(\sigma)]\longrightarrow(0,\infty)
\]
such that
\[
 u_\sigma(t_0)=1,\quad v_\sigma(t_0)=\rho,\qquad
 u_\sigma(t_1)=u_*,\quad v_\sigma(t_1)>w,
\]
\[
 t_1(\sigma)\longrightarrow\infty\quad(\sigma\downarrow0),
\qquad
\Ric_{\widehat g_{\mathrm{neck},m}}>0 ,
\]
where \(\widehat g_{\mathrm{neck},m}\) is defined by
\eqref{eq:unified-base-3-39} with
\(u=u_\sigma\) and \(v=v_\sigma\).
\end{lemma}

\begin{proof}
We write \(u=u_\sigma\) and \(v=v_\sigma\).  We now choose the
parameters and these two functions $u$, $v$, so that: the corresponding path
of slice metrics must remain positively curved, and the coefficient of
the leading term in the radial Ricci curvature must remain positive.
By compactness, there exists $\mu>0$, such that all sectional curvatures of
$\widetilde g_{u,v_0(u)}$, $1\le u\le u_*$, are at least
$1+2\mu$.  By Smooth dependence on the parameters and compactness, there exists
\(\delta_\alpha>0\), such that
\[
 \sec_{\widetilde g_{u,\rho u^{-1/a}}}\geq1+\mu
 \quad\text{for}\quad
 1\leq u\leq u_*,\quad
 \alpha_0\leq a\leq\alpha_0+\delta_\alpha .
\]
Choose $0<\varepsilon_0<1$ so that
\[
 \frac{\alpha_0}{1-\varepsilon_0}
 <\min\{m,\alpha_0+\delta_\alpha\}.
\]
Then choose $\sigma_0>0$ so that, for every
$0<\sigma\le\sigma_0$,
\begin{equation}
\alpha=\frac{1+\sigma}{1-\varepsilon_0}\alpha_0
 <\min\{m,\alpha_0+\delta_\alpha\},
 \qquad
 v_\alpha(u)=\rho u^{-1/\alpha},\qquad 1\le u\le u_*,           \label{eq:unified-base-3-48}
\end{equation}
and the sectional curvatures of
$\widetilde g_{u,v_\alpha(u)}$ are at least $1+\mu$.  The parameters
we chosen so far satisfy the following:
\begin{equation}
1<u_*<\frac1w,\qquad
 1<\alpha_0<\alpha<m,\qquad
 0<\varepsilon_0<1,\qquad 0<\sigma\le\sigma_0,\qquad \mu>0.      \label{eq:unified-base-3-49}
\end{equation}

Fix a cut-off function
\[
 \chi\in C^\infty([0,\infty),[0,1]),\qquad
 \chi(0)=0,\quad \chi'(0)>0,\quad \chi'\geq0,
\]
\[
 \chi(x)>0\ \text{ for } x>0,\qquad \chi(x)=1\ \text{ for } x\geq1.
\]
Choose \(t_0\) large, denote \(L=\log(2t_0)\), and define
\begin{equation}
\Gamma(t)=\frac{L}{t(\log t)^2}
 \chi\!\left(\frac{t-t_0}{t_0}\right),
 \qquad t\geq t_0.                                                \label{eq:unified-base-3-50}
\end{equation}
Set \(q(t)=L/[t(\log t)^2]\), then
\[
 \Gamma'=q'\chi+\frac q{t_0}\chi',\qquad
 \Gamma'+\frac{2\Gamma}{t}
 =q\left[
 \frac{\chi'}{t_0}
 +\chi\left(\frac1t-\frac2{t\log t}\right)\right].
\]
The continuous function
\[
 x\longmapsto (1+x)\chi'(x)+\frac12\chi(x)
\]
has a positive minimum on \([0,1]\).  Since \(\chi=1\) on
\([1,\infty)\), there are constants \(c,C>0\), independent of the constant
\(t_0\), such that
\begin{align}
 0\le\Gamma&\le\frac{CL}{t(\log t)^2}, \qquad
 |\Gamma'|\le\frac{CL}{t^2(\log t)^2},                           \label{eq:unified-base-3-51}\\
\Gamma'+\frac{2\Gamma}{t}
 &\ge\frac{cL}{t^2(\log t)^2}.                                   \label{eq:unified-base-3-52}
\end{align}
Moreover, since \(\chi=1\) for \(t\geq2t_0\),
\begin{equation}\label{eq:I t0 bound}
 1=\int_{2t_0}^{\infty}\frac{L}{t(\log t)^2}\,\dd t
 \leq I_{t_0}:=\int_{t_0}^{\infty}\Gamma(t)\,\dd t
 \leq\int_{t_0}^{\infty}\frac{L}{t(\log t)^2}\,\dd t
 =\frac{L}{\log t_0}\leq2
\end{equation}
after increasing \(t_0\).
Define
\begin{equation}
\beta=\frac{(1-\varepsilon_0)(\log\rho-\log w)}
              {I_{t_0}},                                         \label{eq:unified-base-3-53}
\end{equation}
and solve
\begin{equation}
\frac{u'}u=\alpha\beta\Gamma,\qquad
 \frac{v'}v=-\beta\Gamma,\qquad u(t_0)=1,\quad v(t_0)=\rho.    \label{eq:unified-base-3-54}
\end{equation}
The two differential equations force $v=\rho u^{-1/\alpha}$, so the
slice metrics follow the path fixed in \eqref{eq:unified-base-3-48}.
The definitions of \(\alpha\) and \(\beta\) give
\[
 \alpha\beta I_{t_0}=(1+\sigma)\log u_*.
\]
Since \(\Gamma>0\) on \((t_0,\infty)\), there is a unique \(t_1\)
such that \(u(t_1)=u_*\), and
\begin{equation}\label{eq:I t0 identity}
 \int_{t_1}^{\infty}\Gamma(t)\,\dd t
 =\frac{\sigma}{1+\sigma}I_{t_0}.
\end{equation}
For all sufficiently small \(\sigma\), the right-hand side of \eqref{eq:I t0 identity} is less than \(1\), whereas
\(\int_{2t_0}^{\infty}\Gamma\,\dd t=1\), hence \(t_1>2t_0\).
Equation \eqref{eq:unified-base-3-50} then yields
\[
 \frac{L}{\log t_1}
 =\frac{\sigma}{1+\sigma}I_{t_0},
\]
and therefore
\begin{equation}
v(t_1)=\rho u_*^{-1/\alpha}>w,\qquad
                         t_1\longrightarrow\infty
                         \quad(\sigma\downarrow0).                \label{eq:unified-base-3-55}
\end{equation}
We now restrict the metric \eqref{eq:unified-base-3-39} to
$[t_0,t_1]\times\mathbb S^m$.

The bounds \(1\leq I_{t_0}\leq2\) from \eqref{eq:I t0 bound} show that
\[
 0<c_\beta\leq\beta\leq C_\beta<\infty .
\]
Here \(c,C,c_i,C_i,c_\beta,C_\beta\) may depend on the fixed parameters
\[
 (m,r,\rho,w,\varepsilon_0,\chi),
\]
but are independent of sufficiently large \(t_0\), of
\(0<\sigma\leq\sigma_0\), and of \(t\in[t_0,t_1]\).
Equations \eqref{eq:unified-base-3-51}--\eqref{eq:unified-base-3-53} give
\begin{equation}
\beta\Gamma\le\frac{CL}{t(\log t)^2},\qquad
 \beta\left|\Gamma'+\frac{2\Gamma}{t}\right|
 \le\frac{CL}{t^2(\log t)^2},\qquad
 \beta^2\Gamma^2\le\frac{CL^2}{t^2(\log t)^4},                \label{eq:unified-base-3-56}
\end{equation}
and
\begin{equation}
\beta\left(\Gamma'+\frac{2\Gamma}{t}\right)
 \ge\frac{cL}{t^2(\log t)^2}.                                   \label{eq:unified-base-3-57}
\end{equation}
Substituting \eqref{eq:unified-base-3-54} into
\eqref{eq:unified-base-3-40} gives
\begin{equation}
\II_t(E,E)=\frac1t-\beta\Gamma
       +\frac{\alpha\beta\Gamma u\eta}{f_u},
 \qquad
 \II_t(V,V)=\frac1t-\beta\Gamma.                              \label{eq:unified-base-3-58}
\end{equation}
It follows from \eqref{eq:unified-base-3-56} and the boundedness of
$u\eta/f_u$ that every product of two slice principal curvatures is
$t^{-2}+O(L/[t^2(\log t)^2])$.  The sectional curvatures of
$\widetilde g_{u(t),v(t)}$ are at least $1+\mu$; hence after increasing $t_0$, the Gauss
equations \eqref{eq:unified-base-3-41}--\eqref{eq:unified-base-3-42}
give, 
\begin{equation}
\sec_{\widehat g_{\mathrm{neck},m}}(E,V)\ge\frac{c_0}{t^2},\qquad
 \sec_{\widehat g_{\mathrm{neck},m}}(V,W)\ge\frac{c_0}{t^2}\quad(m=3).      \label{eq:unified-base-3-59}
\end{equation}
whereas
\begin{equation}
|\sec_{\widehat g_{\mathrm{neck},m}}(\partial_t,E)|
 +|\sec_{\widehat g_{\mathrm{neck},m}}(\partial_t,V)|
 +|\Ric_{\widehat g_{\mathrm{neck},m}}(\partial_t,E)|
 \le\frac{C_0L}{t^2(\log t)^2}.                                 \label{eq:unified-base-3-60}
\end{equation}
Here the last estimate uses the boundedness of
$\eta\tan\varphi/f_u^2$.  Then the identity 
\eqref{eq:unified-base-3-46}, together with
\eqref{eq:unified-base-3-59}--\eqref{eq:unified-base-3-60}, imply that
$\Ric_{\widehat g_{\mathrm{neck},m}}(E,E)$ and
$\Ric_{\widehat g_{\mathrm{neck},m}}(V,V)$ are at least
$c_1/t^2$.

Finally, substituting \eqref{eq:unified-base-3-54} into \eqref{eq:unified-base-3-47} gives
\begin{equation}
\begin{aligned}
 \Ric_{\widehat g_{\mathrm{neck},m}}(\partial_t,\partial_t)
 &=\beta\left(m-\alpha\frac{u\eta}{f_u}\right)
          \left(\Gamma'+\frac{2\Gamma}{t}\right)\\
 &\quad-\beta^2\Gamma^2
        \left(m+\alpha(\alpha-2)\frac{u\eta}{f_u}\right).
\end{aligned}                                                     \label{eq:unified-base-3-61}
\end{equation}
If $\eta<0$, then $u\eta/f_u<0$; if $\eta\ge0$, then
$0\le u\eta/f_u\le1$.  Since $\alpha<m$, the first coefficient in
\eqref{eq:unified-base-3-61} has a positive uniform lower bound, while the second
is uniformly bounded.  Consequently,
\begin{equation}
\Ric_{\widehat g_{\mathrm{neck},m}}(\partial_t,\partial_t)
 \ge\frac{c_3L}{t^2(\log t)^2}
     -\frac{C_3L^2}{t^2(\log t)^4}>0.                             \label{eq:unified-base-3-62}
\end{equation}
Moreover,
\begin{equation}
\Ric_{\widehat g_{\mathrm{neck},m}}(\partial_t,\partial_t)
   \Ric_{\widehat g_{\mathrm{neck},m}}(E,E)
 -\Ric_{\widehat g_{\mathrm{neck},m}}(\partial_t,E)^2
 \ge\frac{c_4L}{t^4(\log t)^2}
     -\frac{C_4L^2}{t^4(\log t)^4}>0.                             \label{eq:unified-base-3-63}
\end{equation}
Thus the $2\times2$ Ricci block on
$\vspan\{\partial_t,E\}$ has positive first diagonal
entry and positive determinant, hence is positive definite.  The
vertical diagonal entry is positive and has no mixed components with
this block.  Therefore
$\Ric_{\widehat g_{\mathrm{neck},m}}>0$ for $m=2,3$.
\end{proof}


We now normalize the neck metric in the form \eqref{eq:unified-base-3-13} on
\([0,\ell]\times\mathbb S^m\). 
\begin{lemma}\label{lem:neck-endpoint-normalization}
Let \(\widehat g_{\mathrm{neck},m}\) be the family of neck metric 
constructed in Lemma~\ref{lem:slow-logarithmic-neck}.  For all
sufficiently small \(\sigma>0\), there exists
\(\kappa,\ell,\lambda>0\), such that, under the change of variables
\[
 s=\kappa(t-t_0),\qquad t=t_0+\frac{s}{\kappa},
\]
the metric \(\kappa^2\widehat g_{\mathrm{neck},m}\) has the form
\eqref{eq:unified-base-3-13} on
\([0,\ell]\times\mathbb S^m\).  The boundary metric on $\{0\}\times \mathbb S^m$ is
\((\rho/\lambda)^2g_{\mathbb S^m}\), with every outward
principal curvature equal to \(-\lambda\).  The boundary metric on $\{\ell\}\times \mathbb S^m$
is \(g_{\mathrm{ang},m}\), with every outward principal
curvature greater than \(1\), with
\[
                         B'(s)>0,\qquad B''(s)<0 .
\]
Such normalization is natural under simultaneous translations
\(z\mapsto z+c\), \(z_0\mapsto z_0+c\), and under orthogonal
transformations of \(\mathbb S^k\).
\end{lemma}

\begin{proof}

Set
\begin{equation}
\kappa=\frac{w}{t_1v(t_1)}.                                     \label{eq:unified-base-3-64}
\end{equation}
Multiply $\widehat g_{\mathrm{neck},m}$ by $\kappa^2$.  Since
$f_{u_*}(\varphi)=A_{\mathrm{ang}}(\varphi)/w$ and
$\kappa t_1v(t_1)=w$, the metric induced
at $t=t_1$ is exactly $g_{\mathrm{ang},m}$.  At $t=t_0$, the identity
$f_1(\varphi)=1$ shows that the induced metric is round of radius
\begin{equation}
\kappa t_0\rho=\frac\rho\lambda,
 \qquad \lambda=\frac1{\kappa t_0}.                             \label{eq:unified-base-3-65}
\end{equation}
Since $\Gamma(t_0)=0$, then by \eqref{eq:unified-base-3-58}, all outward principal curvatures on $\{0\}\times \mathbb S^m$
are $-\lambda$.  Indeed, principal curvatures are divided by
\(\kappa\) under the constant rescaling
\(\widehat g\mapsto\kappa^2\widehat g\).  The  vertical
principal curvatures  on $\{\ell\}\times \mathbb S^m$ are
\begin{equation}
\frac{v(t_1)}{w}\bigl(1-\beta t_1\Gamma(t_1)\bigr),             \label{eq:unified-base-3-66}
\end{equation}
and the meridional one is 
\begin{equation}
\frac{v(t_1)}{w}\bigl(1-\beta t_1\Gamma(t_1)\bigr)+\frac{t_1v(t_1)}{w}
   \frac{\eta(\varphi)u'(t_1)}{f_{u_*}(\varphi)}.                \label{eq:unified-base-3-67}
\end{equation}
Because $u\eta/f_u$ is bounded on the compact parameter set, all 
principal curvatures on $\{\ell\}\times \mathbb S^m$ are bounded below by
\begin{equation}
\frac{v(t_1)}{w}\left[1-C\beta t_1\Gamma(t_1)\right].           \label{eq:unified-base-3-68}
\end{equation}
As $\sigma\downarrow0$,
\begin{equation}
\frac{v(t_1)}{w}\longrightarrow
       u_*^{\varepsilon_0/\alpha_0}>1,\qquad
 t_1\Gamma(t_1)=\frac{L}{(\log t_1)^2}\longrightarrow0.         \label{eq:unified-base-3-69}
\end{equation}
Thus \eqref{eq:unified-base-3-68} is greater than one after reducing $\sigma$.
Finally,
\begin{equation}
(tv)'=v(1-\beta t\Gamma)>0,\qquad
 \frac{(tv)''}{tv}
 =-\beta\left(\Gamma'+\frac{2\Gamma}{t}\right)
       +\beta^2\Gamma^2<0                                        \label{eq:unified-base-3-70}
\end{equation}
for large $t_0$.  Set
\[
 s=\kappa(t-t_0),\qquad \ell=\kappa(t_1-t_0),\qquad
 t=t_0+\frac{s}{\kappa},
\]
and choose
\[
 A(s,\varphi)=\kappa tv(t)f_{u(t)}(\varphi),
 \qquad B(s)=\kappa tv(t).
\]
In the $s$-coordinate, $\kappa^2\widehat g_{\mathrm{neck},m}$ is precisely
\eqref{eq:unified-base-3-13}, and \eqref{eq:unified-base-3-70} gives $B'>0$ and $B''<0$.
Every function used in the construction depends on the ambient
\(z\)-coordinate only through \(z-z_0\), and depends on the
\(\mathbb S^k\)-variable only through \(g_{\mathbb S^k}\).  This proves
the asserted naturality.
\end{proof}

\begin{proof}[Proof of Proposition~\ref{prop:unified-neck}]
Choose \(r_*=r_{\mathrm{amb}}\), where
\(r_{\mathrm{amb}}\) is introduced in
Lemma~\ref{lem:neck-ambient-curvature}.  For
\(0<r<r_*\), choose
\[
 0<w<w_0(r)
\]
sufficiently small that \(w^{m-1}<\rho^m\), which is
\eqref{eq:unified-base-3-10}.  Since
Lemma~\ref{lem:concave-warping} gives the desired auxiliary function
\(\mathcal R\), Lemma~\ref{lem:neck-ambient-curvature} then proves
\textup{(i)}.

By \eqref{eq:unified-base-3-5}, \(\mathcal R(r)=w\tan r\).
  Hence
Lemmas~\ref{lem:angular-sphere-geometry}
and~\ref{lem:angular-latitude-form} prove \textup{(ii)}.
Since \(r<\pi/10\) and \(\pi\rho<r\), we have
\[
                 0<\rho<\frac1{10}<\cos r .
\]
Together with \eqref{eq:unified-base-3-10}, these parameters satisfy
hypotheses of Lemma~\ref{lem:angular-slice-interpolation}.
Lemmas~\ref{lem:variable-angular-ricci}
and~\ref{lem:slow-logarithmic-neck} then produce a family of
Ricci-positive metrics $\widehat g_{\mathrm{neck},m}$.
Lemma~\ref{lem:neck-endpoint-normalization} proves all the
claims in \textup{(iii)}.

A simultaneous translation \(z\mapsto z+c\) and
\(z_0\mapsto z_0+c\) preserves the difference \(z-z_0\) and therefore
carries the construction centered at \(o_{z_0}\) isometrically to the
construction centered at \(o_{z_0+c}\).  Moreover, every metric and
auxiliary function depends on the \(\mathbb S^k\)-variable only through
\(g_{\mathbb S^k}\).  Hence the entire construction is natural under
translations in the \(z\)-direction and orthogonal transformations of
\(\mathbb S^k\).
\end{proof}





\bibliographystyle{amsplain}
\bibliography{references}

\end{document}
